\documentclass[11pt]{article}
\usepackage[a4paper,margin=2.2cm]{geometry}
\usepackage{amsmath,amssymb}
\usepackage{booktabs,longtable,array,tabularx}
\usepackage[T1]{fontenc}
\usepackage{lmodern}
\usepackage{microtype}
\usepackage{graphicx}
\usepackage[strings]{underscore}
\usepackage{float}
\usepackage{xcolor}
\usepackage[hidelinks]{hyperref}
\sloppy
\setlength{\parskip}{0.45em}
\setlength{\parindent}{0pt}
\setcounter{tocdepth}{1}
\renewcommand{\arraystretch}{1.08}
\newcolumntype{L}{>{\raggedright\arraybackslash}X}
\newcommand{\bros}{\beta_r}
\newcommand{\gap}{\Delta E_h(0)^{*}}

\usepackage{xr}
\externaldocument[M-]{methods}
\newcommand{\compname}{Methods}
\title{Figure generation}
\author{}
\date{}
\begin{document}
\maketitle
\tableofcontents


\section*{What this document is}
\addcontentsline{toc}{section}{What this document is}

This document describes how every display item of the manuscript was generated:
the population it is drawn from, the selection applied to it, the solver settings,
the quantities plotted and the numbers quoted in its caption. It is the companion
to \texttt{methods.pdf}, which describes the model itself, its numerical treatment
and the classification of parameter combinations, and which this document refers
to throughout. Section, table and equation numbers given as ``of the Methods''
are numbers in that document.

Figures are named by the index they carry in the submitted dossier: \textbf{Fig.~1}
to \textbf{Fig.~5} in the main text, \textbf{Extended Data Fig.~1} to
\textbf{Extended Data Fig.~7} in the Extended Data, and
\textbf{Supplementary Fig.~1} to \textbf{Supplementary Fig.~3} in the
Supplementary Information. Where one analysis produces panels that now sit in two
documents, the section is titled with both.

\textbf{Fig.~1} is the model diagram. It is drawn by hand rather than computed and
so has no entry here; the model it depicts is Section~\ref{M-sec:model} of the Methods.

The conventions used throughout -- what \emph{Tregs OFF} and \emph{Tregs ON} mean,
what an initial condition and a final state are, the units, and the names of the
classes and branches of parameter combinations -- are those set out under
``Conventions'' at the start of the Methods. Every number below was recomputed from
the stored data by a script; the script or expression is named in a comment in the
LaTeX source (comments beginning \texttt{\% from:}).

\section{Overview of the figures}

Table~\ref{tab:overview} lists each figure, the population it is drawn from, its
size, and the section that describes it. Table~\ref{tab:fig-script-data}
(Section~\ref{sec:fig-script-data}) maps each figure to its script and data files.

\begin{table}[H]
\centering
\caption{Overview: figures, populations, sizes and the sections that describe them.}
\label{tab:overview}
\footnotesize
\begin{tabularx}{\linewidth}{@{}l L L L l@{}}
\toprule
Figure & Content & Population & $N$ & Section\\\midrule
Fig.~2 & efferocytosis knockout & ungated draw of (kinetic vector, $\bros$, $r_P$) units, 8 arms, 21 initial conditions each & 10{,}500 units; C: 2{,}537 sterile (unit, $K_g$) pairs from 629 units; D, E: 3{,}294 pathogenic pairs from 836 units; F, G: 1{,}228 units (836 production, 392 additional), 180{,}515 (unit, $e$, $K_g$) pairs & \ref{sec:fig1}\\
Figs.~3A, 3B; Extended Data Figs.~2A, 2B & Tregs shift both attractors and widen the range of recoverable injury & A, B: the two example sterile help combinations (one bistable, one ON\_mono) of the gated sterile branch; C, D: the two example pathogenic help combinations (one bistable, one ON\_mono) of the gated pathogenic help branch & 4 combinations, 10 initial conditions each & \ref{sec:fig3}\\
Fig.~3C; Extended Data Fig.~2C & load balance, bistable class & gated bistable Treg-help combinations, rescue gap above a minimum & 442 sterile, 180 pathogenic help (20 drawn per branch) & \ref{sec:fig4}\\
Fig.~4 & damage-matched comparison & every gated pathogenic combination in which the pathogen persists without Tregs, no condition on the Treg effect; 3 initial conditions & 75{,}000 drawn, 52{,}010 run, 8{,}789 pass the filter (every panel); B, C: 1{,}000; D: one unit & \ref{sec:fig6}\\
Fig.~5; Extended Data Fig.~6 & Treg stability and specificity & stratified draw from the load-balance eligibility (Fig.~3C, Extended Data Fig.~2C), Treg-help class, rescue probability falling in both properties & 977 drawn, 966 retained, 717 crossing, 440 used & \ref{sec:fig7}\\
Extended Data Fig.~1 & two forms of rescue, sterile and pathogenic help: attractor shifts and rescue gaps & the Tregs-help class of the gated sterile and gated pathogenic help branches & sterile help 43{,}375; A, C: 33{,}487; B: 9{,}884; pathogenic help 17{,}495; D, F: 13{,}725; E: 1{,}270 & \ref{sec:figS01}\\
Extended Data Fig.~3 & load balance, ON\_mono class & gated ON\_mono Treg-help combinations & 453 sterile, 500 pathogenic help (20 drawn per branch) & \ref{sec:figS02}\\
Extended Data Fig.~4 & what a stronger ROS response buys and costs without Tregs, over the sweep's $(\beta_r, r_P)$ grid & gated pathogenic combinations whose Tregs-OFF system is bistable with an increasing threshold at all nine $\beta_r$ levels & 12{,}211 parameter sets $\times$ 9 levels, 2 stable steady states each & \ref{sec:figS03}\\
Extended Data Fig.~5 & three classifications of the damage-matched pairs & the 8{,}789 pairs of Fig.~4 & 8{,}789 (15 eliminating, 8{,}774 persistent) & \ref{sec:figS04}\\
Extended Data Fig.~7 & where the rescue probability is gained, and its spread & the combinations of Fig.~5, both branches & 440 of 717 crossing & \ref{sec:figS05}\\
Supplementary Fig.~1 & basin boundary at $t=0$ & gated sterile and pathogenic combinations where Tregs help & A: one combination; B: 16{,}990 of the 89{,}328 sterile combinations measured, help class, 2{,}763--2{,}802 rows per coordinate; C: 17{,}495 pathogenic help, 2{,}367--2{,}714 rows & \ref{sec:fig5}\\
Supplementary Fig.~2 & two forms of rescue, pathogenic harm: phase portraits, attractor shifts and rescue gap & the two example combinations (A, B) and the gated pathogenic harm branch (C--E) & 2 combinations, 10 initial conditions each; 5{,}246; C, E: 1{,}742; D: 1{,}180 & \ref{sec:figS06}\\
Supplementary Fig.~3 & Treg treatments where Tregs harm & gated harm combinations after three further steps & A: 250 help + 250 harm; B: 189; C: 1{,}745; D--G: one combination & \ref{sec:figS07}\\
\bottomrule
\end{tabularx}
\end{table}

%% ===========================================================================
\section{Generation of Fig.~2: efferocytosis knockout}
\label{sec:fig1}

%% ===========================================================================
%% fig01.tex -- Figure 2: the efferocytosis knockout
%% Code:  pipeline/06_effero_knockout.R  (data)   figure_scripts/Fig_2.R  (figure)
%% Every number below was recomputed on 2026-10-01 from ./data/effero_knockout.rds
%% by figure_scripts/Fig_2.R and is stored in ./source_data/Fig_2_numbers.csv
%% (quantity,value), ./source_data/Fig_2_panelCDE.csv and
%% ./source_data/Fig_2_panelFG.csv. Comments "% from: <quantity>" name the row.
%% ===========================================================================

%% ---- from docs/methods_parts/fig01.tex, lines 12-424
The main parameter sweep switches only regulatory T cells on and off; efferocytic
clearance is present in every one of its integrations. Fig.~2 asks what the
system does when clearance is removed, and therefore rests on a separate
experiment in which an efferocytosis knockout arm is integrated next to the intact
model. The model equations are those of Section~\ref{M-sec:model} of the Methods; only the elements
specific to this experiment are defined here. Regulatory T cells are inactive in
every integration of this experiment ($\theta = 0$).

\subsection{The knockout}
\label{sec:fig1-ko}

In the model, efferocytic clearance enters through the load-dependent clearance rate
$g(L)$ of Eq.~\ref{M-eq:g} of the Methods ($a_0 = 0.0045$, $a_{\mathrm{floor}} = 0.0002$, $n_g = 4$;
Table~\ref{M-tab:core-const} of the Methods), and $g$ appears in
exactly three terms: the clearance of apoptotic cells, $g\,(M_0+M_1)\,E_a$; the
clearance of necrotic cells, $g\,(M_0+M_1)\,E_n$; and the efferocytosis-driven
reversion of activated macrophages, $\tau_{\mathrm{eff}}\, g\, E_a M_1$ (an M1-to-M0
flux). The knockout multiplies $g$ by a factor $\varepsilon_{\mathrm{KO}} \in \{0,1\}$
($\varepsilon_{\mathrm{KO}}=1$ intact, $0$ knockout; \texttt{eff\_ko} in
\texttt{lib/skeleton\_defs.R}) and changes nothing else. Under knockout
$g \equiv 0$, so all three terms vanish together; the floor $a_{\mathrm{floor}}$ is
removed as well, i.e.\ there is no residual clearance. Because $g$ is then identically
zero, the knockout arm does not depend on $K_g$.

The knockout is deliberately not implemented as $a_0 = a_{\mathrm{floor}} = 0$. The
normalizing constants of the apoptotic and necrotic pools are defined from the
intact clearance rate $a_0$ (Eq.~\ref{M-eq:KEaKEn} of the Methods),
and the danger-signal production term contains $f_{E_n} = E_n/(K_{E_n} + E_n)$.
Setting $a_0 = 0$ would send $K_{E_n} \to \infty$ and $f_{E_n} \to 0$, i.e.\ it
would also switch off danger-signal production from necrotic debris, which is a
different perturbation. In the knockout used here $K_{E_a}$ and $K_{E_n}$ keep
their intact values, and only removal by efferocytosis is abolished.

\subsection{Sampling design}
\label{sec:fig1-design}

The sampling unit is a triple (kinetic vector, $\beta_r$, $r_P$), where the
kinetic vector is one row of the five stored kinetic-parameter grids
(\texttt{data/tau\_grid\_v6\_1..5.rds}, $5\times 2{,}000 = 10{,}000$ vectors of the 22
kinetic parameters of Table~\ref{M-tab:core-kin} of the Methods), $\beta_r$ (\texttt{beta\_ros}, ROS
production per M1 macrophage) takes the nine values $0.0015, 0.0025, \dots, 0.0095$,
and $r_P$ (\texttt{r\_L\_P}, the pathogen growth rate, applied to both the
luminal and the lamina-propria pathogen) takes the seven
values $0, 0.08, 0.16, \dots, 0.48$. Units with $r_P = 0$ form the \emph{sterile}
branch; units with $r_P > 0$ form the \emph{pathogenic} branch. These are the kinetic grids and the $\bros$ and $r_P$ levels of the main sweep
(Section~\ref{M-sec:sampling} of the Methods); unlike the sweep, the draw includes grids 4 and 5. All
rates are per model time unit; the populations are in the units of
Section~\ref{M-sec:units} of the Methods (carrying capacities $K_E = K_M = K_T = 100$).

Units were drawn stratified by $r_P$ (\texttt{pipeline/06\_effero\_knockout.R},
\texttt{N\_PER = 1500}, \texttt{SEED = 1}): for each of the seven $r_P$ levels,
$1{,}500$ kinetic vectors were drawn without replacement from the $10{,}000$, and each
was paired with a $\beta_r$ level drawn uniformly with replacement from the nine.
Every $r_P$ stratum therefore carries the same weight. This gives
$10{,}500$ units ($1{,}500$ sterile, $9{,}000$ pathogenic) over $6{,}789$ distinct
kinetic vectors (a vector can recur across strata, never within one).
% from: units, units_sterile, units_pathogenic, kinetic_vectors_distinct
% (the draw was re-run with UNITS_ONLY=true and reproduces all 10,500 stored units exactly)
The sample was drawn from the raw grid, not from combinations already found to be
bistable, so that units whose low-burden state could be destroyed by the knockout
are not removed before the experiment.

Each unit is integrated under eight \emph{arms}: efferocytosis intact
($\varepsilon_{\mathrm{KO}} = 1$) at each of the seven half-saturation loads
$K_g \in \{15, 20, 25, 30, 35, 40, 45\}$, and the knockout
($\varepsilon_{\mathrm{KO}} = 0$). We call one unit at one $K_g$ with efferocytosis ON a
\emph{(unit, $K_g$) pair} (73,500 in total: 10,500 sterile, 63,000 pathogenic) and one unit
under the knockout a \emph{knockout system} (10,500). A pair is judged on its
efferocytosis-ON line only; the knockout line of the same unit is shown next to it but
plays no part in the selection.
% from: systems_sterile, systems_pathogenic, ko_lines_sterile, ko_lines_pathogenic

\paragraph{Initial conditions.} Each arm is integrated from the 21 initial conditions
of the injury line (Eq.~\ref{M-eq:icline} of the Methods), indexed by the initial epithelial health
$e \in \{0, 4.95, 9.90, \dots, 99\}$ (21 equally spaced values, against the 11 probes of
the sweep): $E_h(0) = e$, $E_n(0) = K_E - e$, $M_0(0) = K_M = 100$, $T_0(0) = K_T = 100$,
and all other states zero ($E_a = M_1 = R = T_r = D = S = C_{LP} = 0$). In the pathogenic
branch the pathogen is seeded as $P_L(0) = 500$, $P_{LP}(0) = 10$; in the sterile
branch $P_L(0) = P_{LP}(0) = 0$. One unit under one arm, over the 21 initial
conditions, is a \emph{line}; the experiment comprises
$10{,}500 \times 8 \times 21 = 1{,}764{,}000$ integrations.

\paragraph{Solver and horizon.} Each integration uses \texttt{deSolve::ode} with method
\texttt{lsodar}, absolute tolerance $10^{-14}$, relative
tolerance $10^{-8}$, maximum internal step $h_{\max} = 5$ and at most $10^5$ steps
(Table~\ref{M-tab:solver} of the Methods),
from $t = 0$ to $t = 12{,}000$, with the extinction rule of the main sweep: a root
function fires when $\max(P_L, P_{LP})$ falls to $10^{-11}$, and the event sets
$P_L = P_{LP} = 0$. The right-hand side evaluates on $\max(y, 0)$ componentwise. The
stored \emph{final state} of an integration is the state at $t = 12{,}000$; it is not a
time average, and no convergence test or fixed-point refinement is applied to it.
An integration that raised an error or stopped before $t = 12{,}000$ produced no
record. We re-integrated three units (504 integrations, including one of the two
failures below) with the pipeline script and reproduced every stored final state to
within $3\times10^{-12}$ in absolute value, and reproduced the failure.

\paragraph{Integration outcomes.} Of the $1{,}764{,}000$ integrations, 2 did not reach
$t = 12{,}000$ (pathogenic, intact $K_g = 20$ and $K_g = 25$) and 3 returned a
non-finite final state (pathogenic, intact $K_g = 35, 40, 45$). All five are dropped,
leaving $1{,}763{,}995$ final states; five pathogenic efferocytosis-ON lines therefore have 20
rather than 21 final states, and every knockout line and every sterile line is
complete.
% from: rows_expected 1764000, rows_read 1763998, rows_missing 2, rows_nonfinite 3,
% rows_used 1763995, systems_pathogenic_lines_lt21 5, ko_lines_*_lt21 0, systems_sterile_lines_lt21 0

\subsection{Low-burden basin in this experiment: a final-state criterion}
\label{sec:fig1-basin}

No fixed points are computed in this experiment. Final states are compared on the
total burden $b$ of Eq.~\ref{M-eq:burden} of the Methods, which is $K_E - E_h$ in the sterile branch.
A line \emph{has a low-burden basin} if the finite final burdens of its
(21, or 20) initial conditions span more than $\texttt{GAP\_CUT} = 20$ burden
units,
\begin{equation}
  \max_e b_e - \min_e b_e > 20 .
  \label{eq:fig01-gap}
\end{equation}
Its \emph{low-burden final state} is the final state of minimum total burden and its
\emph{high-burden final state} the final state of maximum total burden. The threshold of
20 is the same number as the gate used for the main sweep, but the test differs from
the sweep's attractor test: it does not use the fixed-point finder, does not check
stability and does not cluster states; it asks only whether integrating the same
system from different points of the injury line ends, at $t = 12{,}000$, in states
more than 20 burden units apart. A line that is still moving at $t = 12{,}000$ is
classified by where it is at that time. Where the figure says ``attractor'' it refers
to these final states.

Applied to the data, the criterion gives the following (Table~\ref{tab:fig01-counts}).
No knockout line has a low-burden basin: the largest spread of a knockout line is
$0.056$ burden units in the sterile branch and $0.00024$ in the pathogenic branch,
i.e.\ under knockout the 21 initial conditions end, to within these spreads, in a
single final state.
% from: ko_lines_sterile_gap_gt_cut 0, ko_lines_pathogenic_gap_gt_cut 0,
% ko_lines_sterile_gap_max 0.05583, ko_lines_pathogenic_gap_max 0.000238
With efferocytosis ON, $2{,}537$ of $10{,}500$ sterile (unit, $K_g$) pairs (24.2\%; from 629
of $1{,}500$ units) and $3{,}294$ of $63{,}000$ pathogenic pairs (5.2\%; from 836 of
$9{,}000$ units) have a low-burden basin; the fraction depends on $K_g$
(Table~\ref{tab:fig01-kg}). The median spread of the qualifying pairs is 33.9
(sterile) and 25.7 (pathogenic) burden units.
% from: systems_*_gap_gt_cut, units_*_with_gap_any_Kg, gap_median_*_kept, frac_*_gap_gt_cut_KgNN

In the sterile branch the low-burden final state of every qualifying pair has
the higher epithelial health ($100\%$). In the pathogenic branch the total burden mixes
epithelial damage and pathogen load: among the $3{,}294$ qualifying pathogenic pairs
the low-burden final state has the higher $E_h$ in 87.5\%, the lower $P_{LP}$ in
49.4\%, and both a higher $E_h$ and a $P_{LP}$ that is not higher in 47.2\%.
% from: sterile_kept_lower_higherEh 1, pathogenic_kept_lower_higherEh 0.8749,
% pathogenic_kept_lower_lowerPLP 0.4942, pathogenic_kept_lower_higherEh_and_noHigherPLP 0.4721
``Low-burden'' in the pathogenic panels therefore means lower total burden, which in
about half of the pairs is obtained with more pathogen than the high-burden state.

\begin{table}[H]
\centering
\caption{Counts at every step of Fig.~2. A system is a unit under one arm; a
(unit, $K_g$) pair is a unit at one $K_g$ with efferocytosis ON. ``Spread'' is
$\max_e b_e - \min_e b_e$ over the line's finite final states.}
\label{tab:fig01-counts}
\begin{tabular}{lrrr}
\toprule
 & sterile & pathogenic & total \\
\midrule
units drawn & 1{,}500 & 9{,}000 & 10{,}500 \\
integrations ($\text{units}\times 8\times 21$) & 252{,}000 & 1{,}512{,}000 & 1{,}764{,}000 \\
\quad did not reach $t = 12{,}000$ & 0 & 2 & 2 \\
\quad non-finite final state & 0 & 3 & 3 \\
final states used & 252{,}000 & 1{,}511{,}995 & 1{,}763{,}995 \\
(unit, $K_g$) pairs, efferocytosis ON & 10{,}500 & 63{,}000 & 73{,}500 \\
\quad with spread $> 20$ (C, D, E) & 2{,}537 & 3{,}294 & 5{,}831 \\
\quad units with spread $> 20$ at $\geq 1$ $K_g$ & 629 & 836 & 1{,}465 \\
knockout systems & 1{,}500 & 9{,}000 & 10{,}500 \\
\quad with spread $> 20$ & 0 & 0 & 0 \\
example candidates (A, B) & 2{,}537 & -- & \\
(unit, $e$, $K_g$) pairs in F, G, production draw & -- & 122{,}891 & \\
\quad with the additional draws (Section~\ref{sec:fig1-FG}) & -- & 180{,}515 & \\
\bottomrule
\end{tabular}
% from: Fig_2_numbers.csv (rows_*, systems_*, units_*, ko_lines_*,
% example_candidates_sterile, FG_pairs). 122,891 = 836 x 7 x 21 - 1 (one intact final
% state of a qualifying unit is among the five dropped); 180,515 = 1,228 x 7 x 21 - 1.
\end{table}

\begin{table}[H]
\centering
\caption{Fraction of (unit, $K_g$) pairs whose efferocytosis-ON line has a low-burden basin (spread $> 20$), by
$K_g$, ungated (1{,}500 sterile and 9{,}000 pathogenic units at every $K_g$). Not
plotted in Fig.~2; it gives the composition of panels C--E by $K_g$.}
\label{tab:fig01-kg}
\begin{tabular}{lrrrrrrr}
\toprule
$K_g$ & 15 & 20 & 25 & 30 & 35 & 40 & 45 \\
\midrule
sterile (\%)    & 30.0 & 35.4 & 35.7 & 31.5 & 21.8 & 11.0 & 3.7 \\
pathogenic (\%) & 1.0 & 2.8 & 5.0 & 6.8 & 8.0 & 7.5 & 5.5 \\
\bottomrule
\end{tabular}
% from: frac_sterile_gap_gt_cut_KgNN, frac_pathogenic_gap_gt_cut_KgNN
\end{table}

\subsection{Panels A and B: one example system}
\label{sec:fig1-AB}

\emph{Choice.} Candidates are sterile (unit, $K_g$) pairs whose efferocytosis-ON line has a
low-burden basin and whose knockout line has none; there are 2{,}537, i.e.\ every
sterile qualifying pair.
The candidate with the largest spread of final burdens is shown: kinetic grid 5,
row 1321, $\beta_r = 0.0065$, $K_g = 20$, $r_P = 0$ (spread 74.8 burden units).
% from: example_candidates_sterile, ex_batch, ex_row_ind, ex_beta_ros, ex_Kg, ex_spread_comp
Its kinetic parameters are listed in \texttt{source\_data/Fig\_2\_numbers.csv}
(\texttt{ex\_par\_*}).

\emph{Integration.} The example is re-integrated by \texttt{figure\_scripts/Fig\_2.R} under
both arms (intact at $K_g = 20$, and knockout), with the solver settings above but
only to $t = 4{,}000$, with output every 2 time units; panel A uses the 8 initial
conditions $e \in \{2, 16, 30, 44, 58, 72, 86, 97\}$, panel B the 80 equally spaced
$e \in [1, 99]$. The horizon is therefore shorter than the $t = 12{,}000$ of the stored
data. For this example the difference is immaterial: the intact final states at
$t = 4{,}000$ ($E_h = 89.30$ and $14.48$) equal the stored $t = 12{,}000$ values to
the printed precision, the knockout final state differs by $0.009$ in $E_h$
($12.285$ against $12.275$), and over the last 100 time units $|dE_h/dt| \le 4\times10^{-6}$
on every panel-A trajectory.
% from: ex_t4000_loE_Eh, ex_t4000_hiE_Eh, ex_t12000_on_Eh_max/min, ex_t4000_ofE_Eh,
% ex_t12000_ko_Eh_min/max, ex_t4000_A_max_dEh_dt_last100 3.6e-6

\emph{Classification.} A panel-A/B integration is assigned to the low-burden basin
if its final $E_h$ exceeds the threshold $\Theta$, the midpoint between the smallest
and the largest final $E_h$ of the 16 panel-A integrations (8 intact, 8 knockout):
$\Theta = (12.28 + 89.30)/2 = 50.79$.
% from: ex_t4000_THR, ex_t4000_ends_Eh_min/max
In panel A, 4 of 8 intact integrations ($e = 58, 72, 86, 97$) end above $\Theta$ and
0 of 8 knockout integrations do. In panel B, 37 of the 80 intact integrations
(all $e \ge 54.34$) end above $\Theta$ and 43 below, with a single switch along the line;
the boundary drawn is the midpoint between the last $e$ below and the first above,
$e = 53.72$. All 80 knockout integrations end below $\Theta$.
% from: ex_t4000_A_on_n_lower 4, ex_t4000_A_off_n_lower 0, ex_B_on_n_lower 37,
% ex_B_on_n_higher 43, ex_B_on_monotone TRUE, ex_B_on_boundary 53.72, ex_B_off_n_lower 0
The stored line of the same system ($t = 12{,}000$, 21 initial conditions) switches
between $e = 49.5$ and $e = 54.45$, consistent with this boundary.
% from: ex_t12000_on_last_higher_Eh0, ex_t12000_on_first_lower_Eh0

\begin{table}[H]
\centering
\caption{Panel A/B example (grid 5, row 1321; $\beta_r = 0.0065$, $K_g = 20$,
sterile): final states of the panel-A integrations at $t = 4{,}000$.}
\label{tab:fig01-example}
\begin{tabular}{llrr}
\toprule
marker in A & integration & $E_h$ & $M_1$ \\
\midrule
low-burden, efferocytosis ON  & intact, $e = 97$   & 89.30 & 1.15 \\
high-burden, efferocytosis ON & intact, $e = 2$    & 14.48 & 56.70 \\
high-burden, efferocytosis OFF & knockout, $e = 2$ & 12.28 & 68.54 \\
\bottomrule
\end{tabular}
% from: ex_t4000_loE_*, ex_t4000_hiE_*, ex_t4000_ofE_*
\end{table}

\subsection{Panels C, D, E: where the final states lie}
\label{sec:fig1-CDE}

\emph{Population.} Every (unit, $K_g$) pair whose efferocytosis-ON line has a
low-burden basin (Eq.~\ref{eq:fig01-gap}), with the same 20-unit rule in both
branches: 2{,}537 sterile pairs from 629 units (C) and 3{,}294 pathogenic pairs from
836 units (D, E). No further filter is applied.
% from: panelC_n_systems, panelC_n_units, panelD_n_systems, panelD_n_units

\emph{Plotted points.} Each pair contributes three points: its low-burden ON final
state (``ON low-burden''), its high-burden ON final state (``ON high-burden''), and the
knockout outcome of the same unit (``OFF''). The knockout outcome is the mean of the
unit's 21 knockout final states; for the units in these panels those 21 states
differ by at most $9.3\times10^{-4}$ (sterile) and $1.5\times10^{-4}$ (pathogenic) in
$E_h$, so the mean is the common final state. Because a unit can qualify at several
$K_g$, its knockout point is repeated once per qualifying $K_g$, so each group has the
same $n$ as the number of pairs. Every qualifying unit has a knockout line, so no
pair lacks its OFF point.
% from: panelC_ko_range_Eh_max 0.00093, panelD_ko_range_Eh_max 0.00015, panel*_off_missing 0
Panel C (sterile) and D (pathogenic) are in the $(E_h, M_1)$ plane, panel E
(pathogenic) in the $(E_h, P_{LP})$ plane.

\emph{Axes.} The window of each panel spans the 0.5\% to 99.5\% quantiles of all three
groups' coordinates, widened by 6\% of that span on each side and clipped to
$[0, 100]$ ($E_h$, $M_1$, $P_{LP}$). This gives $E_h \in [0, 100]$, $M_1 \in [0, 100]$ (C);
$E_h \in [0, 88.0]$, $M_1 \in [0, 100]$ (D); $E_h \in [0, 88.0]$, $P_{LP} \in [0, 94.3]$ (E).
% from: Fig_2_panelCDE.csv xlim_*, ylim_*

\emph{Density bands.} Each group is drawn as nested highest-density regions of a
bivariate Gaussian kernel density estimate (\texttt{MASS::kde2d}, default
bandwidth, $120\times120$ grid over the window extended by 10\% on each side). The
region enclosing a fraction $p$ of the group is bounded by the density contour at
the $(1-p)$ quantile of the estimated density evaluated at the group's own points,
for $p = 0.90, 0.70, 0.50, 0.25$ (lighter to darker fill); the $p = 0.90$ contour is
also outlined. Because each group's bands are defined by its own mass, groups of
different spread are comparable.

\emph{Markers.} Diamonds mark density modes, not medians. On the same kernel surface,
candidate peaks are interior grid cells not lower than any of their 8 neighbors, with height
at least 5\% of the highest peak and more than 0.12 apart (Euclidean distance in
coordinates scaled by the window width in each axis) from every higher retained
peak. Every point of the group is assigned to its nearest candidate peak (same scaled
distance), and the share of points so assigned is the peak's \emph{mass}. Peaks with
mass $\ge 0.25$ are retained; if none reaches 0.25 the most massive peak alone is
retained. Retained peaks are ranked by mass; the first is the filled diamond, a
second (if any) the open diamond; at most two are drawn. Arrows join the main OFF
peak to the main peak of each ON group, are trimmed to the marker edge, and are
omitted when shorter than 2.4 times that trim. In panels C and D one peak is retained
per group; in three of these six groups (C OFF, C ON high-burden, D ON high-burden)
no peak reaches mass 0.25 and the retained peak is the most massive one. In E the OFF
and ON high-burden groups have two retained peaks each (Table~\ref{tab:fig01-CDE}).
The main OFF and ON high-burden peaks of E lie close together near $P_{LP} = 0$. The masses are those of a nearest-peak partition of the
points and depend on the number of candidate peaks (7--14 per group in C, D; 2--8 in E).
% from: Fig_2_panelCDE.csv peak*_*, n_peak_candidates, n_peaks_kept, n_peaks_drawn

\begin{table}[H]
\centering
\caption{Panels C--E: the three groups of final states. $n$ = number of (unit, $K_g$) pairs
(each OFF point repeated once per qualifying $K_g$). Median and interquartile range
(IQR) of the coordinates; main (and second) density peak with its assigned mass.
Peak coordinates are kernel-grid cell centers.}
\label{tab:fig01-CDE}
\scriptsize
\setlength{\tabcolsep}{3pt}
\begin{tabular}{llrrrrrll}
\toprule
panel & group & $n$ & \multicolumn{2}{c}{$E_h$} & \multicolumn{2}{c}{$y$} & main peak & 2nd peak \\
 & & & median & IQR & median & IQR & $(E_h, y)$, mass & $(E_h, y)$, mass \\
\midrule
C ($y = M_1$) & OFF          & 2{,}537 & 22.1 & 11.7--37.7 & 72.9 & 40.9--92.2 & (12.2, 94.9), 0.24 & -- \\
              & ON high      & 2{,}537 & 32.5 & 18.0--49.4 & 43.1 & 21.6--69.4 & (55.5, 11.2), 0.16 & -- \\
              & ON low       & 2{,}537 & 71.1 & 55.2--86.2 & 6.0 & 2.4--14.2 & (92.9, 2.1), 0.27 & -- \\
D ($y = M_1$) & OFF          & 3{,}294 & 27.2 & 15.0--42.4 & 75.0 & 47.3--92.2 & (21.4, 94.9), 0.28 & -- \\
              & ON high      & 3{,}294 & 29.3 & 19.1--41.0 & 34.8 & 10.6--65.1 & (26.7, 7.1), 0.20 & -- \\
              & ON low       & 3{,}294 & 48.5 & 36.7--60.2 & 17.2 & 9.4--32.9 & (58.6, 9.2), 0.32 & -- \\
E ($y = P_{LP}$) & OFF       & 3{,}294 & 27.2 & 15.0--42.4 & 0.08 & 0.00--5.19 & (11.6, 0.08), 0.54 & (47.1, 9.6), 0.46 \\
              & ON high      & 3{,}294 & 29.3 & 19.1--41.0 & 4.81 & 0.00--67.6 & (16.9, 0.08), 0.41 & (31.1, 68.6), 0.28 \\
              & ON low       & 3{,}294 & 48.5 & 36.7--60.2 & 9.19 & 1.64--24.4 & (48.9, 2.9), 0.28 & -- \\
\bottomrule
\end{tabular}
% from: source_data/Fig_2_panelCDE.csv (all columns), rounded
\end{table}

\subsection{Panels F and G: pathogen load with and without efferocytosis}
\label{sec:fig1-FG}

\emph{Population.} Pathogenic units with a low-burden basin at one or more $K_g$:
1{,}228 units, 836 from the production draw and 392 from two additional draws made
for these panels only (next paragraph). For these units all seven intact $K_g$ arms
and all initial conditions are used, including the $K_g$ values at which the unit has
no low-burden basin. Sterile units are excluded ($P_{LP} \equiv 0$).
% from: FG_units 1228, FG_units_production 836, FG_units_extra 392

\emph{Additional draws for panels F and G.} In the production draw the fraction of
pathogenic units with a low-burden basin is highest at small $r_P$ (18.3\% at
$r_P = 0.08$) and lowest at large $r_P$ (4.7\% at $0.40$, 4.8\% at $0.48$), which leaves 3
to 39 such units per ($\bros$, $r_P$) cell. To fill
the thin cells, two further batches of pathogenic units were drawn
(\texttt{pipeline/06b\_effero\_fg\_extra\_units.R}) and integrated exactly as above
(\texttt{pipeline/06\_effero\_knockout.R} with \texttt{UNITS\_FILE}; same arms, initial
conditions, solver and horizon). For each of the 54 cells, the fraction of units with
a basin was estimated from the production draw as $(m+1)/(n+2)$, with $m$ units with a
basin among $n$ drawn; draws were then allocated one at a time to the cell with the
smallest expected number of units with a basin, counting the draws of earlier batches.
Within a cell, kinetic vectors were drawn uniformly without replacement from the grid
rows the cell did not yet contain, so the old and new units of a cell are exchangeable
random draws from the grid. Batch 1 (seed 2): 4{,}000 units, all integrated (4 of
672{,}000 integrations did not reach $t = 12{,}000$). Batch 2 (seed 3): 2{,}600 units
planned, of which 2{,}300 were integrated before the run's time limit (one non-finite
final state dropped). Of these 6{,}300 units, 392 have a low-burden basin (251 and
141), and with them every cell holds 13 to 48 units with a basin. The additional units
enter panels F and G only; every other number of this section and of
Table~\ref{tab:fig01-counts} rests on the production draw.
% from: FG_units_extra_drawn 6300, FG_units_extra 392; logs_pipeline/06_effero_fg_extra{,2}.log

\emph{Paired quantity.} For a unit, an initial condition $e$ and an intact $K_g$,
\begin{equation}
  \Delta P_{LP} = \max\!\big(P_{LP}^{\mathrm{ON}, K_g}(e),\, 0.1\big)
                - \max\!\big(P_{LP}^{\mathrm{OFF}}(e),\, 0.1\big),
  \label{eq:fig01-dplp}
\end{equation}
where both terms are final states at $t = 12{,}000$ from the same initial condition
and ON/OFF denote intact/knockout. The floor of 0.1 makes a pair in which both arms
end without pathogen an exact tie. There are 180{,}515 such pairs ($25{,}787$ or
$25{,}788$ per $K_g$); 23.5\% are exact ties, 76.5\% positive (more lamina-propria
pathogen with efferocytosis intact) and 0.08\% negative; the pooled median is 7.26.
Ties are included in every statistic.
% from: FG_pairs, FG_pairs_per_Kg_min/max, FG_frac_ties 0.2346, FG_frac_positive 0.7646,
% FG_frac_negative 0.000825, FG_median_all 7.263

\emph{Statistic.} At each level $x$ of the driver on the horizontal axis (panel F:
$\beta_r$, pooling all $r_P > 0$; panel G: $r_P$, pooling all $\beta_r$), the
other driver defines strata (panel F: the six $r_P$ levels; panel G: the nine
$\beta_r$ levels). The selection for bistability leaves the ($\beta_r$, $r_P$) cells
unequal in size (13 to 48 units with the additional draws), so every cell is
\emph{subsampled} to the same number of units, $k = 13$, the size of the smallest cell.
A subsample of a level draws $k$ units without replacement from each of its strata
(F: $6k = 78$ units; G: $9k = 117$ units), keeps all $7\times21$ ($K_g$, $e$) pairs of
each drawn unit, and takes the plain median of $\Delta P_{LP}$ over the pooled pairs.
Every level of the other driver, every $K_g$ and every initial condition is thereby
equally represented, without weights. The plotted value $S(x)$ is the median of these
medians over 1{,}000 independent subsamples, so that it does not depend on one
particular draw. For example, the point of panel G at $r_P = 0.40$ is obtained as follows.
\begin{enumerate}
\item \emph{One draw.} Pick 13 combinations at random from each of the nine $\bros$
  levels at $r_P = 0.40$: $13 \times 9 = 117$ combinations. Each combination has
  $7\ K_g \times 21$ initial conditions $= 147$ values of $\Delta P_{LP}$, so the draw
  holds $117 \times 147 = 17{,}199$ values. Their median is one number.
\item \emph{Repeat.} Step 1 is done 1{,}000 times, each time with a new random pick of
  13 combinations per $\bros$ level, which gives 1{,}000 medians.
\item \emph{The plotted point.} The point is the median of those 1{,}000 medians.
\end{enumerate}
A point of panel F is obtained in the same way from 13 combinations at each of the six
$r_P$ levels at that $\bros$ ($13 \times 6 = 78$ combinations per draw). Only 13 to 48
combinations exist per ($\bros$, $r_P$) cell, so the draws overlap; repeating them
adds no data and only removes the dependence on one random pick. Every unit carries all $7\times21$ pairs (one pair of the production
draw is missing, see Table~\ref{tab:fig01-counts}). The script also offers the
stratum-weighted median that uses every unit (\texttt{FG\_STAT=balanced}; it agrees with
$S(x)$ to within 0.1 at every level) and the earlier size-weighted median over $K_g$ of
the per-$K_g$ median (\texttt{FG\_STAT=kg}).
% from: source_data/Fig_2_panelFG.csv (med, n_min_stratum); numbers FG_subsample_k 13; figure_scripts/Fig_2.R, sub_band()

\emph{Uncertainty.} The band is a 95\% percentile interval from 400 subsamples drawn
as above but with replacement within each stratum ($k = 13$ units per stratum), each
reduced to the plain median of its pooled pairs; the 2.5\% and 97.5\% quantiles are
reported. Units, not pairs, are drawn because the pairs of one unit share a parameter
vector. The interval therefore describes the median of a balanced sample of 13 units
per cell. The random seed is set to 1 once, before the nine $\beta_r$ levels (in
increasing order), which are followed by the six $r_P$ levels. No test is performed.
The subsampled median computed within each $K_g$ separately (200 subsamples each,
seed 2) spreads widely (Table~\ref{tab:fig01-FG}, last column); the figure does not
show this spread.

\begin{table}[H]
\centering
\caption{Panels F and G: plotted subsampled median $S$ of $\Delta P_{LP}$ (ON $-$ OFF;
13 units per stratum) with 95\% interval, number of units available, size of the
smallest stratum (units) and number of (unit, $e$, $K_g$) pairs per level, and the
range of $S$ computed within each of the seven $K_g$. Production and additional units
pooled.}
\label{tab:fig01-FG}
\begin{tabular}{llrrrrrr}
\toprule
panel & level & units & smallest stratum & pairs & $S$ & 95\% CI & per-$K_g$ $S$ \\
\midrule
F & $\beta_r = 0.0015$ & 146 & 13 & 21{,}461 & 18.86 & 13.67--27.07 & 8.29--42.79 \\
  & $0.0025$ & 142 & 16 & 20{,}874 & 11.93 & 8.54--16.58 & 4.87--39.57 \\
  & $0.0035$ & 156 & 20 & 22{,}932 & 8.80 & 5.20--13.47 & 4.13--34.29 \\
  & $0.0045$ & 120 & 18 & 17{,}640 & 8.27 & 4.93--11.16 & 0.80--33.92 \\
  & $0.0055$ & 131 & 15 & 19{,}257 & 8.25 & 4.77--11.90 & 2.27--28.05 \\
  & $0.0065$ & 141 & 17 & 20{,}727 & 6.26 & 3.09--12.07 & 0.96--22.82 \\
  & $0.0075$ & 126 & 16 & 18{,}522 & 5.24 & 2.71--7.58 & 0.81--14.41 \\
  & $0.0085$ & 136 & 16 & 19{,}992 & 2.29 & 0.52--5.96 & 0.15--8.18 \\
  & $0.0095$ & 130 & 14 & 19{,}110 & 1.59 & 0.31--6.28 & 0.06--9.03 \\
G & $r_P = 0.08$ & 275 & 23 & 40{,}425 & 2.53 & 0.25--4.57 & 0.00--8.94 \\
  & $0.16$ & 193 & 19 & 28{,}370 & 4.49 & 1.55--7.56 & 0.32--20.58 \\
  & $0.24$ & 177 & 14 & 26{,}019 & 8.15 & 6.10--10.28 & 1.59--27.92 \\
  & $0.32$ & 227 & 16 & 33{,}369 & 7.78 & 4.64--10.81 & 2.26--25.74 \\
  & $0.40$ & 194 & 13 & 28{,}518 & 10.30 & 7.31--13.93 & 5.82--35.81 \\
  & $0.48$ & 162 & 14 & 23{,}814 & 12.95 & 10.03--15.75 & 5.94--38.97 \\
\bottomrule
\end{tabular}
% from: source_data/Fig_2_panelFG.csv (n_units, n_min_stratum, n_pairs, med, lo, hi, perKg_min, perKg_max), rounded
\end{table}

\subsection{Reproduction}
%% ---- from docs/methods_parts/fig01.tex, lines 425-430 (the sentence on a second draw is moved to the assembly report)
The data are produced by
\texttt{N\_PER=1500 SEED=1 Rscript pipeline/06\_effero\_knockout.R}
(1{,}764{,}000 integrations; about 19.5 h on 10 cores, about 195 core-hours, in the
original run). The
additional units of panels F and G are drawn by
\texttt{N\_EXTRA=4000 SEED=2 Rscript pipeline/06b\_effero\_fg\_extra\_units.R} and
\texttt{N\_EXTRA=2600 SEED=3 UID0=200000 PRIOR\_UNITS=./data/effero\_fg\_extra\_units.rds}
\texttt{OUT=./data/effero\_fg\_extra2\_units.rds Rscript pipeline/06b\_effero\_fg\_extra\_units.R},
and integrated by \texttt{pipeline/06\_effero\_knockout.R} with \texttt{UNITS\_FILE} set
to each unit table (\texttt{OUT\_STEM=./data/effero\_fg\_extra} and
\texttt{./data/effero\_fg\_extra2}; \texttt{N\_CORES=6}, \texttt{CHUNK\_SIZE=100},
\texttt{DEADLINE\_H} as the time limit; about 10.3 min per 100 units on 6 cores).
The figure, together with every number in this section, is produced by
\texttt{Rscript figure\_scripts/Fig\_2.R} (under a minute).

\subsection{Legend}

%% ---- from docs/methods_parts/fig01.tex, lines 433-469
\paragraph{Fig.~2.}
\textbf{Removing efferocytosis removes the low-burden final state.}
\textbf{(A)} Trajectories in the $(E_h, M_1)$ plane of one sterile system (kinetic
grid 5, row 1321; $\beta_r = 0.0065$, $K_g = 20$) from eight initial conditions on the
injury line ($E_h(0) = 2, 16, 30, 44, 58, 72, 86, 97$; black crosses on the horizontal
axis), integrated to $t = 4{,}000$ with efferocytosis intact (solid, ON) and knocked
out ($g \equiv 0$; dashed, OFF). Arrowheads give the direction of motion. Line color
gives the final state: blue, final $E_h$ above $\Theta = 50.8$ (lower burden); red
(ON) or gray (OFF), below. Filled circles mark the final states of the integrations
from $E_h(0) = 97$ (ON) and $E_h(0) = 2$ (ON and OFF). With efferocytosis OFF all eight
integrations end at the same high-burden state ($E_h = 12.3$, $M_1 = 68.5$); with
efferocytosis ON the four integrations from $E_h(0) \ge 58$ end at a low-burden state
($E_h = 89.3$, $M_1 = 1.1$) and the other four at a high-burden state ($E_h = 14.5$,
$M_1 = 56.7$).
\textbf{(B)} Final state of the same system against the initial condition, for 80
equally spaced $E_h(0) \in [1, 99]$: upper row ON, lower row OFF. With efferocytosis
ON, initial conditions above $E_h(0) = 53.7$ (white mark) end in the low-burden
state; with efferocytosis OFF every initial condition ends in the high-burden state.
\textbf{(C--E)} Selection, identical for the sterile and the pathogenic branch: every (unit, $K_g$)
pair whose 21 efferocytosis-ON final states at $t = 12{,}000$ span more than 20
burden units: (C) sterile, 2{,}537 pairs from 629 units, $(E_h, M_1)$; (D)
pathogenic, 3{,}294 pairs from 836 units, $(E_h, M_1)$; (E) the same pathogenic pairs
in $(E_h, P_{LP})$. Per pair: the minimum-burden ON final state (blue, ``ON low-burden''), the maximum-burden ON final state (red, ``ON
high-burden'') and the knockout final state of the same unit (gray, ``OFF''). Shaded
bands are highest-density regions enclosing 90, 70, 50 and 25\% of each group;
diamonds mark the most massive density peak of each group (filled) and, where it
holds at least a quarter of the group, a second peak (open). Arrows run from the
main OFF peak to the two main ON peaks and are omitted where shorter than the markers.
``Attractor'' in the key denotes these final states; no fixed points were computed. No knockout system in the experiment has final
states spanning more than 20 burden units.
\textbf{(F, G)} Pathogenic units with a low-burden final state at one or more $K_g$
(1{,}228: 836 from the production draw and 392 from additional draws aimed at the
sparsely filled ($\bros$, $r_P$) cells): median paired change in lamina-propria pathogen at the final state, $\Delta
P_{LP} = \max(P_{LP}^{\mathrm{ON}}, 0.1) - \max(P_{LP}^{\mathrm{OFF}}, 0.1)$, same unit
and initial condition, against (F) $\beta_r$ (in units of $10^{-3}$) and (G)
$r_P$. Every ($\bros$, $r_P$) cell is subsampled to 13 units, so every level of the pooled
driver ($r_P$ in F, $\beta_r$ in G) is equally represented; points are the median over
1{,}000 subsamples. Bands: 95\% intervals from 400 subsamples drawn with replacement. Dotted line: $\Delta P_{LP} = 0$.

%% ===========================================================================
\section{Generation of Fig.~3A, 3B and Extended Data Fig.~2A, 2B: Tregs shift both attractors and widen the range of recoverable injury}
\label{sec:fig3}
%% TODO: Methods text for the new Figure 3 (figure_scripts/Fig_3.R). Panels and populations
%% are those of the former Supplementary Figure S01 bistable/ON_mono phase portraits.
%% TODO: the shared design of these phase portraits (integration, coloring, example
%% rule) is still documented in Section~\ref{sec:figS01}, together with the example
%% combinations of Table~\ref{tab:core-s01-ex} (rows Figure 3 A and Figure 3 B) and the
%% legend paragraph for Figure 3 kept in that section's Legend subsection; move what
%% belongs here when this section is written.

%% ===========================================================================

\emph{Note.} The shared design of these phase portraits -- the integration, the
colouring and the rule that picks the example combination -- is documented in
\ref{sec:figS01} (Extended Data Fig.~1), together with the example combinations of
Table~\ref{tab:core-s01-ex} (its rows for Fig.~3A and Fig.~3B) and the legend
paragraph kept in that section. The panels and populations of Fig.~3A, 3B and of
Extended Data Fig.~2A, 2B are those of the phase portraits described there; this
section is a placeholder until that material is moved here.

\section{Generation of Fig.~3C and Extended Data Fig.~2C: load balance during the injury, bistable combinations}
\label{sec:fig4}

%% ===========================================================================
%% fig02_03.tex -- Methods for Figure 4, Supplementary Figure S02 and Figure 5.
%% Every number below was recomputed on 2026-10-01 from ./data/ by
%%   pipeline/20_load_balance.R   (trajectories, full-grid summaries)
%%   pipeline/24_numbers_load_balance_boundary.R -> source_data/load_balance_numbers.csv, source_data/boundary_numbers.csv
%%   figure_scripts/Supplementary_Fig_1.R               -> source_data/Supplementary_Fig_1_example*.csv
%% unless a comment says otherwise.
%% ===========================================================================

%% ---- from docs/methods_parts/fig02_03.tex, lines 13-216
\subsection{Quantities}

All dynamics are those of the 13-state model of Section~\ref{M-sec:model} of the Methods, with
efferocytosis intact and the two conditions differing only in $\theta$
(Tregs OFF: $\theta=0$; Tregs ON: $\theta=1$). Let
$L = E_a + E_n$ be the load of dead epithelial cells. Summing the $E_a$ and $E_n$
equations, the secondary-necrosis terms $\mu_{E_a}E_a$ cancel and, because $M_0+M_1$ is
conserved,
\begin{equation}
  \frac{dL}{dt} \;=\;
  \underbrace{E_h\Big(\tau_r R + \tau_{P,B}\,\frac{P_{LP}}{K_{LP,P}}\Big)}_{\text{killing}}
  \;-\;
  \underbrace{g(L)\,(M_0+M_1)\,L}_{\text{clearance}},
  \label{eq:lb}
\end{equation}
with $g(L)$ of Eq.~\ref{M-eq:g} of the Methods ($a_0 = 0.0045$, $a_{\mathrm{floor}} = 0.0002$, $n_g = 4$). The pathogen term is
identically zero on the sterile branch. $K_g$ is the load at which $g$ is halfway between
$a_{\mathrm{floor}}$ and $a_0$; we call $L > K_g$ the regime of overwhelmed efferocytosis
(low $g$) and $L < K_g$ that of functional efferocytosis (high $g$). A trajectory
\emph{resolves} if $g(L)/a_0 \ge 0.95$ at some output time; with the constants above this is
$L \le 0.4847\,K_g$ (from $(1+(L/K_g)^4)^{-1} \ge (0.95-a_{\mathrm{floor}}/a_0)/(1-a_{\mathrm{floor}}/a_0)$).
``Resolves'' is therefore a statement about the path, not about the final state: a
trajectory that dips below $0.4847\,K_g$ and returns to the overwhelmed regime counts as
resolving (one such case is listed below). The final-state counterpart, $g(L(T))/a_0\ge0.95$ with $T = 12{,}000$,
is reported separately in Table~\ref{tab:lb}. A trajectory is \emph{at rest} if
$|dL/dt| < 10^{-3}$ at every output time in $[11{,}900,\,12{,}000]$.

\subsection{Population}

Combinations are taken from the sterile and the pathogenic help branches of the
rescue classification (files \texttt{rescue\_<branch>.rds}, produced by the gate and
classification of Sections~\ref{M-sec:gate} and~\ref{M-sec:classes} of the Methods; one row per combination of
kinetic vector, $K_g$, $\beta_r$ and $r_P$). On the initial-condition line
$y(e)$ of Eq.~\ref{M-eq:icline} of the Methods, each condition has a threshold
$E_h(0)^*$ below which the tissue reaches its high-burden attractor; the rescue gap is
$\gap = E_h(0)^*_{\mathrm{OFF}}-E_h(0)^*_{\mathrm{ON}}$ (Eq.~\ref{M-eq:gap} of the Methods). Bistable combinations have a threshold
under both conditions; ON\_mono combinations have one under
Tregs OFF only, the Tregs-ON system reaching the low-burden attractor from every probed
$e$. The selection is applied step by step (Table~\ref{tab:lbgate}):
\begin{enumerate}
\item every combination of the rescue file, joined to its kinetic vector;
\item \emph{worst-probe deviation}: the largest deviation, over the probes of the
  initial-condition line, of a probe's final state from the fixed point it was assigned
  to must be at most $10^{-2}$ under both conditions (the maximum, not the median that the
  sweep's trust flag uses: a row with one probe far from every listed fixed point otherwise
  passes);
\item bistable: both thresholds increasing in $e$, gap not censored, and
  $\gap \ge \Delta_{\min}$ with $\Delta_{\min} = 5$ (sterile) or 2 (pathogenic help);
  ON\_mono: Tregs-OFF threshold increasing and located;
\item 500 combinations drawn at random per class (R \texttt{sample}, seed 1, bistable then
  ON\_mono), or all if fewer;
\item of those, the combinations of the Treg-help class, so that the two rows of each figure
  are gated alike: the pathogenic branch is Treg-help by construction, the sterile gate is not,
  and the sterile draw is reduced from 500 to 442 (bistable) and 453 (ON\_mono). Because the
  draw of step 4 is uniform over the gated class, what remains is a uniform draw from its
  Treg-help part.
\end{enumerate}
%% from: pipeline/24_numbers_load_balance_boundary.R, block "gate"
\begin{table}[H]
\centering
\caption{Selection behind Fig.~3C, Extended Data Fig.~2C and Extended Data Fig.~3 (and the eligibility of Fig.~5): combinations remaining after each step, starting from the gated classes of Table~\ref{M-tab:gate} of the Methods.}
\label{tab:lbgate}
\begin{tabular}{lrrrr}
\toprule
 & \multicolumn{2}{c}{sterile} & \multicolumn{2}{c}{pathogenic help}\\
step & bistable & ON\_mono & bistable & ON\_mono\\
\midrule
rescue file & 78{,}169 & 11{,}159 & 16{,}225 & 1{,}270\\
worst-probe deviation $\le 10^{-2}$, both conditions & 78{,}057 & 11{,}124 & 16{,}197 & 854\\
threshold direction(s) increasing & 78{,}051 & 11{,}124 & 13{,}701 & 854\\
gap not censored / threshold located & 78{,}051 & 11{,}124 & 13{,}701 & 854\\
$\gap \ge \Delta_{\min}$ (5 sterile, 2 pathogenic) & 606 & -- & 180 & --\\
of these, Treg-help class & 541 & 9{,}869 & 180 & 854\\
drawn (seed 1) & 500 & 500 & 180 & 500\\
integrated to $t = 12{,}000$, both conditions & 500 & 500 & 180 & 500\\
Treg-help among those integrated (drawn and summarized) & 442 & 453 & 180 & 500\\
\bottomrule
\end{tabular}
\end{table}
The rescue-gap threshold restricts the bistable class to the upper tail of the gap
distribution (Section~\ref{M-sec:gap} of the Methods); in the pathogenic branch it leaves 180
combinations, all of which are used. No integration failed: all 3{,}360 integrated
trajectories reached $t=12{,}000$; the 3{,}150 of the Treg-help class are the ones summarized
and drawn.

\subsection{Initial condition, solver and horizon}

Each combination is integrated once per condition from a single point of the
initial-condition line, the midpoint of its rescued interval,
$e_{\mathrm{mid}} = (E_h(0)^*_{\mathrm{OFF}}+E_h(0)^*_{\mathrm{ON}})/2$ for bistable combinations
and $e_{\mathrm{mid}} = E_h(0)^*_{\mathrm{OFF}}/2$ (midway between total injury and the
Tregs-OFF threshold) for ON\_mono. By construction the Tregs-OFF system starts in its
high-burden basin and the Tregs-ON system in its low-burden basin, up to the resolution
of the thresholds. Integration (Table~\ref{M-tab:solver} of the Methods): \texttt{deSolve::ode} with \texttt{lsodar}, absolute and
relative tolerances $10^{-14}$ and $10^{-8}$, maximum step 5, at most $4\times10^{5}$ steps
per output interval, the pathogen extinction root (both $P_L$ and $P_{LP}$ set to 0 when
$\max(P_L,P_{LP})$ falls below $10^{-11}$) and the non-negativity clamp of the sweep;
output every 0.5 time units on $[0, 12{,}000]$. Every per-trajectory quantity of
Table~\ref{tab:lb} is computed on this full 0.5 grid.

\subsection{Per-trajectory definitions}

On the output grid $t_k$: $L_0 = L(0)$; $L_{\min} = \min_k L(t_k)$ and its time
$t_{\min}$ (the global minimum over $[0,12{,}000]$); the load removed before the minimum,
$(L_0 - L_{\min})/L_0$; ``killing exceeds clearance while $L>K_g$'': $dL/dt>0$ at some $t_k$
with $L(t_k) > K_g$; the peak of $M_1$ while $L>K_g$, $\max\{M_1(t_k): L(t_k)>K_g\}$;
``$L$ reaches $K_g$'': $L(t_k) \le K_g$ for some $k$; the minimum of $dL/dt$ over the whole
trajectory; ``ends at $L>K_g$'': $L(12{,}000)>K_g$; the final and the peak value of
$P_{LP}/K_{LP,P}$. The paired comparisons use, for each combination, the time
$t_{\min}^{\mathrm{OFF}}$ of the Tregs-OFF minimum: $M_1$ and $L$ of both conditions are read
at that time.

\subsection{Endpoint statistics}

%% from: pipeline/24_numbers_load_balance_boundary.R, block "table" (source_data/load_balance_numbers.csv)
\begin{table}[H]
\centering
\caption{The load balance from the midpoint injury, per class and branch, one trajectory per
combination and condition. Counts unless stated; medians over combinations. Computed on the
full 0.5 output grid. Columns are the four branch-and-class combinations drawn in Fig.~3C, Extended Data Fig.~2C and Extended Data Fig.~3.}
\label{tab:lb}
\small
\resizebox{\linewidth}{!}{%
\begin{tabular}{lrrrr}
\toprule
 & \multicolumn{2}{c}{sterile} & \multicolumn{2}{c}{pathogenic help}\\
 & bistable & ON\_mono & bistable & ON\_mono\\
\midrule
$n$ (combinations) & 442 & 453 & 180 & 500\\
starts at $L>K_g$ & 442 & 453 & 180 & 500\\
$L_0/K_g$, median (range) & 3.12 (1.93--5.97) & 3.45 (1.98--5.69) & 2.27 (1.11--3.20) & 2.46 (1.76--4.27)\\
\multicolumn{5}{l}{\emph{Tregs OFF}}\\
$L_{\min} > K_g$ & 441 & 453 & 169 & 498\\
$L_{\min}/K_g$, median & 2.07 & 2.50 & 1.82 & 2.05\\
load removed before the minimum, median & 33.2\% & 26.0\% & 19.5\% & 14.5\%\\
$t_{\min}$, median & 20.5 & 16.0 & 11.8 & 9.0\\
killing $>$ clearance while $L>K_g$ & 442 & 453 & 180 & 499\\
ends at $L>K_g$ & 442 & 453 & 180 & 498\\
$L(T)/K_g$, median (IQR) & 10.13 (7.65--12.87) & 9.48 (7.13--12.20) & 6.80 (6.03--7.81) & 7.79 (6.40--10.34)\\
resolves ($g/a_0\ge0.95$ at some time) & 0 & 0 & 1 & 2\\
at rest at $t=12{,}000$ & 442 & 453 & 180 & 500\\
\multicolumn{5}{l}{\emph{Tregs ON against Tregs OFF, same combination}}\\
$L/K_g$ at $t_{\min}^{\mathrm{OFF}}$, median ON / OFF & 2.00 / 2.07 & 2.48 / 2.50 & 1.79 / 1.82 & 2.03 / 2.05\\
$M_1$ at $t_{\min}^{\mathrm{OFF}}$, median ON / OFF & 13.6 / 19.1 & 12.9 / 15.4 & 25.6 / 31.0 & 8.1 / 8.7\\
lower $M_1$ at $t_{\min}^{\mathrm{OFF}}$ with Tregs ON & 442 & 453 & 180 & 499\\
peak $M_1$ while $L>K_g$, median OFF / ON & 68.7 / 14.3 & 42.3 / 15.1 & 80.9 / 33.0 & 67.9 / 27.8\\
lower peak $M_1$ while $L>K_g$ with Tregs ON & 442 & 453 & 180 & 500\\
deeper minimum of $dL/dt$ with Tregs ON & 442 & 453 & 180 & 499\\
$\min dL/dt$ ratio ON/OFF, median & 2.37 & 2.59 & 2.41 & 2.68\\
\multicolumn{5}{l}{\emph{Tregs ON}}\\
killing $>$ clearance while $L>K_g$ & 101 & 310 & 91 & 433\\
$L$ reaches $K_g$ & 442 & 453 & 180 & 497\\
resolves ($g/a_0\ge0.95$ at some time) & 442 & 453 & 177 & 497\\
ends with $g/a_0 \ge 0.95$ & 442 & 453 & 177 & 497\\
ends at $L<K_g$ & 442 & 453 & 180 & 497\\
$L(T)/K_g$, median & 0.080 & 0.028 & 0.153 & 0.118\\
time to $g/a_0 \ge 0.95$, median & 54.5 & 130.5 & 53.0 & 244.0\\
at rest at $t=12{,}000$ & 442 & 453 & 180 & 500\\
\multicolumn{5}{l}{\emph{Pathogen, $P_{LP}/K_{LP,P}$}}\\
OFF final $<10^{-4}$ / $<10^{-2}$ & & & 71.7\% / 85.0\% & 78.0\% / 95.2\%\\
ON final $>10^{-2}$ / $>10^{-1}$ & & & 66.1\% / 5.0\% & 59.8\% / 28.6\%\\
final, median OFF / ON & & & 0 / 0.028 & 0 / 0.039\\
peak, median OFF / ON & & & 0.261 / 0.261 & 0.375 / 0.399\\
\bottomrule
\end{tabular}}
\end{table}

\emph{Exceptions.} All 895 sterile trajectory pairs follow the pattern
``Tregs OFF ends at rest in the overwhelmed regime, Tregs ON resolves and ends at rest with
$g/a_0\ge0.95$''. In the pathogenic branch nine trajectories do not
(\texttt{pipeline/24\_numbers\_load\_balance\_boundary.R}, block \texttt{exceptions}):
(i) three bistable Tregs-ON trajectories cross $K_g$ and come to rest at
$L = 0.50$--$0.62\,K_g$ ($g/a_0 = 0.88$--$0.94$): below $K_g$, but short of the 0.95 criterion;
(ii) three ON\_mono Tregs-ON trajectories never reach $K_g$ ($\min L/K_g = 2.3$--$3.2$) and come
to rest at $L = 3.8$--$7.7\,K_g$, i.e.\ the Tregs-ON system has a high-burden attractor
reachable from $e_{\mathrm{mid}}$ although it was classified ON\_mono;
(iii) one bistable Tregs-OFF trajectory dips to $L = 0.456\,K_g$ (so it ``resolves'' by the
path criterion) and returns to rest at $L=16.3\,K_g$;
(iv) two ON\_mono Tregs-OFF trajectories resolve and end at $L = 0.18$ and $0.24\,K_g$. Both have
$E_h(0)^*_{\mathrm{OFF}} = 0.039$, the midpoint of the first bisection bracket $[0, 0.078]$, so
$e_{\mathrm{mid}} = 0.0195$ lies within the threshold's resolution and need not be on the
high-burden side.

\emph{ON\_mono is a statement about the initial-condition line} (Section~\ref{M-sec:classes} of the Methods):
the grid- and endpoint-seeded finder finds two stable Tregs-ON fixed points in 763 of
11{,}159 sterile and 581 of 1{,}270 pathogenic help ON\_mono combinations (column
\texttt{n\_field\_stable\_on} of \texttt{selected\_big.rds}). ON\_mono therefore means that no
probed initial condition of the line reaches a high-burden attractor under Tregs ON, not
that such an attractor does not exist; exception (ii) above is a case in which one is
reached from $e_{\mathrm{mid}}$, a point between the probes.

\subsection{Figure construction}

Paths are drawn in the plane of $L/K_g$ (logarithmic axis) against $dL/dt$
(Eq.~\ref{eq:lb}, cells per time unit). Fig.~3C and Extended Data Fig.~2C shows the bistable class and
Extended Data Fig.~3 the ON\_mono class. The two arms are overlaid in one
panel per branch (dashed, Tregs OFF; solid, Tregs ON): the sterile branch in
Fig.~3C and Extended Data Fig.~3A, the pathogenic branch in Extended Data
Fig.~2C and Extended Data Fig.~3B. Per class and branch, 20 combinations are drawn at random from the
Treg-help combinations of the integrated population (seed 1, re-set per branch), the same combinations in both columns.
For drawing only, the stored series is thinned by motion: a point is kept when
$\log_{10}(L/K_g)$ has moved by more than 0.004 or $dL/dt$ by more than 0.02 since the last kept
point, and in any case every 10 time units, so that fast passages are drawn at full
resolution; the last output point is always kept. Points with $L/K_g \le 10^{-3}$ are not
drawn (this truncates the tail of 49 sterile and 2 pathogenic stored series of the 3{,}150, all
Tregs ON). Circles mark the initial condition; triangles the last stored point ($t = 12{,}000$),
filled when the trajectory is at rest (all of them are). Sterile paths and markers are colored
by $M_1$ (linear, 0--100); pathogenic paths and markers by $P_{LP}/K_{LP,P}$ (logarithmic,
$10^{-4}$--1, values below $10^{-4}$ drawn at $10^{-4}$); one color scale per branch, shared by
Tregs OFF and ON. The shaded region is $L>K_g$; the solid vertical line is $L = K_g$ and the dashed
horizontal line $dL/dt = 0$. The $x$ range is set per panel by the data drawn; the $y$ range is
shared by the four panels of a figure. No statistic is computed on the drawn subset; all numbers
are those of Table~\ref{tab:lb}.

\subsection{Legend}

%% ---- from docs/methods_parts/fig02_03.tex, lines 218-232
\paragraph{Fig.~3C and Extended Data Fig.~2C.}
The load balance during the injury, bistable combinations. Each path is one parameter
combination integrated from the midpoint of its rescued interval,
$(E_h(0)^*_{\mathrm{OFF}}+E_h(0)^*_{\mathrm{ON}})/2$, plotted as the dead-cell load
$(E_a+E_n)/K_g$ (logarithmic) against its rate of change, epithelial killing minus clearance
(Eq.~\ref{eq:lb}). \textbf{(A)} Sterile, Tregs OFF; \textbf{(B)} sterile, Tregs ON, the same 20
combinations; paths and markers colored by activated macrophages $M_1$. \textbf{(C)} Pathogenic
help, Tregs OFF; \textbf{(D)} pathogenic, Tregs ON, the same 20 combinations; paths and
markers colored by $P_{LP}/K_{LP,P}$. Circles: initial condition; triangles: state at
$t=12{,}000$, filled when at rest. Shaded: $E_a+E_n > K_g$ (overwhelmed efferocytosis, low $g$);
arrows label the overwhelmed and functional regimes. Under Tregs OFF the load first falls,
turns while still above $K_g$ and rises to rest in the overwhelmed regime; under Tregs ON it
crosses $K_g$, the balance becomes strongly negative as clearance recovers, and the path comes to
rest far below $K_g$. Populations: 442 sterile and 180 pathogenic bistable combinations, all of the Treg-help class
(Table~\ref{tab:lb}); 20 per branch drawn.

%% ===========================================================================
\section{Generation of Fig.~4: damage-matched comparison of Tregs OFF and Tregs ON}
\label{sec:fig6}

The experiment asks what a brake on inflammation buys when its cost in epithelial damage is
held fixed: at the same cumulative epithelial damage, can the system with regulatory T cells
(Tregs) run a stronger ROS response than the system without, and what does that do to the
pathogen? All trajectories are solutions of the 13-state model of Section~\ref{M-sec:model} of the Methods; the
arms differ only in the Treg induction switch, $\theta = 0$ (Tregs OFF) or $1$
(Tregs ON), and, where stated, in the ROS production rate $\beta_r$. Nothing is interpolated
across parameter sets, fitted or imputed: every quantity is read from a simulated trajectory.
Code: \texttt{pipeline/43\_damage\_matched\_gated\_units.R} (the sample),
\texttt{pipeline/40\_damage\_matched.R} with \texttt{pipeline/cluster/damage\_matched\_gated.sh}
(the experiment), \texttt{pipeline/41\_damage\_matched\_trajectories.R} (time courses for
panels B and C), \texttt{pipeline/42\_first\_peak.R} (the first peak of each arm),
\texttt{pipeline/48\_gated\_fourth\_arm.R} (the fourth arm),
\texttt{lib/fig\_damage\_matched.R} (the filter), \texttt{lib/fig\_demo\_case.R} (panel D),
\texttt{figure\_scripts/Fig\_4.R} and \texttt{figure\_scripts/Extended\_Data\_Fig\_5.R}.

\subsection{Units of analysis and sampling design}

\paragraph{Combinations and units.} A \emph{combination} is a kinetic parameter vector (one
row of the sampled parameter grids of Section~\ref{M-sec:sampling} of the Methods, identified by \texttt{tau\_id})
together with an efferocytosis half-saturation load $K_g \in \{15, 20, \dots, 45\}$, a
pathogen growth rate $r_P \in \{0.08, 0.16, \dots, 0.48\}$ and a
ROS production rate $\beta_r \in \{0.0015, 0.0025, \dots, 0.0095\}$, the nine levels of the
parameter sweep. A \emph{unit} is a combination together with an initial condition
$E_h(0) \in \{20, 50, 80\}$; the remaining states start as in the sweep
($E_n(0) = K_E - E_h(0)$ with $K_E = 100$, $M_0(0) = 100$, $T_0(0) = 100$, $P_L(0) = 500$,
$P_{LP}(0) = 10$, all others zero). The reference level of a unit is its own $\beta_r$, written
$\beta_{\mathrm{OFF}}$ where the distinction matters.

\paragraph{Population.} The population is every pathogenic combination that passed the
bistability gate of the Tregs OFF system (Sections~\ref{M-sec:gate} and \ref{M-sec:classes} of the Methods): 57{,}497 combinations in
11{,}955 ($\texttt{tau\_id}, K_g, r_P$) cells over 1{,}035 parameter-grid rows.
% from: pipeline/checks/fig06_methods_numbers.R, section 1
No condition is placed on what Tregs do to the combination: combinations in which Tregs
are neutral or harmful at their own $\beta_r$ (by the steady-state criterion of
Section~\ref{M-sec:classes} of the Methods) are included, and the label is recorded only to describe the
sample. With three
initial conditions this gives 172{,}491 possible units.

\paragraph{The persistence gate.} The question only arises where the system without Tregs
does not clear the infection, so a unit enters the sample only if the Tregs OFF system is
predicted to leave the pathogen at a positive level. The prediction is read from the sweep's
fixed points: the attractor reached from $E_h(0)$ is the low-burden one if $E_h(0)$ lies at or
above the Tregs OFF threshold $E_h(0)^*_{\mathrm{OFF}}$ and the high-burden one otherwise, and the
unit is eligible if that attractor has $P_{LP} > 10^{-3}$. 111{,}055 units (64.4\%) in
10{,}269 cells are eligible. The prediction was verified on every unit run: the
% from: pipeline/checks/fig06_methods_numbers.R, section 1
experiment re-tests it on its own Tregs OFF trajectory and records a unit whose pathogen is
eliminated after all without searching (status \texttt{off\_eliminates}); none of the 52{,}010
units run was of this kind, and the realized final $P_{LP}$ of the Tregs OFF run exceeds
$10^{-3}$ in all of them (quartiles 1.37, 3.93, 9.56).
% from: pipeline/checks/fig06_methods_numbers.R, section 2

\paragraph{Sample.} 75{,}000 eligible units were drawn uniformly at random without replacement
(R \texttt{sample.int}, seed 1, on the eligible list ordered by \texttt{tau\_id}, $K_g$,
$r_P$, $\beta_r$, $E_h(0)$) and stored in the order drawn, so that any prefix of the file
is itself a uniform sample. They cover 10{,}005 cells over 914 grid rows. By the Treg
effect of their own combination (Section~\ref{M-sec:classes} of the Methods) they are 48{,}036 neutral, 19{,}249
help and 7{,}715 harm; by $r_P = 0.08, \dots, 0.48$: 26{,}055, 17{,}041, 11{,}866, 8{,}749,
6{,}474, 4{,}815; by $K_g = 15, \dots, 45$: 1{,}999, 5{,}542, 9{,}459, 13{,}685, 16{,}614, 16{,}423,
11{,}278; by $E_h(0) = 20, 50, 80$: 18{,}492, 25{,}834, 30{,}674. The experiment was run as
% from: pipeline/checks/fig06_methods_numbers.R, section 1
400 array tasks of 187--188 units each (unit $i$ in task $(i-1) \bmod 400$) with an eight-hour
limit; every task saved each unit as it finished and every task reached the limit, leaving
52{,}010 units (69.3\%) complete. Because the units are in random order, the completed set is a
uniform subsample of the draw; all counts below refer to it.

\paragraph{Integration.} Every run is integrated with \texttt{deSolve::lsodar} (absolute
tolerance $10^{-14}$, relative tolerance $10^{-8}$, maximum step 5, at most $10^5$ internal
steps) from $t = 0$ to $T_{\max} = 12{,}000$, with the pathogen extinction floor of
Section~\ref{M-sec:units} of the Methods ($P_L$ and $P_{LP}$ set to zero by a root event once both fall below
$10^{-11}$). Output is recorded at $t = 0$ and at 400 times log-spaced between 0.05 and
12{,}000; every integral is a trapezoidal sum over these output times.

\paragraph{Four arms.} (i) OFF($\beta_r$): Tregs OFF at the reference level; its cumulative
damage is the cost to be matched. (ii) ON($\beta_r$): Tregs ON at the same level, ``the brake
alone''; it isolates the effect of Tregs at fixed aggression.
(iii) ON($\beta_r^{*}$): Tregs ON at the level $\beta_r^{*}$ found by the search below, at which
its cumulative damage equals that of OFF($\beta_r$). (iv) OFF($\beta_r^{*}$): Tregs OFF at the
same $\beta_r^{*}$, ``the aggression alone''; it is integrated afterwards for the units in which
ON($\beta_r^{*}$) eliminates the pathogen and for panel D, and shows what the stronger response
would cost without the brake.

\subsection{The cost that is matched}

Damage is the epithelial deficit $K_E - E_h(t)$ and the cost of a run over a window $[0, T]$ is
$D(T) = \int_0^{T} (K_E - E_h(t))\,dt$. A run has settled at its own convergence time
$t_{\mathrm{conv}}$, the first output time beyond which every one of the 13 states stays within
1\% of its final value, or within $0.01$ units for states near zero: within
$\max(0.01\,|x_i(T_{\max})|,\ 0.01)$ of its value at $T_{\max}$. The matching window of a pair is
$[0, t_{\mathrm{match}}]$ with
$t_{\mathrm{match}} = \max\{t_{\mathrm{conv}}[\mathrm{OFF}(\beta_r)],\ t_{\mathrm{conv}}[\mathrm{ON}(\beta)]\}$,
the later of the two settling times: the episode ends once both runs have settled, and both are
charged over the same window.
Because $t_{\mathrm{conv}}$ of the Tregs ON run depends on the $\beta$ under evaluation, the
window is recomputed at every evaluation of the search and the reported match is that of the
final pair on its own window. In the pairs of the figure (below) the window has median 648
(IQR 263--1{,}170; range 67.2--12{,}000), $t_{\mathrm{conv}}$ of OFF($\beta_r$) has median 248, and
ON($\beta_r^{*}$) sets the window in 76.6\% of pairs. The relative cost difference of a
% from: pipeline/checks/fig06_methods_numbers.R, section 4
Tregs ON run at $\beta$ is
\[
r(\beta) = \frac{D_{\mathrm{ON}(\beta)}(t_{\mathrm{match}}) - D_{\mathrm{OFF}(\beta_r)}(t_{\mathrm{match}})}{D_{\mathrm{OFF}(\beta_r)}(t_{\mathrm{match}})} .
\]

\subsection{The search for $\beta_r^{*}$}

The cost is not monotone in $\beta$: raising $\beta$ kills more epithelium through ROS but also
more pathogen, and where the pathogen drives the damage a higher $\beta$ can lower the cost.
$\beta_r^{*}$ is therefore defined as the \emph{largest} level at which the Tregs ON cost
crosses the Tregs OFF cost from below (\texttt{pipeline/40\_damage\_matched.R}): above it, every level tested costs Tregs ON more
damage than Tregs OFF. The procedure is
\begin{enumerate}
\item \emph{Scan.} $r(\beta)$ is evaluated on the geometric grid $\beta_r \cdot 1.03^k$ over
  $[0.0005, 10]$ ($\beta_r$ is a grid point).
\item \emph{Probe.} Every local minimum of $r$ on the grid that lies in $(0, 0.02)$
  receives six extra evaluations between its neighbors, so that a dip below the target
  narrower than a grid step is not missed.
\item \emph{Bracket and bisect.} The last scanned point (grid or probe) with $r \le 0$ and the next one bracket the
  crossing; bisection on $\log\beta$ (at most 30 halvings) targets $|r| \le 10^{-4}$ and accepts
  $|r| \le 10^{-3}$ (\texttt{TOL}). If no point has $r \le 0$ the unit is \texttt{no\_match\_lo}
  (Tregs ON is costlier at every level); if the last point does, \texttt{no\_match\_hi}; if
  bisection ends above \texttt{TOL}, the bracket has closed on a discontinuity of the cost (a
  basin boundary) and the unit is \texttt{no\_match\_jump}.
\item \emph{Jump test.} $r$ is evaluated at $\beta_r^{*} \times \{0.97, 0.99, 1.01, 1.03\}$ and the
  change of $r$ across $\pm 3\%$, in percentage points, is recorded (\texttt{jump\_pts}), as is
  the number of upward crossings on the scanned profile (\texttt{n\_cross\_up}).
\end{enumerate}
The status is \texttt{matched\_same} if $|\beta_r^{*}/\beta_r - 1| < 10^{-9}$, \texttt{matched\_up} if
$\beta_r^{*} > \beta_r$ and \texttt{matched\_down} otherwise. Each unit took a median of 343 evaluations (IQR 341--343, range 334--390). Of the 52{,}010 units, 46{,}371 are \texttt{matched\_up}, 3{,}875 \texttt{matched\_same},
% from: pipeline/checks/fig06_methods_numbers.R, section 2
27 \texttt{matched\_down}, 969 \texttt{no\_match\_lo} and 768 \texttt{no\_match\_jump}; no unit is
\texttt{no\_match\_hi}. The achieved error $|r(\beta_r^{*})|$ has median $3.9\times10^{-5}$ and
99th percentile $1.0\times10^{-4}$ (maximum $9.8\times10^{-4}$); in absolute terms the two
damages differ by a median 0.87 units on a median target of 24{,}500. The search range
% from: pipeline/checks/fig06_methods_numbers.R, section 2
extends far beyond the sweep: $\beta_r^{*}$ exceeds the largest sweep level (0.0095) in 18{,}756
matched units (maximum 0.676).

\paragraph{The direction of $\beta_r^{*}$.} Whether $\beta_r^{*}$ lies above or below $\beta_r$ is
decided by the brake alone: if ON($\beta_r$) causes less damage than OFF($\beta_r$) there is
damage to spend and $\beta_r^{*} > \beta_r$; if it causes more, the deficit can only be made up by
lowering $\beta$, or by raising it far enough to kill the pathogen. ON($\beta_r$) causes less
damage over the run than OFF($\beta_r$) in 45{,}635 units (87.7\%). Of the 6{,}375 units in which
it causes more, 5{,}032 are nevertheless matched upward (the second route), 965 find no match
at all (\texttt{no\_match\_lo}: 965 of the 969), and 23 are matched below $\beta_r$; the 27
\texttt{matched\_down} units are all neutral-labeled, and 24 of them have $\beta_r^{*}$ within
3\% of $\beta_r$.
% from: pipeline/checks/fig06_methods_numbers.R, section 2

\subsection{Selection, with counts}
\label{sec:dm_filter}

The pairs of Fig.~4 satisfy, in this order (\texttt{lib/fig\_damage\_matched.R}):
\begin{enumerate}
\item a matched status with $|r(\beta_r^{*})| \le 10^{-4}$ (\texttt{MATCH\_TOL}): 50{,}273 units are
  matched, of which 49{,}757 reached this target of the bisection;
\item a robust match: no jump at $\beta_r^{*}$ (\texttt{jump\_pts} $\le 10$) and a single upward
  crossing of the cost profile (\texttt{n\_cross\_up} $= 1$), so that $\beta_r^{*}$ does not depend on
  which branch of a discontinuous profile the search landed on (355 and 574 of the 49{,}757
  fail these, 915 at least one; 48{,}842 remain);
\item \emph{the brake alone saves at least 5\% of the cumulative damage}:
  $D_{\mathrm{ON}(\beta_r)}(T_B) \le 0.95\, D_{\mathrm{OFF}(\beta_r)}(T_B)$
  (\texttt{DAM\_MIN} $=0.05$), where
  $T_B = \max\{t_{\mathrm{conv}}[\mathrm{OFF}(\beta_r)], t_{\mathrm{conv}}[\mathrm{ON}(\beta_r)]\}$
  is the convergence window of that pair, so the criterion is read over a window of the same
  construction as the match it selects for. Both arms of all 48{,}842 robust pairs are
  re-integrated to their own $T_B$ for this
  (\texttt{pipeline/44\_window\_damage.R}; $T_B$ median 447, quartiles 240 and 1{,}133, set by
  ON($\beta_r$) in 56.4\% of pairs). This is the condition under which the brake creates headroom;
  without it a pair can be matched only by lowering $\beta$ or by the pathogen-killing route, and
  the comparison is no longer about the headroom.
\end{enumerate}
8{,}789 pairs remain, and every panel of Fig.~4 and Extended Data Fig.~5 uses all of them.
In all of them $\beta_r^{*} > \beta_r$, the smallest ratio being 1.053.
% from: pipeline/checks/fig06_methods_numbers.R, section 3
By the Treg effect at the pair's own $\beta_r$ they are 7{,}204 help
(median $\beta_r^{*}/\beta_r$ 1.86), 658 harm (2.39) and 927 neutral (1.67); the label is
recorded and reported but not used for selection.

\paragraph{What the floor replaces.} Criterion 3 stands in place of two earlier ones, a timing
criterion (``Tregs act within the run'': the $M_1$ of ON($\beta_r$) departs from that of
OFF($\beta_r$) by more than 1\% before $T_{\max}$) and a ratio cut ($\beta_r^{*} \ge 1.05\,\beta_r$,
imposed to drop pairs whose $\beta_r^{*}$ was indistinguishable from $\beta_r$). Both are implied by
the floor and neither is needed. The timing criterion was an effect size in disguise: of the
15{,}272 pairs it alone removed, the brake saved a median 0.04\% of the cumulative damage
(maximum 0.66\%) and 15{,}271 were neutral-labeled at their own $\beta_r$. Their Treg channel is
inert rather than their challenge mild --- their median $K_g$, $r_P$, $\beta_r$ and $E_h(0)$ equal
those of the pairs that pass, and their cumulative damage and peak load without Tregs are within
10\% --- but $S$ decays 7.5 times faster ($\tau_S$ median 0.109 against 0.0145), so the loop gain
$\tau_{10}(\sigma_S/\tau_S)(\tau_{t,a}/\tau_{t,d})$ is 81 times lower. A floor on the damage the
brake saves removes them in the currency the experiment matches. The ratio cut is then redundant,
since no pair above the floor has $\beta_r^{*} < 1.05\,\beta_r$. The conclusions do not turn on the
level of the floor: at 1\%, 2\%, 5\% and 10\% the sample holds 16{,}532, 13{,}016, 8{,}789 and
6{,}142 pairs, and the pathogen integrated over the matching window is lower with Tregs in
86.6\%, 87.4\%, 89.3\% and 90.3\% of them.

\paragraph{The same floor read over the whole run.} Criterion 3 was first read from the stored
integrals over the whole run, $D_{\mathrm{ON}(\beta_r)}(T_{\max})$ against
$D_{\mathrm{OFF}(\beta_r)}(T_{\max})$, which is not the integral the match is made over. Both
versions are now available for every robust pair, and they select almost the same sample: the two
savings correlate at 0.994 (Spearman 0.989), 8{,}720 pairs pass both floors, 69 only over the
window, 923 only over the run, and 34{,}511 fail both. The samples behave alike, the window one
first: 8{,}789 pairs against 9{,}643, median $\beta_r^{*}/\beta_r$ 1.89 against 1.76, the pathogen
over the matching window lower in 89.3\% against 89.3\%, the steady-state load higher in 92.5\%
against 92.5\%, and the epithelium higher in 95.3\% against 95.2\%. The window version is the one
used, because the quantity the criterion asks about -- the damage the brake saves -- is then the
quantity the experiment spends.
% from: scratch recount (window rule, all 48,842 robust pairs)
% from: scratch recount (damage-floor audit)

\subsection{The first peak of the pathogen}
\label{sec:first_peak}

The \emph{first peak} of a run is the first strict local maximum of $P_{LP}$ on the output grid:
the first time the load turns over, whatever happens afterwards. It is computed for OFF($\beta_r$)
and ON($\beta_r^{*}$) of every pair by \texttt{pipeline/42\_first\_peak.R}, which re-integrates both
arms with the experiment's parameters, initial condition and solver on a 2{,}000-point log-spaced
grid. Both arms have one in all 8{,}789 pairs; in OFF($\beta_r$) it is reached at a median time
9.00 and has quartiles 21.3, 33.8, 53.5 in value.
% from: pipeline/checks/fig06_methods_numbers.R, section 4

No prominence rule is applied and none is needed. Over the 25{,}568 series of the sample selected
over the whole run --- 9{,}643 pairs, 8{,}720 of them also in the present 8{,}789 --- the plain
first local maximum and one required to fall 5\% below the peak before any higher value pick the
same point in 12{,}781 pairs (both arms); in the three exceptions the first maximum is a shoulder
on a climb to a higher attractor, and the comparison it enters is the same either way. The same
check on the 254 series of the 127 pairs that are new to the present sample finds no
disagreement.

\paragraph{Why not the maximum over the run.} \texttt{pipeline/40\_damage\_matched.R} also stores
$\mathrm{peak}\,P_{LP}$, the maximum over $[0, T_{\max}]$. For a run that climbs to a high-burden
attractor and stays there, that maximum \emph{is} the attractor, so comparing it with the other
arm's genuine peak compares a steady state with a peak. On that measure 18 of the 8{,}789 pairs
have more pathogen ``at the peak'' with Tregs; in all 18 the ON($\beta_r^{*}$) arm has no transient
peak at all, so its maximum is its own attractor, reached at a median $t = 2{,}231$ against 11 for
the Tregs OFF maximum. On the first peak none remain. The two measures still coincide where a run
has no transient peak at all: the first peak lies within 1\% of the final load in 44 pairs for
ON($\beta_r^{*}$) and 4 for OFF($\beta_r$).
Fig.~4 and Extended Data Fig.~5 use the first peak throughout; the divergence time of earlier versions of this
experiment (the first output time at which $M_1$ of ON($\beta_r$) departed from that of
OFF($\beta_r$) by more than 1\%) is not used anywhere, and neither are the integrals split at it.
% from: pipeline/checks/fig06_methods_numbers.R, section 4

\subsection{Trajectory subset for panels B and C}

The experiment stores per-unit summaries only. For panels B and C, 1{,}000 pairs were drawn
uniformly at random without replacement (seed 1) from the 8{,}789 pairs of the filter, and their
three runs
re-integrated with the experiment's parameters, initial condition and solver on a finer output
grid ($t = 0$ and 2{,}000 log-spaced times from 0.05 to 12{,}000). All 1{,}000 returned three
complete runs, so panels B and C are drawn on the whole draw.
The first peak of OFF($\beta_r$) is recomputed on this grid (median 8.95, IQR
6.52--12.7, range 2.84--41.9). Time is expressed as $\tau = t/t_{\mathrm{peak}}$, so that
% from: pipeline/checks/fig06_methods_numbers.R, section 4
$\tau = 1$ is that peak in every pair, and each pair's series are linearly interpolated in
$\log t$ onto a common grid of 400 log-spaced values of $\tau$ in $[0.1, 100]$, a span every
pair covers with recorded output.

\subsection{Plotted quantities}

\paragraph{Panel A (the headroom).} Per reference level $\beta_r$, the median (point) and IQR
(bar) of $\beta_r^{*}$ over the 8{,}789 pairs at that level, on logarithmic axes, against the
diagonal $\beta_r^{*} = \beta_r$; the label is the median per-pair ratio $\beta_r^{*}/\beta_r$
(Table~\ref{tab:fig04_panelB}).

\paragraph{Panel B (order of events).} For each of the 1{,}000 pairs and each $\tau$: the pathogen
$P_{LP}$ of OFF($\beta_r$) and of ON($\beta_r$), scaled jointly to $[0,1]$ by the minimum and
maximum of the two series over the drawn $\tau$ range; the Treg level $T_r$ of ON($\beta_r$)
and the suppression $M_1^{\mathrm{OFF}} - M_1^{\mathrm{ON}}$ at the same $\beta_r$, each scaled to
$[0,1]$ by its own range. Lines are medians and ribbons the IQR over the 1{,}000 pairs at each
$\tau$. No series in this panel is a magnitude. At $\tau = 1$, the suppression has reached a
median 0.12\% of its eventual maximum (90th percentile 1.9\%) and exceeds a tenth of it in 20 of
the 1{,}000 pairs: the peak forms before the brake takes hold.
% from: pipeline/checks/fig06_methods_numbers.R, section 4

\paragraph{Panel C (pathogen ratio).} For each of the 1{,}000 pairs and each $\tau$, the ratios
$P_{LP}^{\mathrm{ON}(\beta_r)}/P_{LP}^{\mathrm{OFF}}$ and $P_{LP}^{\mathrm{ON}(\beta_r^{*})}/P_{LP}^{\mathrm{OFF}}$,
each pathogen value floored at $10^{-12}$. The median and quartiles over pairs are taken of the
log ratio and back-transformed; the horizontal line marks 1. The median ratio of ON($\beta_r^{*}$)
is 0.982 at $\tau = 0.5$, 0.836 at $\tau = 1$, reaches its minimum 0.430 at $\tau = 2.18$, and is
0.782, 0.907, 1.076 and 1.094 at $\tau = 5$, 10, 50 and 100; that of ON($\beta_r$) is 1.000 at
$\tau = 1$ and 1.108, 1.364, 1.804 and 1.881 at $\tau = 5$, 10, 50 and 100.
% from: pipeline/checks/fig06_methods_numbers.R, section 5

\paragraph{Panel D (the demonstration unit).} One unit (index 38614 of the sample:
\texttt{tau\_id} 31555, $K_g = 45$, $r_P = 0.08$, $E_h(0) = 20$, $\beta_r = 0.0075$,
$\beta_r^{*} = 0.0279$, i.e.\ $3.72\,\beta_r$; harm-labeled at its own $\beta_r$), one of the 15 pairs
of the 8{,}789 in which only ON($\beta_r^{*}$) eliminates, chosen for lying at the medians of that
set ($\beta_r^{*}/\beta_r$ 3.72 against a median of 3.56; cost of the aggression without the brake
1.60 against a median of 1.62). Its four arms
are re-integrated on a 3{,}000-point log-spaced grid and $M_1$, $P_{LP}$ and $E_h$ are drawn
against time; the first peak of OFF($\beta_r$) at $t = 15.8$ and $t_{\mathrm{match}} = 3{,}257$ are marked. Final
values: $P_{LP}$ 22.9 (OFF, $\beta_r$), 65.7 (ON, $\beta_r$), 0 (ON and OFF at $\beta_r^{*}$,
eliminated at $t = 162$ and 151 on the grid of the figure; 165 for ON($\beta_r^{*}$) on the
experiment's 400-point grid); $E_h$ 39.6, 52.7, 46.6 and 3.4; $M_1$ 9.9, 2.0, 2.3 and 58.0.
The cumulative damage of OFF($\beta_r^{*}$) over $[0, t_{\mathrm{match}}]$ is 1.60 times that of OFF($\beta_r$).
% unit 38614 since 2026-11-18 (lib/fig_demo_case.R output); the previous unit, 4108, had an
% inverted Tregs-ON threshold and so was outside the 15 pairs the text describes.

No bootstrap, test or correlation is used in Fig.~4; all summaries are medians and quartiles
over pairs. The pairs span a region of parameter space whose distribution is not known, so
proportions are proportions of the sample, not frequencies in a population.

\subsection{Endpoint statistics}

\begin{table}[ht]
\centering
\caption{Panel A: $\beta_r^{*}$ per reference level, all 8{,}789 pairs of the filter.}
% from: pipeline/checks/fig06_methods_numbers.R, section 3
\label{tab:fig04_panelB}
\begin{tabular}{rrcrr}
\hline
$\beta_r$ & $n$ & $\beta_r^{*}$, median (IQR) & $\beta_r^{*}/\beta_r$, median & 90th pct. \\
\hline
0.0015 &   219 & 0.00280 (0.00206--0.00468) & 1.87 & 4.66 \\
0.0025 &   491 & 0.00497 (0.00369--0.00714) & 1.99 & 3.57 \\
0.0035 &   779 & 0.00644 (0.00515--0.00896) & 1.84 & 3.36 \\
0.0045 &   916 & 0.00832 (0.00669--0.0113) & 1.85 & 3.32 \\
0.0055 & 1{,}034 & 0.0106 (0.00810--0.0149)  & 1.92 & 3.97 \\
0.0065 & 1{,}136 & 0.0123 (0.00972--0.0172)  & 1.89 & 3.96 \\
0.0075 & 1{,}289 & 0.0143 (0.0111--0.0203)   & 1.91 & 4.62 \\
0.0085 & 1{,}402 & 0.0161 (0.0125--0.0227)   & 1.89 & 4.41 \\
0.0095 & 1{,}523 & 0.0178 (0.0138--0.0242)   & 1.87 & 4.46 \\
\hline
all & 8{,}789 & & 1.89 & 4.04 \\
\hline
\end{tabular}
\end{table}
Over all pairs $\beta_r^{*}/\beta_r$ has median 1.89, IQR 1.47--2.64 and maximum 39.8.
% from: pipeline/checks/fig06_methods_numbers.R, section 3

\paragraph{The pathogen at equal damage (all 8{,}789 pairs).}
The first peak of $P_{LP}$ is lower with Tregs at $\beta_r^{*}$ than without Tregs at $\beta_r$ in
every pair, by a median 11.8\% and by as much as 62.6\%. The integral over the matching window is
lower in 89.3\% (median reduction 6.3\%, maximum 99.6\%), and over the whole run in 14.1\%. The
pathogen is eliminated (extinction floor reached, or final $P_{LP} < 10^{-3}$)
% from: pipeline/checks/fig06_methods_numbers.R, sections 4 and 5
by ON($\beta_r^{*}$) in 15 pairs from 10 cells, and by the brake alone, ON($\beta_r$), in 47; by
construction OFF($\beta_r$) eliminates in none. In the 15, $\beta_r^{*}/\beta_r$ has quartiles 2.84,
3.56, 6.69, the residual pathogen without Tregs 3.86, 9.23, 10.8, the elimination time 260, 338,
562, and the final $E_h$ is higher with Tregs at $\beta_r^{*}$ than without Tregs by 6.0--7.3
units (quartiles).
% from: pipeline/checks/fig06_methods_numbers.R, section 7

\paragraph{The reversal pairs read over the shorter window.}
In 897 of the 8,789 pairs (10.2\%) both the integral over the matching window and the
steady-state pathogen come out higher with Tregs at $\beta_r^{*}$. The matching window is the
longer of the two convergence times, so it is generous to whichever arm settles later: a small
residual difference in $P_{LP}$, accumulated over a long tail, can outweigh a large early
reduction. In these pairs the Tregs-ON arm takes a median 8.3 times longer to settle than the
Tregs-OFF arm, and the residual gap is small in absolute terms (median 0.88 units of
$K_{LP,P} = 100$, against 0.17 units in the pairs that do not reverse).
\texttt{pipeline/45\_short\_window.R} therefore re-integrates both arms of each reversal pair,
OFF($\beta_r$) and ON($\beta_r^{*}$), and accumulates $P_{LP}$ by the trapezoid rule to
$\min\{t_{\mathrm{conv}}^{\mathrm{OFF}}, t_{\mathrm{conv}}^{\mathrm{ON}(\beta_r^{*})}\}$, the
window both arms have earned, on the same 600-point logarithmic grid used elsewhere in this
experiment. Over that window the cumulative pathogen load is lower with Tregs in 874 of the
897 pairs (97.4\%), by a median 27.8\%. The 23 exceptions are not explained by the tail: their
residual gap is markedly larger (median 1.38 against 0.88 units) and their two convergence times
are nearly equal (ratio median 1.20 against 8.65), so the early reduction is outweighed before the
long tail begins rather than by it. In all 23 the first peak is still lower with Tregs at
$\beta_r^{*}$, as it is in every pair of the set. This is a diagnostic for the reversal pairs
alone; the damage match itself, and every other quantity in this section, is read over the
matching window as defined above.
% from: pipeline/45_short_window.R

\paragraph{The fourth arm.} For each of the 15 eliminating pairs, OFF($\beta_r^{*}$) was
integrated and its cumulative damage over the pair's window compared with that of OFF($\beta_r$).
The aggression alone also eliminates the pathogen in all 15, but at a cost: 1.49--2.27 times the
damage of OFF($\beta_r$) (quartiles; above 1.2 in all 15, above 1.5 in 11), with the epithelium
ending a median 36 units lower than in OFF($\beta_r$) and 43 lower than in ON($\beta_r^{*}$). So in every one of these pairs the brake is what makes
% from: pipeline/checks/fig06_methods_numbers.R, section 7
the clearing response affordable. Of the 15, 10 are harm-labeled and 5 help at
their own $\beta_r$: by the steady-state criterion the brake alone makes most of them worse
(it eliminates in none of them and leaves more pathogen), and only together with the stronger
response it licenses does it clear the infection.

\section{Generation of Fig.~5 and Extended Data Fig.~6: Treg stability and specificity}
\label{sec:fig7}

% ============================================================================
% fig05.tex -- the imperfect Treg: Figure 7 and Supplementary Figures S05-S07.
% Code: pipeline/60_imperfect_sample.R, pipeline/61_imperfect_sweep.R
%       (cluster launcher pipeline/cluster/imperfect_sweep.sh),
%       pipeline/62_imperfect_read.R, pipeline/63_imperfect_rescue_prob.R,
%       lib/fig_imperfect_treg.R, figure_scripts/Fig_5.R, figure_scripts/Extended_Data_Fig_7.R,
%       figure_scripts/Extended_Data_Fig_6.R, figure_scripts/Supplementary_Fig_3.R.
% Every number below was recomputed on 2026-10-01 by
%       pipeline/checks/imperfect_numbers.R   (cited as [N§k], k = section of
%       its printed output) unless another source is given.
% ============================================================================

%% ---- from docs/methods_parts/fig05.tex, lines 16-442
In every other analysis the Treg is the \emph{baseline} Treg: it is fully specific,
in that induction reads the true commensal fraction of the lamina propria, and fully
stable, in that activated Tregs revert to the inactive pool at the constant rate
$\tau_{t,d}$. Fig.~5 and Extended Data Figs.~6 and~7 relax both properties. All other terms of the model are those
of Section~\ref{M-sec:model} of the Methods; only the two Treg equations change.

\paragraph{Phenotypic instability.}
With $\phi \in [0,1]$ (\texttt{phi\_instability}), the reversion rate of
activated Tregs increases with the danger signal $D$ relative to its largest
attainable steady level $\sigma_D/\tau_D$ (the source term of $dD/dt$ is at most
$\sigma_D$ because both fractions in it are at most one):
\begin{equation}
  \tau_{t,d}^{\mathrm{eff}}(t) \;=\; \tau_{t,d} \;+\; \phi\,\frac{D(t)}{\sigma_D/\tau_D}.
  \label{eq:imp-phi}
\end{equation}
$\phi = 0$ recovers the baseline Treg. The modification is deterministic. We
report the stability $1-\phi$, so that larger values mean a better Treg on both
axes.

\paragraph{Reduced antigen specificity.}
With $\eta \in [0,1]$ (\texttt{treg\_eff}), induction reads a perceived
commensal fraction $\tilde C_{\mathrm{rat}}(t)$ in place of
$C_{\mathrm{rat}} = f_C/(f_P+f_C)$ (with $f_P = P_{LP}/K_{LP,P}$,
$f_C = C_{LP}/K_{LP,C}$, and $C_{\mathrm{rat}} = 0$ when both vanish). Time is
divided into unit intervals $[k, k+1)$, $k = 0, 1, \dots$; on each interval
\begin{equation}
  \tilde C_{\mathrm{rat}}(t) \;=\; b_k\, C_{\mathrm{rat}}(t) \;+\; (1-b_k)\, u_k,
  \qquad t \in [k, k+1),
  \qquad b_k \sim \mathrm{Bernoulli}(\eta),\; u_k \sim U(0,1),
  \label{eq:imp-eta}
\end{equation}
with all $b_k$ and $u_k$ mutually independent. With probability $\eta$ the
Treg reads the true commensal share during that interval; otherwise it reads a
value drawn uniformly at random and unrelated to the tissue state. $\eta = 1$ recovers the baseline Treg; at $\eta = 0$ induction is driven entirely by the
uniform draws. In the implementation (\texttt{make\_noise\_fun} in
\texttt{lib/ODE\_SYSTEM\_v6.R}) a vector of $\lfloor T_{\max}\rfloor + 2$
pairs $(p_k, u_k)$ of independent $U(0,1)$ variates is drawn before each
integration, $b_k = \mathbf{1}[p_k < \eta]$, and the right-hand side looks up
the pair of $\lfloor t\rfloor$ through a step function
(\texttt{make\_ode\_system(noise\_fun)}). Because $p_k < 1$ always,
$\eta = 1$ runs are exactly deterministic.

\paragraph{Modified Treg equations.}
With $\theta\in\{0,1\}$ switching Treg induction off or on,
\begin{align}
  \frac{dT_0}{dt} &= -\,\theta\,\tau_{t,a}\,M_1\,\tilde C_{\mathrm{rat}}(t)\,T_0
                     \;+\; \Big(\tau_{t,d} + \phi\,\frac{D}{\sigma_D/\tau_D}\Big)\,T_r, \\
  \frac{dT_r}{dt} &= \phantom{-}\,\theta\,\tau_{t,a}\,M_1\,\tilde C_{\mathrm{rat}}(t)\,T_0
                     \;-\; \Big(\tau_{t,d} + \phi\,\frac{D}{\sigma_D/\tau_D}\Big)\,T_r ,
\end{align}
so $T_0 + T_r$ remains conserved. Every other equation is unchanged; in
particular $\tilde C_{\mathrm{rat}}$ enters only Treg induction, and the
efferocytic clearance is the load-dependent switch
$g(L)$ of Eq.~\ref{M-eq:g} of the Methods. This is the only analysis in the
paper whose dynamics are stochastic, and the only one in which $\phi$ and
$\eta$ differ from $0$ and $1$.

\subsection{Units: parameter combination and initial condition}

A \emph{unit} is one parameter combination -- a kinetic parameter vector
(identified by \texttt{tau\_id}, i.e.\ \texttt{(batch, row\_ind)} of the sampled
kinetic grid), together with the efferocytosis midpoint $K_g$, the pathogen
growth rate $r_P$ (\texttt{r\_L\_P}, applied in the lumen and the LP) and
the ROS production rate $\beta_r$ -- plus one initial condition
$E_h(0)$ on the initial-condition line. Units are drawn
(\texttt{pipeline/60\_imperfect\_sample.R}) from the gated combinations of the
sterile and pathogenic help branches, using the eligibility and the injury of
the load-balance analysis of Fig.~3C and Extended Data Fig.~2C (Section~\ref{sec:fig4}, Table~\ref{tab:lbgate}):

\begin{enumerate}
\item \emph{Fixed-point assignment.} In both conditions (Tregs OFF and Tregs ON),
  the largest per-probe deviation of the initial-condition-line endpoints
  from their assigned fixed points is at most $10^{-2}$.
\item \emph{Bistable combinations} (two or more Tregs-ON attractors on the line):
  both thresholds run in the increasing direction, the rescue gap
  $\gap$ (Eq.~\ref{M-eq:gap} of the Methods) is not censored and
  $\gap \ge \texttt{GAP\_MIN}$, with $\texttt{GAP\_MIN} = 5$ (sterile) and
  $2$ (pathogenic help).
\item \emph{ON\_mono combinations} (one Tregs-ON attractor
  from every probed initial condition): the Tregs-OFF threshold
  $E_h(0)^*_{\mathrm{OFF}}$ is located (finite) and runs in the increasing
  direction.
\end{enumerate}
The initial condition is the midpoint of the rescued interval,
\begin{equation}
  E_h(0) = e_{\mathrm{mid}} =
  \begin{cases}
    \big(E_h(0)^*_{\mathrm{ON}} + E_h(0)^*_{\mathrm{OFF}}\big)/2 & \text{bistable},\\
    E_h(0)^*_{\mathrm{OFF}}/2 & \text{ON\_mono},
  \end{cases}
\end{equation}
so that, up to the resolution of the thresholds, the Tregs-OFF system starts in
its high-burden basin and the baseline-Treg system in its low-burden basin. All
231 drawn bistable units satisfy
$E_h(0)^*_{\mathrm{ON}} < E_h(0) < E_h(0)^*_{\mathrm{OFF}}$ and all 746 ON\_mono units
have $E_h(0) = E_h(0)^*_{\mathrm{OFF}}/2$ exactly. The remaining state
% from: pipeline/checks/imperfect_numbers.R, section 2
variables follow the initial-condition line (Eq.~\ref{M-eq:icline} of the Methods): $E_n(0) = K_E - E_h(0)$, $M_0(0) = K_M = 100$,
$T_0(0) = K_T = 100$ and zero otherwise, except the pathogen inoculum
$P_L(0) = 500$, $P_{LP}(0) = 10$ on the pathogenic help branch
(\texttt{seed()} in \texttt{lib/skeleton\_defs.R}). Every eligible sterile
combination has $r_P = 0$ and every eligible pathogenic help combination
$r_P > 0$; each branch is integrated with its own seeding, so the sterile units
carry no pathogen at any time.

\paragraph{Eligible counts.}
Table~\ref{tab:imp-elig} gives the count after each eligibility step. No combination is lost in the join to the parameter table, and every
% from: pipeline/checks/imperfect_numbers.R, section 1
ON\_mono combination that passes step 1 also passes step 3; the binding
conditions are the deviation filter and the gap threshold.

\begin{table}[H]
\centering
\caption{Eligibility of units for the stability--specificity sweep (counts of parameter
combinations). The steps and counts agree with Table~\ref{tab:lbgate}, computed
independently; the last row gives the occupied $(K_g, \bros, r_P)$ strata. Source:
\texttt{pipeline/checks/imperfect\_numbers.R}, \S1.}
\label{tab:imp-elig}
\begin{tabular}{lrrrr}
\toprule
 & \multicolumn{2}{c}{sterile} & \multicolumn{2}{c}{pathogenic help}\\
 & bistable & ON\_mono & bistable & ON\_mono\\
\midrule
combinations in the branch file, criterion~(3) & 78{,}169 & 11{,}159 & 16{,}225 & 1{,}270\\
deviation $\le 10^{-2}$ in both arms  & 78{,}057 & 11{,}124 & 16{,}197 &   854\\
$\gap \ge$ \texttt{GAP\_MIN} (5; 2)       &    606 &   --   &    180 &   --  \\
both directions increasing, uncensored &   606 &   --   &    180 &   --  \\
OFF threshold located and increasing   &   --   & 11{,}124 &   --   &   854\\
\midrule
eligible                                &    606 & 11{,}124 &    180 &   854\\
strata occupied                         &     51 &     63 &     91 &   165\\
\bottomrule
\end{tabular}
\end{table}

\subsection{Stratified draw}

A stratum is a $(K_g, \beta_r, r_P)$ cell of the design, with
$K_g \in \{15, 20, \dots, 45\}$ (7 levels), $\beta_r \in
\{0.0015, 0.0025, \dots, 0.0095\}$ (9 levels) and $r_P = 0$ (sterile) or
$r_P \in \{0.08, 0.16, \dots, 0.48\}$ (pathogenic help, 6 levels). The sterile
eligible pool occupies all 63 possible strata (3 to 370 eligible combinations
per stratum, median 200) and the pathogenic help pool 189 of the 378 possible
strata (1 to 32, median 2). Within each branch and stratum the eligible
% from: pipeline/checks/imperfect_numbers.R, section 1
combinations (both modes pooled) were randomly permuted (R
\texttt{set.seed(1)}, \texttt{sample}) and the first $\min(n, 5)$ taken
($K_{\mathrm{PER}} = 5$). Bistable combinations not yet drawn were then added
until the stratum held two bistable units, where the stratum had them
(\texttt{BISTABLE\_FLOOR} $= 2$); this added 119 units over 70 strata (92
sterile, 27 pathogenic help). The allocation is therefore equal per stratum
and capped by availability, not proportional to the eligible pool, and the
bistable share is raised above its share of the pool (5\% of sterile and 17\%
of pathogenic help eligible combinations, against 25\% and 23\% of the drawn
units). The draw gives 977 units (Table~\ref{tab:imp-draw}); rerunning the
sampler reproduces the stored units exactly. Units per stratum: sterile 3 (1
stratum), 5 (11), 6 (10), 7 (41); pathogenic help 1 (74 strata), 2 (22), 3
(16), 4 (6), 5 (52), 6 (11), 7 (8).
% from: pipeline/checks/imperfect_numbers.R, section 2

\begin{table}[H]
\centering
\caption{Drawn units by branch and rescue mode, with the initial condition
$E_h(0) = e_{\mathrm{mid}}$ (quartiles).}
\label{tab:imp-draw}
\begin{tabular}{llrrrrr}
\toprule
branch & mode & units & strata & $E_h(0)$ Q1 & median & Q3\\
\midrule
sterile         & bistable & 100 &  51 & 8.39 & 16.23 & 30.14\\
sterile         & ON\_mono & 305 &  63 & 2.25 &  8.34 & 18.30\\
pathogenic help & bistable & 131 &  91 & 8.14 & 18.98 & 31.91\\
pathogenic help & ON\_mono & 441 & 165 & 4.39 & 12.17 & 23.73\\
\midrule
total           &          & 977 & 252 &      &       &      \\
\bottomrule
\end{tabular}
\end{table}

Rescue gaps of the drawn bistable units range from 5.00 to 48.59 (median 7.03)
in the sterile branch and from 2.03 to 18.44 (median 3.36) in the pathogenic help
branch. Pathogenic help units concentrate at large $K_g$ (9, 44, 59, 93, 129,
145, 93 units at $K_g = 15, \dots, 45$) and small $r_P$ (295, 162, 54, 35, 15,
11 at $r_P = 0.08, \dots, 0.48$); sterile units are spread almost evenly over
$K_g$ (43--62 per level).
% from: pipeline/checks/imperfect_numbers.R, section 2

\subsection{Simulation design}

For every unit, the Tregs-ON system ($\theta = 1$) was integrated
on the $11 \times 11$ grid $(\phi, \eta) \in \{0, 0.1, \dots, 1\}^2$ (121 cells),
each cell in $R_{\mathrm{REP}} = 10$ replicates with independent noise
realisations, and the Tregs-OFF system ($\theta = 0$, where
neither $\phi$ nor $\eta$ can act) once per unit, giving
$977 \times (121 \times 10 + 1) = 1{,}183{,}147$ planned integrations
(\texttt{pipeline/61\_imperfect\_sweep.R}). The work unit is one
\emph{item} (unit, replicate), i.e.\ the 121 cells of one replicate; at the
start of each item the generator is seeded with
$20260917 + 1000\cdot\mathrm{unit} + \mathrm{replicate}$, and the noise vectors
are drawn in a fixed order: the Tregs-OFF run first (replicate 1 only), then
the grid with $\phi$ in the outer and $\eta$ in the inner loop. Items were
distributed over Slurm array tasks on a cluster (four waves of 400 tasks,
waves 1--2 sterile, 3--4 pathogenic help).

\paragraph{Solver.}
Each run integrates the 13 states of Section~\ref{M-sec:model} of the Methods plus the
cumulative perceived signal $\int_0^t \tilde C_{\mathrm{rat}}\,ds$ (not used
further) from $t = 0$ to $T_{\max} = 12{,}000$, with output every time unit, using
\texttt{deSolve::ode} with method \texttt{lsodar}, $\mathrm{atol} = 10^{-14}$,
$\mathrm{rtol} = 10^{-8}$, at most $2\times10^{5}$ steps per output interval, the
pathogen extinction floor (a root on $\max(P_L, P_{LP}) = 10^{-11}$ that sets
both to zero), and maximum step size $h_{\max} = 1$, against $h_{\max} = 5$ in
every deterministic analysis. $h_{\max} = 1$ equals the length of the
intervals on which $\tilde C_{\mathrm{rat}}$ is constant; the jump times are not
passed to the solver as critical times, so a step can still straddle a jump.
% from: pipeline/checks/imperfect_hmax_check.R: 3 units x 2 noise
% realisations, TMAX = 3000, hmax 1 vs 0.1, max relative difference of tail L 9.1e-10
Reducing $h_{\max}$ to 0.1 changed the tail-mean dead-cell load by at most a
relative $9\times10^{-10}$ and no outcome, in a check on three units with split
cells (two noise realisations each, $T_{\max} = 3000$).
The right-hand side is evaluated at the state clamped to be non-negative.
Cost is dominated by the noise: the median solve time of one run was 125\,s at
$\eta = 0$, 48\,s at $\eta = 0.8$ and 2.1\,s at $\eta = 1$; the stored runs
account for 26,919 core-hours.
% from: pipeline/checks/imperfect_numbers.R, section 3

\paragraph{Tail statistics.}
Under $\eta < 1$ the perceived signal jumps every time unit and the trajectory
never comes to rest on a fixed point, so runs are summarized by statistics over
the second half of the output grid, $t = 6001, 6002, \dots, 12{,}000$ (the last
$\lfloor 12001/2\rfloor = 6000$ output rows): for every state $x$ the tail mean
$\bar x$, the tail (sample) variance and the tail maximum. No endpoint is
matched to a fixed point in this analysis.

\subsection{Run and unit filters}

\texttt{pipeline/62\_imperfect\_read.R} merges the per-task output. A run is
removed if the solver failed or stopped before $T_{\max}$, if $\bar E_h$ or
$\bar P_{LP}$ is not finite, if $\bar P_{LP}$ or $\bar C_{LP}$ exceeds $10^4$,
or if any of $\bar E_h$, $K_E - \bar E_h$, $\bar M_0$, $\bar M_1$, $\bar T_r$
exceeds 101 (the capacity 100 plus round-off slack; $E_a$ and $E_n$ have no
capacity and are not bounded). A unit is kept only if each of its 121 cells
retains at least one replicate and its Tregs-OFF run exists.

The merged file holds 1,176,120 Tregs-ON runs and 966 Tregs-OFF runs.
% from: pipeline/checks/imperfect_numbers.R, section 3
The missing runs are 50 complete items (50 $\times$ 121 = 6,050 Tregs-ON runs),
all in the pathogenic help waves (24 in wave 3, 26 in wave 4), spread over 33
units with at most 4 missing items per unit; every item present contains all
121 cells. Because only the filtered runs were transferred from the cluster, the
merged file cannot distinguish an item that never completed from one whose runs
were all removed by the filter; no partially filtered item exists, and every
stored run lies inside the bounds (largest tail means: $E_h$ 100,
$K_E - E_h$ 94.6, $M_0$ 100, $M_1$ 98.0, $T_r$ 92.4, $P_{LP}$ 46.4, $C_{LP}$
94.5). Eleven of the 50 missing items are replicate 1, which also carries the
Tregs-OFF run; these 11 units (all pathogenic help: 10 ON\_mono, 1 bistable)
are dropped by the Tregs-OFF requirement, although each still has all 121 cells
with 6 to 9 replicates. No unit is dropped for an incomplete grid. Of the 966
retained units, 48 were short of the full ten replicates, between them missing 68
(unit, replicate) items. The 37 items belonging to the 440 units of Fig.~5 were
re-run locally and merged, so every cell of those units now carries ten replicates;
31 items remain missing, all in units outside Fig.~5. Of the 116,886 cells of the
966 retained units, 114,224 have 10 replicates, 1,815 have 9, 605 have 8 and 242 have 7. The Tregs-OFF run enters no plotted
% from: pipeline/checks/imperfect_numbers.R, section 4
quantity; it is required only because the cell table carries it.

\subsection{Outcome: probability of rescue}

A replicate $r$ of cell $(\phi, \eta)$ of unit $i$ is \emph{rescued} when its
tail-mean dead-cell load lies below the half-saturation point of the
efferocytosis switch,
\begin{equation}
  \bar L_{i,r}(\phi,\eta) \;=\; \bar E_{a} + \bar E_{n} \;<\; K_{g,i},
\end{equation}
i.e.\ when clearance is on the high-capacity side of $g(L)$. The probability of
rescue of a cell is the fraction of its retained replicates that are rescued,
\begin{equation}
  p_i(\phi,\eta) \;=\; \frac{1}{n_i(\phi,\eta)}\sum_{r=1}^{n_i(\phi,\eta)}
  \mathbf{1}\big[\bar L_{i,r}(\phi,\eta) < K_{g,i}\big],
\end{equation}
with $n_i(\phi,\eta) \in \{6,\dots,10\}$
(\texttt{pipeline/63\_imperfect\_rescue\_prob.R}). The criterion uses the
dead-cell load only; tail pathogen is not part of it. In the pathogenic help branch
the median tail-mean $P_{LP}$ of rescued runs is 3.24 (90th percentile 13.8,
maximum 46.4), and 258,129 of 366,502 rescued runs have $\bar P_{LP} > 1$.
% from: pipeline/checks/imperfect_numbers.R, section 6
The outcome is sharply separated from the threshold: rescued runs have
$\bar L/K_g$ with median 0.117 and 99th percentile 0.406, non-rescued runs have
$\bar L/K_g \ge 1.003$ with first percentile 3.24.
% from: pipeline/checks/imperfect_numbers.R, section 6

\paragraph{Replicate spread.}
At $\eta = 1$ the replicates of a cell are identical (largest across-replicate
standard deviation of $\bar L$: $1.6\times10^{-11}$). Over the 106,260 cells with $\eta < 1$, the
across-replicate coefficient of variation of $\bar L$ has median $9.5\times10^{-4}$
and 95th percentile 0.349; restricted to cells whose replicates all agree on the
outcome it has median $8.7\times10^{-4}$ and 95th percentile $4.5\times10^{-3}$. Within a cell the noise therefore either leaves the outcome unchanged
% from: pipeline/checks/imperfect_numbers.R, section 5
or splits the replicates between the two basins: 5,788 of 116,886 cells
(4.95\%) are split ($0 < p < 1$): 883 cells at $\eta = 0$, 909 at $\eta = 0.1$, then
falling with $\eta$ to 166 at $\eta = 0.9$ and 0 at $\eta = 1$. The within-run tail
coefficient of variation of $K_E - E_h$ has median 0.0042 and 99th percentile
0.030 (maximum 6.26).
% from: pipeline/checks/imperfect_numbers.R, section 5

\paragraph{Crossing units.}
For each unit, $\bar p_i$ is the mean of $p_i$ over the 121 cells. Units with
$\bar p_i = 0$ (never rescued in any replicate of any cell) or $\bar p_i = 1$
(rescued in every replicate of every cell) carry no information about either
axis and are excluded; the remaining \emph{crossing} units ($0 < \bar p_i < 1$)
are 717 of the 966 retained units. Three further filters, described next, reduce
them to the 440 units of every panel. No other filter is applied: in particular
no settled-tail filter and no weighting enter Fig.~5, Extended Data Fig.~6 or
Extended Data Fig.~7, and every retained unit carries equal weight.
Table~\ref{tab:imp-flow} gives the flow from the draw to the panels.

\paragraph{Treg effect.}
Units are drawn from the sterile and pathogenic-help branches only, so no drawn
unit belongs to the harm class, in which Tregs raise the total burden of the
low-burden attractor. The sterile branch, however, is not split by Treg effect
at all (\texttt{BRANCHES} in \texttt{pipeline/60\_imperfect\_sample.R}), and its
eligible pool carries both labels; 36 of the 405 drawn sterile units are
labeled \emph{neutral}, meaning that the shift of the low-burden attractor is
smaller than the tolerance \texttt{EFFECT\_TOL} $= 1$ burden unit that
separates help from harm in \texttt{pipeline/03\_select.R}. Of these, 26 are
crossing units and they are excluded, so every unit of every panel is one in
which the Treg both lowers the burden of the resolved state and decides whether
the injury resolves at all. Applied alone this filter changes no reported
quantity by more than 0.012, moving $\bar P(0,0)$ from 0.081 to 0.084 and the
marginal at zero stability from 0.200 to 0.204. Applied before the crossing test instead of after it, this filter leaves 930 units,
of which 3 are never rescued, 236 always and 691 crossing, and the final set of 440
is the same; the main text states the counts in that order.

\paragraph{Minimum range.}
A unit whose marginals barely move carries no shape for either axis to be read
from, however often it is rescued. The test is on the per-unit marginals
$m^s_i$ and $m^\eta_i$, the curves the panels draw, and not on the 121 cells: a
unit rescued in every cell but one has cells spanning the whole of $[0,1]$ while
its marginals move by $1/11$, so a test on the cells cannot see it. A unit is
kept when the larger of its two marginal ranges is at least 0.5, which removes
167 of the 691 units that survive the effect filter; of those, the median
marginal range is 0.16 along stability and 0.09 along $\eta$. The larger of the
two is used rather than each separately, so that a unit whose stability curve
spans the axis is kept however flat its $\eta$ curve is. A per-axis minimum
would not be neutral between the axes: the per-unit range has median 0.98 along
stability and 0.15 along $\eta$, so requiring 0.25 on both would drop 349 of the
691, 294 of them failing on $\eta$ alone and 2 on stability alone, and would
leave a population in which the two properties are interchangeable (areas under
the two mean curves 0.482 and 0.486, against 0.281 and 0.299 before it).

\paragraph{Monotonicity.}
The grid measures how the probability of rescue declines as the two properties
are degraded away from the baseline Treg, which presupposes that the baseline
Treg is the best Treg of that unit's own grid. At $(1-\phi, \eta) = (1,1)$ the
run is deterministic, and in 716 of the 717 crossing units it is rescued, so the
baseline corner cannot be beaten; in the remaining unit it is not rescued while
other cells are. A unit can nonetheless fail the premise in the interior, where
degrading a property raises the probability of rescue, and the units in which it
does cannot be read as a decline. Units whose per-unit marginal $m^s_i$ or
$m^\eta_i$ (the means over the other axis that the marginal panels draw) falls at
any step of its axis are therefore excluded: 84 of the 524 units that survive the
two filters above.

Four of the 84 fall by more than 0.10, and all six fail in the same way. Every
cell at intermediate stability is rescued and the failures lie in the top row,
$\phi = 0$, at low $\eta$: in unit 168 the top row is not rescued for
$\eta \le 0.6$ and the other 110 cells are rescued in every replicate; units 79, 279 and
306 have the same shape over a narrower range of $\eta$ (units 179 and 292 show it
too but are removed earlier, for too small a marginal range). A
Treg that never destabilises and reads the commensal fraction at random holds
its suppression indefinitely, and the inflammation needed to clear the debris
does not develop; the same parameter combination resolves whenever $\phi > 0$,
because the danger signal returns the Treg to the inactive pool. This is the
regime in which less suppression gives a better outcome, and it is the reason
the filter is applied rather than a nuisance it removes.

The other 80 fall by small amounts, and at this replicate depth a small fall
cannot be separated from the sampling noise. A cell's probability moves in steps
of $1/n_i \ge 1/10$, and of the 1{,}139 downward steps along stability over the
717 crossing units, 1{,}128 have a split cell ($0 < p_i < 1$) at one end and the
median fall is exactly 0.1, one replicate of ten in one cell of 121. Taking each
unit's surface, forcing it monotone (iterated isotonic regression on rows and
columns) and re-drawing ten replicates per cell from it 200 times flags 162 of
the 717 on average, against the 148 observed, so the data are consistent with
the six being the only units that genuinely fail the premise. The rate at which
a monotone unit is flagged also rises with how often it is rescued, from 0.04
for units with $\bar p_i < 0.2$ to 0.62 for units with $\bar p_i > 0.8$, so the
filter removes readily rescued units preferentially and the reported surface
lies below the unfiltered one throughout: $\bar P(0,0)$ falls from 0.081 to
0.014, the marginal at zero stability from 0.200 to 0.090 and the marginal at
$\eta = 0$ from 0.246 to 0.179, while the shape of both curves is unchanged.
Setting \texttt{DROP\_NEUTRAL=0}, \texttt{MIN\_RANGE=0} and
\texttt{DROP\_NONMONO=0} (\texttt{lib/fig\_imperfect\_treg.R}) reproduces the
unfiltered surface.

% from: pipeline/checks/imperfect_numbers.R, sections 3, 6

\begin{table}[H]
\centering
\caption{From drawn units to the units in each panel. ``No OFF run'': units
lost because their replicate-1 item (which carries the Tregs-OFF run) is absent.
``Neutral'': crossing units whose Treg effect is not help. ``Flat'': units whose
larger marginal range is below 0.5. ``Falling'': units whose per-unit marginal
falls at some step of either axis.
``Used'' is the population of every panel.}
\label{tab:imp-flow}
\resizebox{\linewidth}{!}{%
\begin{tabular}{llrrrrrrrrrr}
\toprule
branch & mode & drawn & no OFF run & retained & never & always & crossing & neutral & flat & falling & used\\
\midrule
sterile         & bistable & 100 &  0 & 100 & 0 &   1 &  99 & 10 &  8 & 42 &  39\\
sterile         & ON\_mono & 305 &  0 & 305 & 0 & 103 & 202 & 16 & 44 &  4 & 138\\
pathogenic help & bistable & 131 &  1 & 130 & 0 &  28 & 102 &  0 & 54 & 18 &  30\\
pathogenic help & ON\_mono & 441 & 10 & 431 & 3 & 114 & 314 &  0 & 61 & 20 & 233\\
\midrule
sterile (Supplementary Fig.~2)          & & 405 &  0 & 405 & 0 & 104 & 301 & 26 &  52 & 46 & 177\\
pathogenic help (Supplementary Fig.~3)  & & 572 & 11 & 561 & 3 & 142 & 416 &  0 & 115 & 38 & 263\\
pooled (Fig.~5, Extended Data Fig.~7)        & & 977 & 11 & 966 & 3 & 246 & 717 & 26 & 167 & 84 & 440\\
\bottomrule
\end{tabular}}
\end{table}

At the baseline Treg ($\phi = 0$, $\eta = 1$) all 405 sterile and 557 of 561
pathogenic help units are rescued; one of the four that are not is a crossing
unit. The Tregs-OFF run itself ends below $K_g$ in 4 of 966 units (all ON\_mono
with $E_h(0) = 0.0195$, i.e.\ a Tregs-OFF threshold
$E_h(0)^*_{\mathrm{OFF}} = 0.039$); all four are rescued in every Tregs-ON cell and are therefore
excluded as ``always''.
% from: pipeline/checks/imperfect_numbers.R, section 6

\subsection{Plotted quantities}

Let $\mathcal C$ be the set of retained units of the panel's population
(pooled, sterile or pathogenic help), $s = 1 - \phi$ the stability, and the grid
values $s, \eta \in \{0, 0.1, \dots, 1\}$. Every unit has equal weight.
\begin{itemize}
\item \emph{Marginal curves} (Extended Data Figs.~6A, 6B): per unit, the mean of $p_i$
  over the other axis, $m^s_i(s) = \frac{1}{11}\sum_\eta p_i(1-s, \eta)$ and
  $m^\eta_i(\eta) = \frac{1}{11}\sum_s p_i(1-s,\eta)$. The black line is the mean
  over $\mathcal C$ at each grid value; the shaded fan is nine nested central
  quantile intervals across units, $[Q_q, Q_{1-q}]$ for
  $q = 0.05, 0.10, \dots, 0.45$ (R \texttt{quantile}, type 7), each drawn with
  opacity 0.06 so that the fill darkens toward the median; the dot-dashed red
  line is the identity.
\item \emph{Aligned marginal curves} (Figs.~5A, 5B). The mean of a set of
  curves cannot distinguish curves that rise gradually from curves that step,
  each at its own place, because averaging over the places blurs the step. Each
  curve is therefore shifted onto its own midpoint: with
  $r_i = \max_x m_i - \min_x m_i$, the midpoint $x^{50}_i$ is the axis value at
  which $m_i$ first reaches $\min_x m_i + r_i/2$, by linear interpolation
  between the two grid values that bracket it, and the curve is redrawn against
  $x - x^{50}_i$. Curves with $r_i = 0$ have no midpoint and are omitted (none along stability,
  76 of 440 along $\eta$). Each shifted curve is interpolated onto
  a common grid in steps of 0.05 (\texttt{approx}, no extrapolation), and the
  mean and the fan are formed at each offset exactly as in A and B.
\item \emph{The range the aligned panels are drawn over.} A curve covers a
  different stretch of the shifted axis depending on where its own midpoint
  sits, so towards the ends the mean and the quantiles would be taken over a
  shrinking and self-selected subset -- the combinations whose midpoint is
  early, which plateau lower than the rest. C and D are therefore drawn only
  over the offsets at which at least 60\% of the curves are defined, which is
  $-0.70$ to $+0.10$ along stability (at least 274 of 440 curves at every offset
  drawn) and $-0.40$ to $+0.50$ along $\eta$ (at least 242 of 364). The vertical
  dot-dashed red line marks the alignment point.
\item \emph{One-step differences} (Extended Data Figs.~7A, 7B):
  $\Delta_s \bar P(s + 0.05, \eta) = \bar P(s + 0.1, \eta) - \bar P(s, \eta)$
  (a $10 \times 11$ array placed at the step midpoints) and
  $\Delta_\eta \bar P(s, \eta + 0.05) = \bar P(s, \eta + 0.1) - \bar P(s, \eta)$
  for $\eta = 0, \dots, 0.9$ (an $11 \times 10$ array), where
  $\bar P(s, \eta) = |\mathcal C|^{-1}\sum_{i\in\mathcal C} p_i(1 - s, \eta)$ is
  the mean surface. Both panels share one white-to-black scale from 0 to the
  largest difference in either.
\item \emph{Single combinations} (Extended Data Figs.~7C--7F): the two marginals $m^s_i$ and
  $m^\eta_i$ of one unit against the axis value, with a vertical rule in the
  color of each curve at its own midpoint $x^{50}_i$. The four are units 195,
  203, 151 and 389, ordered by the midpoint of their stability curve (0.27,
  0.46, 0.68 and 0.92), so that the quantity the aligned panels of Fig.~5
  remove reads left to right across the row.
\end{itemize}
Because $p_i \in [0,1]$, the identity line marks a response proportional to the
axis value; a mean curve below it on most of the axis (smaller area under the
curve than $0.5$) means most of the rescue is gained near full stability or full specificity. The
mean surface $\bar P$ itself, printed tile by tile, is
Table~\ref{tab:imp-surface}; it is not a panel of the paper; it is drawn by
\texttt{imperfect\_treg\_main()} in \texttt{lib/fig\_imperfect\_treg.R}.

\subsection{Results behind the panels}

Table~\ref{tab:imp-surface} gives the pooled surface. No one-step
difference along $\eta$ or along stability is negative in any population; the
only zeros lie in the top row, where $\bar P$ has already reached 1. The largest single step is the last step of stability,
$s = 0.9 \to 1$ ($\phi = 0.1 \to 0$), which alone raises $\bar P$ by 0.291--0.422
(pooled; by $\eta$), 33--51\% of the whole rise from $s = 0$ to $s = 1$ at that
$\eta$. Table~\ref{tab:imp-steps} summarizes the steps and the curves per
population.
% from: pipeline/checks/imperfect_numbers.R, section 7

\begin{table}[H]
\centering
\caption{The mean probability of rescue $\bar P(1-\phi,\eta)$ over the 440
retained units (both branches), the surface that Extended Data Figs.~7A, 7B difference and that
\texttt{imperfect\_treg\_main()} (\texttt{lib/fig\_imperfect\_treg.R}) draws. Rows: stability $1-\phi$;
columns: $\eta$.
}
\label{tab:imp-surface}
\small
\begin{tabular}{l*{11}{r}}
\toprule
$1-\phi$ \textbackslash{} $\eta$ & 0 & 0.1 & 0.2 & 0.3 & 0.4 & 0.5 & 0.6 & 0.7 & 0.8 & 0.9 & 1\\
\midrule
1.0 & 0.726 & 0.780 & 0.827 & 0.892 & 0.953 & 0.994 & 1.000 & 1.000 & 1.000 & 1.000 & 1.000\\
0.9 & 0.359 & 0.407 & 0.454 & 0.495 & 0.531 & 0.585 & 0.619 & 0.647 & 0.670 & 0.685 & 0.709\\
0.8 & 0.235 & 0.272 & 0.317 & 0.357 & 0.388 & 0.421 & 0.452 & 0.490 & 0.515 & 0.533 & 0.561\\
0.7 & 0.160 & 0.189 & 0.221 & 0.251 & 0.285 & 0.331 & 0.365 & 0.395 & 0.428 & 0.443 & 0.461\\
0.6 & 0.078 & 0.120 & 0.155 & 0.189 & 0.216 & 0.241 & 0.279 & 0.312 & 0.343 & 0.382 & 0.393\\
0.5 & 0.045 & 0.063 & 0.094 & 0.127 & 0.163 & 0.192 & 0.217 & 0.246 & 0.275 & 0.299 & 0.325\\
0.4 & 0.018 & 0.035 & 0.049 & 0.080 & 0.104 & 0.136 & 0.166 & 0.189 & 0.224 & 0.252 & 0.284\\
0.3 & 0.008 & 0.013 & 0.024 & 0.047 & 0.068 & 0.087 & 0.112 & 0.136 & 0.161 & 0.192 & 0.225\\
0.2 & 0.003 & 0.004 & 0.010 & 0.025 & 0.045 & 0.065 & 0.078 & 0.104 & 0.123 & 0.148 & 0.173\\
0.1 & 0.000 & 0.002 & 0.005 & 0.009 & 0.020 & 0.043 & 0.056 & 0.075 & 0.098 & 0.120 & 0.139\\
0.0 & 0.000 & 0.000 & 0.002 & 0.004 & 0.010 & 0.020 & 0.036 & 0.055 & 0.078 & 0.094 & 0.111\\
\bottomrule
\end{tabular}
\end{table}

\begin{table}[H]
\centering
\caption{One-step differences and marginal curves per population. ``Area'':
trapezoidal area under the mean marginal curve on $[0,1]$ (identity: 0.5).
}
\label{tab:imp-steps}
\small
\resizebox{\linewidth}{!}{%
\begin{tabular}{lrrr}
\toprule
 & pooled & sterile & pathogenic help\\
 & (Fig.~5, Extended Data Fig.~7) & (Supplementary Fig.~2) & (Supplementary Fig.~3)\\
\midrule
units used                                      & 440 & 177 & 263\\
step along $\eta$: median (min--max)            & 0.025 (0.000--0.065) & 0.031 (0.000--0.079) & 0.021 (0.000--0.055)\\
step along $1-\phi$: median (min--max)          & 0.058 (0.000--0.422) & 0.060 (0.000--0.399) & 0.053 (0.000--0.438)\\
step $1-\phi$: $0.9\to1$, range over $\eta$     & 0.291--0.422 & 0.282--0.399 & 0.297--0.438\\
$\bar P$ at $(1-\phi,\eta) = (0,0)$ / $(1,1)$   & 0.000 / 1.000 & 0.000 / 1.000 & 0.000 / 1.000\\
mean curve along $1-\phi$ at $0, 0.5, 0.9, 1$   & 0.037, 0.186, 0.560, 0.925 & 0.052, 0.205, 0.571, 0.906 & 0.028, 0.174, 0.553, 0.937\\
mean curve along $\eta$ at $0, 0.5, 0.9, 1$     & 0.148, 0.283, 0.377, 0.398 & 0.136, 0.296, 0.415, 0.438 & 0.157, 0.275, 0.352, 0.372\\
area, $1-\phi$ / $\eta$                         & 0.257 / 0.277 & 0.272 / 0.291 & 0.246 / 0.268\\
\bottomrule
\end{tabular}}
\end{table}

Along stability the mean curve is flat-to-convex and rises steeply only over the
last step; the per-unit median of $m^s_i$ is 0 up to $s = 0.5$ (pooled) and
reaches 0.65 at $s = 0.9$ and 1.00 at $s = 1$, where the first quartile is 0.91.
Along $\eta$ the mean curve rises in steps of nearly equal size (pooled, 0.023
at the first step and 0.021 at the last) and ends at 0.398, because averaging
over all stabilities includes the rows $s < 1$ where even $\eta = 1$ rarely
rescues. The median of $m^\eta_i$ is $1/11 = 0.091$ at $\eta = 0$, the value of
a unit rescued only in the $s = 1$ row, and 0.273 at $\eta = 1$.

\subsection{Legend}

%% ---- from docs/methods_parts/fig05.tex, lines 446-455
\paragraph{Fig.~5.}
\emph{Stability and specificity change the probability of rescue in different
ways.} 440 combinations of both branches (177 sterile, 263 pathogenic help),
each integrated from its injury to $t = 12{,}000$ on an $11\times11$ grid of
Treg stability $1-\phi$ and antigenic specificity $\eta$, ten noise replicates
per cell. A rescue is a run whose tail-mean dead-cell load $E_a + E_n$ is below
$K_g$, and the marginal of a combination is $P(\text{rescue})$ averaged over the
other axis. \textbf{A}, the marginals along stability: nine nested central
quantile bands across combinations (5--95\% to 45--55\%, the fill darkening
towards the median), the mean over combinations (black) and the identity
(dot-dashed red). \textbf{B}, the same along $\eta$. \textbf{C} and \textbf{D},
the same curves with each shifted so that its own midpoint lies at zero
(dot-dashed red), drawn over the offsets at which at least 60\% of the curves
are defined. Aligning removes the spread of midpoints, which is what the mean in
A and B averages over, so a population of thresholds at different places can be
told from a population of curves that rise gradually.

%% ===========================================================================
\section{Generation of Extended Data Fig.~1: two forms of rescue, sterile and pathogenic help branches}
\label{sec:figS01}

Fig.~3A, 3B (sterile help) and Extended Data Fig.~2A, 2B (pathogenic help) and Supplementary Fig.~2 (pathogenic harm) hold the phase portraits of one design, and Extended Data Fig.~1 (sterile help in A--C, pathogenic help in D--F) and Supplementary Fig.~2
(pathogenic harm) the attractor shifts and rescue gaps of the same design, drawn once per
branch from the classes and rescue gaps of
Sections~\ref{M-sec:classes} and~\ref{M-sec:gap} of the Methods. The design, the definitions of all
panels and the tables for all of them are given here; Sections~\ref{sec:fig3}
and~\ref{sec:figS06} give what is specific to Fig.~3A, 3B and Extended Data Fig.~2A, 2B and to Supplementary Fig.~2.

%% ---- from docs/methods_parts/core.tex, lines 608-703
Scripts \texttt{figure\_scripts/Fig\_3.R} and \texttt{figure\_scripts/Extended\_Data\_Fig\_2.R} (Fig.~3A, 3B and Extended Data Fig.~2A, 2B, sterile and pathogenic help phase portraits),
\texttt{figure\_scripts/Extended\_Data\_Fig\_1.R} (sterile and pathogenic help shifts and gaps) and
\texttt{figure\_scripts/Supplementary\_Fig\_2.R} (pathogenic harm phase portraits, shifts and gaps), all drawing with
\texttt{lib/fig\_skeleton.R}. They read the skeleton and rescue files of their branch
and the parameter grids, and integrate only the example combinations.

\emph{Populations.} The bistable violins and the gap panels (Extended Data Figs.~1A, 1C, 1D, 1F; Supplementary Figs.~2C, 2E) use the bistable combinations with a measured rescue
gap; the ON\_mono violins (Extended Data Figs.~1B, 1E; Supplementary Fig.~2D) use all ON\_mono combinations (mono\_A, which would add further panels,
is empty in every branch). Each branch draws one Tregs-effect class: the sterile branch
holds help and neutral combinations in one file, and \texttt{figure\_scripts/Fig\_3.R}, \texttt{figure\_scripts/Extended\_Data\_Fig\_2.R} and
\texttt{figure\_scripts/Extended\_Data\_Fig\_1.R} keep the
help ones, so that all branches show a single class. Starting from the branch
(Table~\ref{M-tab:gate} of the Methods): sterile help 33,491 bistable, of which 4 have no gap, so 33,487
drawn, and 9,884 ON\_mono; pathogenic
help 16,225 bistable, of which 2,443 have opposite directions and 57 no Tregs-ON
threshold, so 13,725 drawn, and 1,270 ON\_mono; pathogenic harm 4,066 bistable, of which
2,313 opposite and 11 none, so 1,742 drawn, and 1,180 ON\_mono.
%% from: census [skel]

%% TODO: the split squares under the x axis (upper-left half the attractor reached with
%% Tregs OFF, lower-right half with Tregs ON, outlined in black inside the rescue gap)
%% of Figure 3 and S06 are not yet described here.
\emph{Phase portraits (Figs.~3A, 3B and Extended Data Figs.~2A, 2B; Supplementary Figs.~2A, 2B).} A (and C) is the bistable combination with the largest rescue gap;
B (and D) is the ON\_mono combination with the largest decrease of the low-burden attractor's total burden,
$\max(b^{\mathrm{lo,OFF}} - b^{\mathrm{lo,ON}})$. Each is integrated from
$e = 0, 10, \dots, 90$ on the initial-condition line under both conditions, to
$t = 12{,}000$ with output every 2 time units (settings of
Section~\ref{M-sec:solver} of the Methods, at most $10^5$ steps), and drawn in $(E_h, M_1)$ in
Figs.~3A, 3B and in $(E_h, P_{LP})$, the plane of the total burden, in Extended Data Figs.~2A, 2B and Supplementary Fig.~2.
Each trajectory is colored by the attractor it reaches, judged against its own
condition: lower basin if its final burden is at most the midpoint of that
condition's low- and high-burden attractor burdens (an ON\_mono Tregs-ON trajectory,
having no Tregs-ON high-burden attractor, is always lower basin). Filled circles are the
attractors stored by the sweep. Gray points mark the initial conditions.

\emph{Attractor shifts (Extended Data Figs.~1A, 1B, 1D, 1E; Supplementary Figs.~2C, 2D).} Tregs ON minus Tregs OFF in the total burden ("comp"), in
$E_h$ ("dEh") and, in the pathogenic panels, in $P_{LP}$ ("dPLP") of the low-burden (L) and
high-burden (H) attractors. For ON\_mono the single Tregs-ON attractor is compared with
each Tregs-OFF attractor. Violins: Gaussian kernel density (R \texttt{density}, default
bandwidth) trimmed to the data range and scaled to a common maximum width; inner box the
interquartile range, point the median, whiskers to the most extreme value within 1.5
interquartile ranges, points beyond drawn individually with horizontal jitter (seed 1).
In the pathogenic panels the graphics device omits category labels that would overlap
(comp H and dPLP L are not printed); the order is comp L, comp H, dEh L, dEh H, dPLP L,
dPLP H.

\emph{Rescue gap (Extended Data Figs.~1C, 1F; Supplementary Fig.~2E).} The fraction of the panel's combinations at each value of the
rescue gap, binned by rounding to the lattice step 0.078125 (a check stops the script if
any value falls outside the bins). The lattice runs unbroken from 0 to the maximum, so every bar is one step wide, and the
fraction is drawn on a logarithmic axis, which keeps bars four decades apart legible; the
minor lattice ticks are not drawn, there being several hundred of them at this step. The
in-panel text gives the minimum, mean and maximum of the gap, and, where any combination
is censored (Tregs ON reaches the low basin from every probed initial condition, so the
gap is only a lower bound), the censored fraction; no combination drawn in Extended Data Fig.~1C,
Extended Data Fig.~1F or Supplementary Fig.~2E is censored. No gap panel is drawn for the ON\_mono class, every
ON\_mono combination being censored, so that class has no resolved gap.

\begin{table}[H]
\centering
\caption{The numbers behind Extended Data Fig.~1 and Supplementary Fig.~2: medians [interquartile
range] of the attractor shifts, Tregs ON $-$ Tregs OFF.}
%% from: census [skel]
\label{tab:core-s01}
\small
\resizebox{\linewidth}{!}{%
\begin{tabular}{llrrr}
\toprule
panels & series & sterile (Extended Data Figs.~1A--1C) & pathogenic help (Extended Data Figs.~1D--1F) & pathogenic harm (Supplementary Figs.~2C--2E)\\\midrule
bistable & $n$ & 33,487 & 13,725 & 1,742\\
 & comp L & $-4.93$ [$-11.72$, $-2.16$] & $-4.41$ [$-9.42$, $-2.03$] & 2.49 [1.57, 3.88]\\
 & comp H & $-12.84$ [$-24.33$, $-5.58$] & $-9.62$ [$-20.01$, $-4.41$] & $-14.62$ [$-18.33$, $-9.88$]\\
 & dEh L & 4.93 [2.16, 11.72] & 4.78 [2.25, 10.28] & $-0.83$ [$-2.00$, 0.35]\\
 & dEh H & 12.84 [5.58, 24.33] & 9.64 [4.41, 20.05] & 14.63 [9.90, 18.59]\\
 & dPLP L & -- & 0.51 [0.00, 2.35] & 6.31 [4.20, 8.96]\\
 & dPLP H & -- & 0.00 [0.00, 0.11] & 0.00 [0.00, 0.71]\\
ON\_mono & $n$ & 9,884 & 1,270 & 1,180\\
 & comp L & $-13.18$ [$-24.70$, $-6.17$] & $-10.31$ [$-20.24$, $-4.91$] & 15.60 [7.17, 27.52]\\
 & comp H & $-51.33$ [$-63.96$, $-39.89$] & $-38.54$ [$-49.84$, $-31.36$] & $-14.39$ [$-22.87$, $-2.95$]\\
 & dEh L & 13.18 [6.17, 24.70] & 12.47 [5.94, 22.29] & $-5.02$ [$-13.85$, 0.11]\\
 & dEh H & 51.33 [39.89, 63.96] & 41.02 [32.98, 51.76] & 24.83 [16.78, 32.76]\\
 & dPLP L & -- & 2.75 [0.26, 8.01] & 22.57 [14.25, 35.82]\\
 & dPLP H & -- & 4.40 [0.85, 11.21] & 27.84 [18.03, 41.61]\\
\midrule
gap & exactly 0 / $>0$ / max & 30.7\% / 69.3\% / 48.59 & 58.7\% / 41.3\% / 33.28 & 65.7\% / 34.0\% / 17.11\\
ON\_mono gap & censored & 100\% (9,884) & 100\% (1,270) & 100\% (1,180)\\
\bottomrule
\end{tabular}}
\end{table}

\begin{table}[H]
\centering
\caption{The example combinations of the phase-portrait panels: key, thresholds, and the stored
attractors ($E_h$, $M_1$, $P_{LP}$; total burden $b$).}
%% from: census [skel] and the console output of the figure scripts
\label{tab:core-s01-ex}
\scriptsize\setlength{\tabcolsep}{2.5pt}
\resizebox{\linewidth}{!}{%
\begin{tabular}{lllll}
\toprule
panel & \texttt{tau\_id}, $K_g$, $\beta_r$, $r_P$ & thresholds OFF / ON & Tregs OFF lo / hi & Tregs ON lo / hi\\\midrule
Fig.~3A & 20723, 15, 0.0065, 0 & 56.60 / 8.01 (gap 48.59) & 53.50, 22.43, 0 / 25.79, 74.29, 0 & 93.69, 1.74, 0 / 73.15, 9.48, 0\\
Fig.~3B & 30070, 30, 0.0035, 0 & 34.96 / -- & 29.90, 18.33, 0 / 7.74, 93.20, 0 & 98.06, 0.15, 0 / --\\
Extended Data Fig.~2A & 30336, 35, 0.0025, 0.08 & 40.04 / 6.76 (gap 33.28) & 69.41, 27.47, 0 / 43.02, 82.59, 0 & 73.32, 22.69, 0 / 50.86, 60.22, 0\\
Extended Data Fig.~2B & 30173, 45, 0.0085, 0.08 & 13.87 / -- & 31.00, 21.03, 1.48 / 9.79, 90.47, 0 & 76.66, 2.35, 1.13 / --\\
Supplementary Fig.~2A & 20745, 35, 0.0065, 0.16 & 34.96 / 17.85 (gap 17.11) & 69.75, 35.49, 0.99 / 47.33, 95.19, 0 & 66.79, 33.88, 5.30 / 51.61, 80.19, 0\\
Supplementary Fig.~2B & 20324, 30, 0.0095, 0.08 & 39.57 / -- & 59.29, 7.11, 0 / 15.07, 58.31, 0 & 61.86, 4.78, 16.89 / --\\
\bottomrule
\end{tabular}}
\end{table}

\subsection{The pathogenic help branch: Extended Data Figs.~2A, 2B and Extended Data Figs.~1D--1F}

%% The former Supplementary Figure S02 is retired; its design, population and legend are
%% absorbed here (phase portraits now in Figure 3 C, D, shifts and gap in S01 D--F).
\emph{Data, selection and panels.} As the sterile branch above, for the pathogenic help
branch: the 14{,}995 retained gated pathogenic
combinations in which Tregs lower the low-burden attractor's total burden by more than 1
(Table~\ref{M-tab:gate} of the Methods; 2{,}500 of the 17{,}495 are excluded by
Section~\ref{M-sec:excl_decreasing} of the Methods, 2{,}443 for an inverted Tregs-ON
threshold and 57 for an unlabelled one). Panels Extended Data Figs.~1D, 1F use all
13{,}725 retained bistable combinations; panel Extended Data Fig.~1E uses all 1{,}270
ON\_mono combinations. Phase portraits (Extended Data Figs.~2A, 2B) are drawn in $(E_h, P_{LP})$, the plane of the total burden,
from $P_{LP}(0) = 10$, and the attractor shifts include $P_{LP}$ (dPLP). The medians and
interquartile ranges of panels Extended Data Figs.~1D, 1E, the gap fractions of panel Extended Data Fig.~1F and the example
combinations of Extended Data Figs.~2A, 2B are the pathogenic help column and the Extended Data Figs.~2A, 2B rows of
Tables~\ref{tab:core-s01} and~\ref{tab:core-s01-ex}.

\subsection{Legend}

%% ---- from docs/methods_parts/core.tex, lines 705-720
%% TODO: this first paragraph is the legend of the main Figure 3; move it to
%% Section~\ref{sec:fig3} when that section is written. The split squares under the
%% x axis are not described in it yet.
\paragraph{Fig.~3A, 3B and Extended Data Fig.~2A, 2B.}
Tregs shift both attractors and widen the range of recoverable injury, in the sterile
(A, B) and pathogenic (C, D) branches where Tregs help. Sterile branch, Tregs help
($r_P = 0$; the 43,375 of 89,328 gated
combinations in which Tregs lower the low-burden attractor's total burden by more than 1).
\textbf{A}, Tregs-ON bistable combinations: phase portrait in $(E_h, M_1)$ of the
combination with the largest rescue gap (\texttt{tau\_id} 20723, $K_g = 15$,
$\beta_r = 0.0065$; gap 48.59), trajectories from $E_h(0) = 0, 10, \dots, 90$
(gray points) with Tregs ON (solid) and OFF (dashed), colored by the basin reached (blue
lower, red higher; darker for Tregs OFF); filled circles, the stored attractors.
\textbf{B}, Tregs-ON monostable (ON\_mono) combinations:
phase portrait of the combination with the largest decrease of the low-burden
total burden (\texttt{tau\_id} 30070, $K_g = 30$, $\beta_r = 0.0035$): every Tregs-ON
trajectory reaches the low-burden attractor.
\textbf{C}, \textbf{D}, as A, B for the pathogenic help branch ($r_P > 0$, Tregs lower the
low-burden attractor's total burden by more than 1; 17,495 gated combinations), with phase portraits in
$(E_h, P_{LP})$ starting at $P_{LP}(0) = 10$.
\textbf{C}: \texttt{tau\_id} 30336, $K_g = 35$, $\beta_r = 0.0025$, $r_P = 0.08$
(gap 33.28). \textbf{D}: \texttt{tau\_id} 30173, $K_g = 45$, $\beta_r = 0.0085$,
$r_P = 0.08$.

\paragraph{Extended Data Fig.~1.}
The attractor shifts and rescue gaps of the two forms of rescue: sterile branch, Tregs
help in A--C ($r_P = 0$; the 43,375 of 89,328 gated combinations in which Tregs lower the
low-burden attractor's total burden by more than 1), pathogenic help branch in D--F (17,495 gated
combinations), with $P_{LP}$ shifts (dPLP) added to D and E.
\textbf{A}, attractor shift Tregs ON $-$ Tregs OFF in the total burden and in $E_h$ of the
low-burden (L) and high-burden (H) attractors of the Tregs-ON bistable combinations
($n = 33{,}487$). \textbf{B}, for the Tregs-ON monostable (ON\_mono) combinations, the
single Tregs-ON attractor against each Tregs-OFF attractor ($n = 9{,}884$); all ON\_mono
combinations are censored (Tregs ON reaches the low basin from every probed initial
condition). \textbf{C}, distribution of
the rescue gap $\Delta E_h(0)^*$ on its lattice (step 0.078125; 31\% exactly 0, 69\%
positive, maximum 48.6). \textbf{D}, \textbf{E}, \textbf{F}, as A, B, C for the
pathogenic help branch: $n = 13{,}725$ in D and F (59\% of gaps exactly 0, 41\% positive,
maximum 33.3) and $n = 1{,}270$, all censored, in E.

\section{Generation of Extended Data Fig.~3: load balance during the injury, ON\_mono combinations}
\label{sec:figS02}

\emph{Data, selection and panels.} Extended Data Fig.~3 is the ON\_mono
counterpart of Fig.~3C and Extended Data Fig.~2C and is generated, selected and drawn exactly as described in
Section~\ref{sec:fig4}. Population (Table~\ref{tab:lbgate}): sterile ON\_mono 11{,}159
gated combinations, 11{,}124 after the worst-probe deviation filter, all with a located,
increasing Tregs-OFF threshold, 9{,}869 of them in the Treg-help class, 500 drawn (seed 1) of
which 453 are Treg-help and are used; pathogenic help ON\_mono 1{,}270, 854 after the deviation
filter, all Treg-help, 500 drawn. Each combination is integrated once per
condition from $e_{\mathrm{mid}} = E_h(0)^*_{\mathrm{OFF}}/2$, midway between total
injury and the Tregs-OFF threshold. The endpoint statistics are the two ON\_mono
columns of Table~\ref{tab:lb}; 20 combinations per branch are drawn for display (seed
1, the same combinations in both columns).

\emph{Caveats specific to this class.} ON\_mono is a classification of the
initial-condition line only (Section~\ref{M-sec:classes} of the Methods): three pathogenic ON\_mono
Tregs-ON trajectories never reach $K_g$ and come to rest at $L = 3.8$--$7.7\,K_g$, i.e.\
the Tregs-ON system has a high-burden attractor reachable from $e_{\mathrm{mid}}$
(exception (ii) of Section~\ref{sec:fig4}); and two pathogenic ON\_mono Tregs-OFF
trajectories resolve because their Tregs-OFF threshold (0.039) lies within the first
bisection bracket (exception (iv)).

\subsection{Legend}

%% ---- from docs/methods_parts/fig02_03.tex, lines 234-237
\paragraph{Extended Data Fig.~3.}
As Fig.~3C and Extended Data Fig.~2C for ON\_mono combinations, integrated from $E_h(0)^*_{\mathrm{OFF}}/2$.
\textbf{(A, B)} Sterile, Tregs OFF and ON; \textbf{(C, D)} pathogenic help, Tregs OFF and
ON; 20 combinations per branch drawn from the Treg-help combinations integrated per branch,
453 sterile and 500 pathogenic (Table~\ref{tab:lb}).

\section{Generation of Extended Data Fig.~4: what a stronger ROS response buys and costs without Tregs}
\label{sec:figS03}

Fig.~4 asks what a brake buys at fixed damage. Extended Data Fig.~4 asks the prior
question --- what raising $\beta_r$ does on its own --- over the sweep's own grid, with Tregs
OFF throughout. Both quantities it draws are properties of a \emph{steady state}, so they are
read from the fixed points the sweep already stores and nothing is integrated
(\texttt{pipeline/45\_beta\_effero\_tradeoff.R}, \texttt{figure\_scripts/Extended\_Data\_Fig\_4.R}).

\paragraph{Population.} Every pathogenic combination whose Tregs-OFF system is bistable
(\texttt{n\_field\_stable\_off} $\ge 2$) with an increasing threshold direction, a finite
$E_h(0)^{*}$ and a trusted fixed-point map. A parameter set $(\texttt{tau\_id}, K_g, r_P)$
enters only if all nine $\beta_r$ levels survive those conditions, so every comparison along a
row is within one parameter set with $\beta_r$ the only difference: 12{,}211 sets, 109{,}899
combinations.

\paragraph{The two quantities.} Each stable steady state of a set counts once, and the chance
plotted is the share of them with the property, averaged over the sets of the cell:
\begin{itemize}
\item \emph{pathogen below $P_{\mathrm{cut}}$}: $P_{LP} < 1$ at the steady state, 1\% of the
  lamina propria carrying capacity. The criterion is a drawing choice and the figure does not
  turn on it: at cut-offs 1, 5 and 10 the chance rises with $\beta_r$ in every row, from
  59.8 to 71.4\%, 72.8 to 87.7\% and 82.3 to 94.7\% at $r_P = 0.08$, and from 18.3 to 36.7\%,
  30.7 to 54.7\% and 38.7 to 58.9\% at $r_P = 0.48$. Setting it to the extinction tolerance
  $10^{-3}$ gives clearance instead, which saturates at the fast end (15.6 to 16.2\% at
  $r_P = 0.48$) and so hides the gradient the figure is about.
\item \emph{epithelium above half}: $E_h > 50$ at the steady state.
\end{itemize}

\paragraph{What it shows.} Across $\beta_r = 0.0015$ to $0.0095$ the pathogen chance rises in
every row (59.8 to 71.4\% at $r_P = 0.08$; 18.3 to 36.7\% at $r_P = 0.48$) and the epithelium
chance falls in every row (38.9 to 6.5\%; 20.7 to 3.8\%). The cost is nearly flat in $r_P$
while the benefit shrinks with it, so the exchange is least favorable against the fastest
pathogens. Read across rows rather than along one, the same grid says that holding pathogen
control fixed costs more epithelium as $r_P$ rises: the smallest $\beta_r$ reaching
$P(P_{LP}^{\mathrm{SS}} < 1) \ge 0.35$ grows from 0.0015 at $r_P = 0.08$ to 0.0035 at
$r_P = 0.32$ and 0.0085 at $r_P = 0.48$, while the epithelium chance at that $\beta_r$ falls
from 0.39 to 0.11 and then 0.04. At $r_P = 0.48$ the pathogen chance never exceeds 0.37, so a
higher target is out of reach at any aggression on the grid.
% from: figure_scripts/Extended_Data_Fig_4.R, P_CUT=1

\paragraph{What it does not say.} The chance is over the steady states a parameter set can
settle on, not over the initial injuries that reach them: it counts what is possible, not how
often it happens. The grid also stops at the sweep's largest level, $\beta_r = 0.0095$, while
the damage-matched experiment of Fig.~4 runs far past it ($\beta_r^{*}$ up to 0.676); the
steady states above the sweep's range were not computed.

\section{Generation of Extended Data Fig.~5: three classifications of the damage-matched pairs}

Each of the 8{,}789 pairs is classified three times by comparing ON($\beta_r^{*}$) with
OFF($\beta_r$): by the first peak of $P_{LP}$ (Section~\ref{sec:first_peak}), by $P_{LP}$
integrated over the matching window, and by the steady-state $P_{LP}$ (the value at $T_{\max}$),
each as less or more pathogen with Tregs. The first peak is lower with Tregs in every pair, so
that classification does not split the sample and the routes separate at the other two. The
15 eliminating pairs are less on all three. Of the 8{,}774 persistent pairs, 597 are less on
all three; 7{,}235 less at the first peak and over the window but more at steady state; 897
less at the first peak only; and 45 less at the first peak and at steady state but more over the
window.
% from: pipeline/checks/fig06_methods_numbers.R, section 6
The steady-state load is higher with Tregs in 92.7\% of persistent pairs, by a median 0.17 units
on a median residual load of 1.89 without Tregs. Over all 8{,}789 pairs the first-peak difference
(OFF minus ON) has median 4.21 units and range 0.16 to 46.7; the integrated reduction median
6.3\% and range $-207$\% to 100\% (the lower end a pair in which the Tregs ON integral is 3.1
times the Tregs OFF one); the steady-state difference median $-0.17$ and range $-16.7$ to 22.9.
The negative tails belong to pairs that are worse over the window and at steady state:
combinations in which the brake tips a nearly cleared infection into a high endemic state, the
harm of Supplementary Figs.~2 and~3 (Section~\ref{sec:figS06}).
% from: pipeline/checks/fig06_methods_numbers.R, section 6

Extended Data Fig.~5 draws the 8{,}789 pairs as ribbons through the three classifications
(a Sankey diagram), one ribbon per route, stacked within each class in a fixed order and
colored by route, with the share of pairs in each class on the bars; above each axis, the
distribution of the corresponding difference (units of $P_{LP}$ for the first peak and the steady
state, per cent for the integral; positive when there is less pathogen with Tregs) as a
histogram on a symmetric logarithmic axis with a bin edge at zero, stacked by route in the same
colors, labeled with the median and the range. The ribbons are drawn largest first so that
the thin routes stay visible.

%% ===========================================================================
\label{sec:figS04}

It is drawn by \texttt{figure\_scripts/Extended\_Data\_Fig\_5.R}.

%% ===========================================================================
\section{Generation of Extended Data Fig.~7: where the rescue probability is gained, and its spread}
\label{sec:figS05}

\emph{Data, selection and panels.} The units, the stochastic simulations, the run and
unit filters, the rescue criterion and the definitions of all plotted quantities are
those of Fig.~5 (Section~\ref{sec:fig7}). Population: the 440 retained units of both
branches (977 drawn, 966 retained, 3 never and 246 always rescued;
Table~\ref{tab:imp-flow}), each with equal weight. Panels A and B are the one-step
differences $\Delta_s \bar P$ and $\Delta_\eta \bar P$; panels C--F are four single
units, each drawn as its two marginals with a rule at the midpoint of each curve.
The marginal curves of the whole population, with their mean and quantile fan, and
the same curves aligned on their midpoints, are Fig.~5. The step
medians and ranges, the last-step increase of stability and the mean curves are the
pooled column of Table~\ref{tab:imp-steps}; the sterile and pathogenic-help columns of
that table give the same quantities per branch, where the pattern is the same (areas
under the two mean curves 0.272 and 0.291 sterile, 0.246 and 0.268 pathogenic help),
which is why the branches are no longer drawn separately.

\subsection{Legend}

%% ---- from docs/methods_parts/fig05.tex, lines 457-466
\paragraph{Extended Data Fig.~7.}
\emph{Where the probability of rescue is gained, and what single combinations
look like} (the 440 combinations of Fig.~5). \textbf{A}, change in
$P(\text{rescue})$ for one grid step along stability,
$P(1-\phi + 0.1, \eta) - P(1-\phi, \eta)$, placed at the step midpoint.
\textbf{B}, the same for one step along $\eta$,
$P(1-\phi, \eta + 0.1) - P(1-\phi, \eta)$. A and B share one gray scale and are
differences of the mean surface, Table~\ref{tab:imp-surface}. \textbf{C--F},
four single combinations (units 195, 203, 151 and 389), each drawn as its two
marginals against the axis value -- $P(\text{rescue})$ averaged over $\eta$
(blue) and over stability (purple) -- with a vertical rule in the color of each
curve at its own midpoint. They are ordered by the midpoint of the stability
curve, 0.27, 0.46, 0.68 and 0.92, which is the spread that aligning removes in
 Figs.~5A and 5B. The branches are pooled; the per-branch values are the
sterile and pathogenic-help columns of Table~\ref{tab:imp-steps}.

\section{\texorpdfstring{Generation of Supplementary Fig.~1: displacement of the basin boundary at $t=0$}{Generation of Supplementary Fig. 1: displacement of the basin boundary at t = 0}}
\label{sec:fig5}

%% ---- from docs/methods_parts/fig02_03.tex, lines 243-438
\subsection{Pool}

The pool is every bistable and ON\_mono combination of the sterile and pathogenic help
branches as they stood after criterion~(3): 89{,}328 sterile (78{,}169 bistable, 11{,}159
ON\_mono) and 17{,}495 pathogenic
(16{,}225 bistable, 1{,}270 ON\_mono). No deviation gate is applied. The pool is put in random
order (seed 20260820) and split into contiguous chunks, one per cluster task
(\texttt{pipeline/21\_corridor\_tzero.R}). The pathogenic pool was measured in full (17{,}495
combinations, 100\%). \emph{The sterile pool was measured only in part}: the first round of
tasks hit the wall-clock limit, and what each task had completed (a prefix of its chunk) was
recovered. This round-1 set holds 16{,}990 of those 89{,}328 sterile combinations (19.0\%;
14{,}895 bistable and 2{,}095 ON\_mono); the remaining 72{,}338 were never measured, and all
sterile numbers below and in Supplementary Fig.~1B rest on this subsample. Because the pool was permuted
before chunking, the round-1 set is a random subsample; it matches the pool in the ON\_mono
share (12.33\% against 12.49\%) and in the rescue-gap distribution of its bistable
combinations (fraction above one lattice step 26.19\% against 26.66\%; two-sample
Kolmogorov--Smirnov $D = 0.0047$, $p = 0.95$). Every measured combination returned a result
(no failures in either branch).
%% from: pipeline/24_numbers_load_balance_boundary.R, blocks "pool", "notes"

\emph{One Treg-effect class.} The two pools are not matched in effect. The sterile branch is
defined by $r_P = 0$ alone and carries its help and neutral combinations together, whereas the
pathogenic branches are defined by the effect, so the 34{,}191 gated pathogenic neutral combinations
belong to no branch and were never measured (Section~\ref{M-sec:classes} of the Methods). Panels B and C
therefore keep the help class alone (\texttt{treg\_effect} $=$ help), so that both show one
class, as every other figure of the paper does. On the pathogenic pool this changes nothing,
the pool being help by construction. In the sterile pool it keeps 8{,}241 of the 16{,}990
measured combinations, 6{,}369 of the 14{,}895 bistable ones, and moves no correlation of
Table~\ref{tab:gainstats} by more than 0.004: with the neutral combinations pooled in, $\rho$ is
0.939, 0.938, 0.937, 0.935 and 0.935 for $M_1, R, D, E_a, E_n$ on 3{,}854--3{,}897 rows, and the
neutral combinations alone give 0.932, 0.932, 0.930, 0.930 and 0.928 on 1{,}091--1{,}095 rows.
The combination of panel A is help-class as well, so the whole figure is one class.
%% from: lib/fig_boundary_gain.R (the join and the filter); the pooled and neutral-only rho
%% from the same function with the filter lifted

\subsection{Fixed points and basin labels}

For each combination and condition, the 11 probes of the initial-condition line
($e = 0, 10, \dots, 90, 99$; Eq.~\ref{M-eq:icline} of the Methods) are integrated to $t = 12{,}000$ and the
fixed-point finder of Section~\ref{M-sec:fp} of the Methods is run with their final states as additional seeds,
with the grid seeds kept on both branches. Stable fixed points are grouped into attractors by
the rule of Section~\ref{M-sec:clustering} of the Methods (two stable points are one attractor when they are
within 1 in both $E_h$ and $P_{LP}$, single linkage); the low-burden attractor is the group of
lowest total burden (Eq.~\ref{M-eq:burden} of the Methods). A state $y$ is labeled by integrating it to
$t=12{,}000$ and assigning the final state to the nearest fixed point in the scaled distance
of Section~\ref{M-sec:ic-line} of the Methods, $\sum_i \big((y_i - f_i)/(|f_i|+1)\big)^2$ over
$E_h, E_a, E_n, P_{LP}, C_{LP}, M_1, R, D, T_r, S, P_L$: label 1 if that fixed point is stable
and in the low-burden group, 0 if it is stable and in another group, missing otherwise (or
if the integration fails). Solver settings as in the sweep, with at most $10^5$ steps per
output interval and output at $0$ and $12{,}000$ only (Table~\ref{M-tab:solver} of the Methods).

\subsection{Base state and scans}

The base state lies on the initial-condition line at
$e_{\mathrm{base}} = E_h(0)^*_{\mathrm{OFF}} + f\,(99 - E_h(0)^*_{\mathrm{OFF}})$ for the first $f$ in
$(0.02, 0.05, 0.10, 0.20, 0.35, 0.50, 0.75, 1)$ at which the Tregs-OFF system is labeled 1;
$f = 0.02$ in every measured combination (16{,}990 sterile, 17{,}495 pathogenic). From the base
state, each coordinate $X \in \{C_{LP}, E_h, E_a, E_n, M_1, R, D\}$ (and $P_{LP}$ on the pathogenic
branch) is set to a value $x$ with all other states at their base values; for $M_1$, $M_0$ is set
to $K_M - x$ so that the macrophage pool is conserved. 80 equally spaced values from the base
value $x^{\mathrm{base}}$ to an upper end are labeled under both conditions. The upper end starts at
$\min\!\big(\max(2\max(x^{\mathrm{hi}}, x^{\mathrm{lo}}, x^{\mathrm{base}}), 1),\, x^{\max}\big)$, where
$x^{\mathrm{hi}}, x^{\mathrm{lo}}$ are the coordinate's values at the Tregs-OFF high- and low-burden
fixed points, and is doubled (at most 10 passes) until each condition shows exactly one label
change or the ceiling $x^{\max}$ is reached ($C_{LP}$ 300, $E_h$ 100, $E_a$ 600, $E_n$ 800,
$M_1$ 100, $P_{LP}$ 300, $R$: $3\beta_rK_M/\mu_R$, $D$: $3\sigma_D/\tau_D$). Values with a
missing label in either condition are dropped. Where a condition has exactly one label change
along the scan, its boundary is located by 30 bisection halvings of the bracketing interval:
$x_{\mathrm{OFF}}$ (Tregs OFF) and $x_{\mathrm{ON}}$ (Tregs ON). Each (combination, coordinate) row is
classified: \emph{measured} (both boundaries located), \emph{no OFF wall} (no single Tregs-OFF
change), \emph{ON absent} (an OFF wall and no Tregs-ON label change anywhere up to the ceiling),
\emph{ON ambiguous} (an OFF wall and more than one Tregs-ON change).

\subsection{Boundary displacement and the rescue gap (panels B, C)}

For $X \in \{E_a, E_n, M_1, R, D\}$, and $P_{LP}$ on the pathogenic branch where it exists,
the plotted quantity is the additional displacement
tolerated with Tregs, relative to the displacement tolerated without them,
\begin{equation}
  G_X = 100\,\frac{x_{\mathrm{ON}} - x_{\mathrm{OFF}}}{x_{\mathrm{OFF}} - x^{\mathrm{base}}}\quad(\%),
  \label{eq:GX}
\end{equation}
defined on measured rows. ($x^{\mathrm{base}} = 0$ for $E_a, M_1, R, D$ and $K_E - e_{\mathrm{base}}$ for
$E_n$; $G_X$ depends on the base convention, which is why it is fixed.) $C_{LP}$ and $E_h$ are measured but not plotted, the first because its boundary is
located in fewer than half the pathogenic rows and the second because the rescue gap is itself
measured along $E_h$. The rescue gap $\gap$ (Eq.~\ref{M-eq:gap} of the Methods) lies on the
bisection lattice of Section~\ref{M-sec:bisection} of the Methods: the thresholds are bisected from a probe
spacing of 10 until the bracket is at most 0.1 wide, so gaps are multiples of
$10/2^7 = 0.078125$ (one sterile gap of 0.0703125 comes from the shorter last bracket
$[90, 99]$). Rows are kept if: the combination is bistable; the row is measured;
$G_X > 0$; the rescue gap is finite, not censored and $> 1.01\times0.078125$ (gaps at the
lattice floor, including exact zeros, are excluded: the bisection did not separate the two
thresholds there, and a logarithmic axis cannot show 0). Table~\ref{tab:gainledger} gives the
counts per step.
%% from: pipeline/24_numbers_load_balance_boundary.R, blocks "ledger", "q_kind"
\begin{table}[H]
\centering
\caption{Rows behind Supplementary Fig.~1B (sterile, the 8{,}241 help-class combinations of the round-1
subsample of 16{,}990) and Supplementary Fig.~1C (pathogenic help, all 17{,}495), bistable combinations,
per coordinate.}
\label{tab:gainledger}
\small
\begin{tabular}{llrrrrrr}
\toprule
 & & $E_a$ & $E_n$ & $M_1$ & $R$ & $D$ & $P_{LP}$\\
\midrule
sterile & bistable rows & 6{,}369 & 6{,}369 & 6{,}369 & 6{,}369 & 6{,}369 & --\\
 & both boundaries located & 6{,}365 & 6{,}363 & 6{,}326 & 6{,}362 & 6{,}364 & --\\
 & and $G_X>0$ & 6{,}365 & 6{,}363 & 6{,}326 & 6{,}362 & 6{,}364 & --\\
 & and gap finite, uncensored, $>0$ & 4{,}352 & 4{,}352 & 4{,}313 & 4{,}351 & 4{,}352 & --\\
 & lost: gap at one lattice step & 1{,}550 & 1{,}550 & 1{,}550 & 1{,}550 & 1{,}550 & --\\
 & \textbf{plotted} & \textbf{2{,}802} & \textbf{2{,}802} & \textbf{2{,}763} & \textbf{2{,}801} & \textbf{2{,}802} & --\\
pathogenic & bistable rows & 16{,}225 & 16{,}221 & 16{,}225 & 16{,}225 & 16{,}199 & 16{,}221\\
 & both boundaries located & 15{,}312 & 15{,}400 & 14{,}113 & 15{,}248 & 14{,}365 & 15{,}209\\
 & and $G_X>0$ & 14{,}792 & 14{,}869 & 14{,}034 & 14{,}819 & 14{,}364 & 14{,}686\\
 & and gap finite, uncensored, $>0$ & 5{,}203 & 5{,}251 & 4{,}743 & 5{,}166 & 5{,}064 & 5{,}154\\
 & lost: gap at one lattice step & 2{,}515 & 2{,}537 & 2{,}376 & 2{,}510 & 2{,}434 & 2{,}494\\
 & \textbf{plotted} & \textbf{2{,}688} & \textbf{2{,}714} & \textbf{2{,}367} & \textbf{2{,}656} & \textbf{2{,}630} & \textbf{2{,}660}\\
\bottomrule
\end{tabular}
\end{table}
In the pathogenic branch the rows lost before the gap filter are: no single Tregs-OFF change
(386, 347, 1{,}526, 434, 1{,}197, 434 for $E_a, E_n, M_1, R, D, P_{LP}$), more than one
Tregs-ON change (527, 474, 446, 522, 572, 578), no Tregs-ON change (0, 0, 140, 21, 65, 0), and
$G_X \le 0$ (520, 531, 79, 429, 1, 523). In the sterile help set at most 43 rows per coordinate are lost
before the gap filter. The drawn points pool 2{,}805 distinct sterile and 2{,}722 distinct
pathogenic combinations.
%% G_X <= 0 counts: measured minus G>0 rows of the ledger

\emph{Statistics.} Per coordinate, Spearman's rank correlation $\rho$ between $G_X$ and the
rescue gap (\texttt{cor.test}, \texttt{exact = FALSE}, asymptotic $p$; the gap is heavily tied on
its lattice) and an ordinary least-squares line of $\log_{10}\gap$ on $\log_{10}G_X$, drawn
over the range of the coordinate's data. Both use all kept rows. The points drawn are 9{,}000
rows sampled without replacement (seed 1) from the kept rows pooled over the
coordinates of each panel, five sterile and six pathogenic (13{,}970 sterile rows,
15{,}715 pathogenic), in gray; the fits are colored by coordinate.
Both axes are logarithmic; the $y$ range is the union of the two panels' ranges and is shared.
%% from: pipeline/24_numbers_load_balance_boundary.R, block "stats"
\begin{table}[H]
\centering
\caption{Supplementary Figs.~1B, 1C: per coordinate, kept rows, Spearman $\rho$ of $G_X$ against the rescue
gap, OLS line $\log_{10}\gap = a_0 + a_1\log_{10}G_X$, and the median (IQR) of $G_X$.
The asymptotic $p$ underflows to 0 for every coordinate except $R$ (pathogenic),
$p = 4.2\times10^{-210}$; all $p < 10^{-200}$.}
\label{tab:gainstats}
\small
\begin{tabular}{llrrrrl}
\toprule
branch & $X$ & $n$ & $\rho$ & $a_0$ & $a_1$ & $G_X$ (\%), median (IQR)\\
\midrule
sterile (round 1) & $M_1$ & 2{,}763 & 0.942 & $-1.702$ & 0.842 & 35.6 (18.4--88.5)\\
 & $D$ & 2{,}802 & 0.940 & $-1.757$ & 0.896 & 34.7 (17.9--85.9)\\
 & $R$ & 2{,}801 & 0.939 & $-1.758$ & 0.935 & 30.0 (15.8--72.3)\\
 & $E_n$ & 2{,}802 & 0.937 & $-1.762$ & 0.931 & 31.3 (16.2--75.0)\\
 & $E_a$ & 2{,}802 & 0.937 & $-1.763$ & 0.933 & 30.9 (16.1--74.5)\\
pathogenic & $E_n$ & 2{,}714 & 0.893 & $-1.684$ & 0.890 & 19.4 (12.7--38.1)\\
 & $E_a$ & 2{,}688 & 0.893 & $-1.683$ & 0.891 & 19.3 (12.5--38.0)\\
 & $P_{LP}$ & 2{,}660 & 0.892 & $-1.634$ & 0.849 & 19.4 (12.7--38.6)\\
 & $M_1$ & 2{,}367 & 0.773 & $-1.459$ & 0.722 & 19.6 (12.8--34.8)\\
 & $D$ & 2{,}630 & 0.702 & $-1.266$ & 0.584 & 21.5 (13.6--42.2)\\
 & $R$ & 2{,}656 & 0.550 & $-0.941$ & 0.448 & 12.3 (6.6--21.6)\\
\bottomrule
\end{tabular}
\end{table}

\emph{Consistency with the sweep's attractors.} For every measured combination the Tregs-OFF
attractor separation along each coordinate found by the boundary stage was compared with the
one stored by the sweep (\texttt{pipeline/22\_corridor\_attractors.R}). In the sterile round-1 set
they agree exactly on every coordinate. In the pathogenic branch they agree to $<10^{-6}$
(relative) on every coordinate but $P_{LP}$; on $P_{LP}$, 5{,}007 combinations differ by more than
1\%, where the differences are of order $10^{-10}$ in absolute value (median
$5.8\times10^{-11}$) on separations near zero. These combinations are kept in Supplementary Fig.~1C, which
plots $P_{LP}$. Excluding them changes $\rho$ by at most 0.06 (to 0.886, 0.886, 0.885, 0.777,
0.761, 0.548 for $E_n, E_a, P_{LP}, M_1, D, R$).

\emph{ON\_mono rows} are not plotted. In them the Tregs-ON boundary is absent within the searched
range in 93.9--96.3\% of rows per coordinate (sterile round 1, both effect classes,
2{,}094--2{,}095 rows) and
39.3--47.0\% (pathogenic, 1{,}018--1{,}267 rows), so $G_X$ is unbounded rather than large there.
%% from: block "mono_B"

\subsection{One combination in time (panel A)}

The combination shown is chosen by rule: among sterile bistable combinations of the round-1 set
with both boundaries along $E_n$ measured and an uncensored gap (14{,}886), candidates are tried
in decreasing rescue gap until one has both boundaries along $E_n$ at $t=0$ in the panel's own
scan (described below); the first candidate passes: kinetic vector 621 of grid 1
(\texttt{tau\_id} 10621), $K_g = 15$, $\beta_r = 0.0015$, $\gap = 29.297$
(the largest gap among those 14{,}886, so also the largest among the help-class subset the
panels A and C are drawn on; its corridor width along $E_n$ is also the largest).
Its base state has $e_{\mathrm{base}} = 39.534$, i.e.\ $E_n^{\mathrm{base}} = 60.466$, and the stored
boundaries are $x_{\mathrm{OFF}} = 61.639$ and $x_{\mathrm{ON}} = 91.958$, width $w = 30.32$. Three initial
conditions differ from the base state in $E_n$ only:
$\max(x_{\mathrm{OFF}} - w/2, x^{\mathrm{base}}) = 60.466$ (below; the base value itself),
$(x_{\mathrm{OFF}}+x_{\mathrm{ON}})/2 = 76.798$ (inside the corridor) and $x_{\mathrm{ON}} + w/2 = 107.118$ (above).
Each is integrated under both conditions on $t=0$ and 1{,}200 logarithmically spaced times from
$10^{-3}$ to 1{,}000, the window of panel A, with the solver settings above. The attractor each
trajectory is assigned to, which the brackets under panel A group them by, is the classifier's
label of its initial condition and not its position at the end of that window; the classifier
integrates to 12{,}000.
%% from: source_data/Supplementary_Fig_1_example.csv

\emph{Boundaries in time (A).} Along the trajectory from the lowest initial condition (which
resolves under both conditions, so both boundaries remain defined), at 60 sample times (20 equally
spaced on $[0,5]$ and 40 logarithmically spaced on $[2, 1000]$) the state of that condition's
trajectory at the nearest output time is taken, $E_n$ alone is varied, and the condition's
boundary is located: 40 values from 0 to $\min(800, \max(4E_n, 50))$, the upper end doubled up to
8 times (ceiling 800); a boundary is accepted only if exactly one label change is found, then
refined by 20 bisection halvings. Labels use the integration, fixed-point assignment and
attractor rule of the scans above, with the fixed points recomputed for this combination
(clustering tolerance 1). The boundary is thus re-located from the state reached at each time; it
is not the $t=0$ boundary transported by the flow. The 60 $\times$ 6 boundary searches are
cached (\texttt{data/cache\_bnd\_v6\_\ldots\_cheb1.rds}) and the rendered figure reads that cache.
At $t=0$ the located boundaries are 61.639 (OFF) and 91.958 (ON), equal to the stored scan to
within $5\times10^{-6}$. All 60 sample times give a boundary under both conditions from the lowest
initial condition. Beyond the window the corridor does not close. Locating both boundaries by
the same rule from the states reached at $t = 10{,}000$ gives a width of 52.68 against 52.26 at
$t = 1{,}000$, both walls having moved outward together by a further
4.2 ($S^{\mathrm{OFF}}$, 43.20 to 47.39) and 4.6 ($S^{\mathrm{ON}}$, 95.47 to 100.07) units of
$E_n$. At that time each boundary is located from its own condition's resting state, which
differ ($E_h = 96.77$ with Tregs OFF and 98.70 with Tregs ON), so the late width compares the two
basins from their own attractors rather than from one shared state as at $t = 0$.
The same question is put to the population of panels A and C
(\texttt{pipeline/25\_corridor\_persistence.R}): 500 combinations per branch are drawn at random
from those plotted along $E_n$, the lowest initial condition is integrated under each condition to
$t = 1{,}000$ and $t = 10{,}000$, and both boundaries are re-located from the state reached, by the
scan and bisection described above. A boundary the scan cannot resolve is recorded as missing, not
as a closed corridor. Every corridor that resolves is open at both times: 500 of 500 sterile, and
497 and 498 of 500 pathogenic. Median widths rise from 0.43 units of $E_n$ at $t = 0$ to 8.88 and
8.84 (sterile) and from 0.24 to 9.31 and 9.33 (pathogenic); 499 of the 500 sterile corridors are
wider at $t = 10{,}000$ than at $t = 0$, the exception being narrower by 0.8\% (2.380 against
2.400) and still open. In the three pathogenic rows that do not resolve, the scan brackets a
single Tregs-ON boundary above the Tregs-OFF wall at both times and only the bisection fails, so
the width is unresolved rather than the corridor closed. Shading: blue on the side from which the system resolves and pink on the side
from which it reaches the high-burden attractor, white between the two boundaries (the
corridor). The time axis is pseudo-logarithmic (linear below 1) and runs downwards; $E_n$ is
linear (the boundaries span a ratio of 3).

\emph{Attractors of the six trajectories (the brackets under A).} Each initial condition is
assigned to an attractor by the classifier, and the brackets under the panel group the ends of the
six trajectories by that assignment, each bracket spanning its group's end points (widened to 5\%
of the axis when narrower). From below, the low-burden attractor under both conditions; from
inside the corridor, the high-burden attractor without Tregs ($E_h = 69.4$, $M_1 = 24.5$) and the
low-burden attractor with Tregs ($E_h = 98.7$, $M_1 = 0.73$); from above, the high-burden
attractor under both (Tregs ON: $E_h = 79.0$, $M_1 = 14.8$). Three of the six therefore reach the
low-burden attractor and three the high-burden one.
%% from: source_data/Supplementary_Fig_1_example_final_states.csv

\subsection{Legend}

%% ---- from docs/methods_parts/fig02_03.tex, lines 440-458
\paragraph{Supplementary Fig.~1.}
Tregs move the basin boundary outward. \textbf{(A)} One sterile bistable combination
($K_g = 15$, $\beta_r = 0.0015$, rescue gap 29.3), with time running downward and
necrotic load $E_n$ on the horizontal axis. Colored curves: three initial conditions that differ
only in $E_n$ (below, inside and above the corridor), Tregs OFF solid and Tregs ON dashed. Dotted
black curves: the basin boundary along $E_n$ re-located at each time from the state of the
lowest trajectory, Tregs OFF plain and Tregs ON with circular markers (labeled
$S^{\mathrm{OFF}}$ and $S^{\mathrm{ON}}$, separatrix, in the key). Blue: resolving side; pink:
side leading to the high-burden attractor; white: the corridor between the two boundaries,
which at $t = 0$ spans $E_n = 61.6$--$92.0$. Brackets under the panel group the six trajectory
ends by the attractor each initial condition is assigned to; only the initial condition inside
the corridor changes outcome with Tregs. \textbf{(B, C)} Additional displacement tolerated
at $t = 0$ with Tregs, $G_X = 100(x_{\mathrm{ON}}-x_{\mathrm{OFF}})/(x_{\mathrm{OFF}}-x^{\mathrm{base}})$, against the
rescue gap, for $X = E_a, E_n, M_1, R, D$, and $P_{LP}$ in the pathogenic panel, in the combinations
where Tregs help: sterile (B; the
help class of the round-1 subsample, 16{,}990 of 89{,}328 combinations measured,
2{,}763--2{,}802 rows per coordinate) and pathogenic (C; 2{,}367--2{,}714 rows per coordinate). Bistable combinations with both boundaries located,
$G_X > 0$ and a rescue gap above one lattice step (0.078125). Gray: 9{,}000 measurements drawn at random from all
coordinates pooled, each one a combination paired with a single coordinate, so a combination
contributes as many points as it has measured coordinates; the fits and the correlations use
every measurement, not the drawn sample. Lines: per-coordinate least-squares fits of $\log_{10}\gap$ on $\log_{10}G_X$; labels:
Spearman $\rho$ per coordinate (all $p < 10^{-200}$). The vertical axis is the rescue gap $\gap$
(labeled $E_{h0}^{\mathrm{OFF}} - E_{h0}^{\mathrm{ON}}$). Logarithmic axes; shared $y$ range.

%% ===========================================================================
%% docs/methods_parts/fig04.tex -- Figure 6 of the manuscript and of the captions
%% document (figure_scripts/Fig_4.R) and Supplementary Figure S04: the damage-matched comparison.
%% Every number below was computed from ./data/ by
%%   Rscript pipeline/checks/fig06_methods_numbers.R
%% and is cited as [N§k], k = the numbered section of its printed output.
%% Notation: \beta_r is the ROS production rate (code: beta_ros, b_off for the reference level).

%% ---- from docs/methods_parts/fig04.tex (whole fragment; the master mirrors it)
\section{Generation of Supplementary Fig.~2: two forms of rescue, pathogenic harm branch}
\label{sec:figS06}

\emph{Data, selection and panels.} As Extended Data Figs.~2A, 2B and Extended Data Fig.~1
(Section~\ref{sec:figS01}), for the pathogenic harm branch: the 2{,}922 retained gated pathogenic
combinations in which Tregs raise the low-burden attractor's total burden by more than 1
(Table~\ref{M-tab:gate} of the Methods; 2{,}324 of the 5{,}246 are excluded by
Section~\ref{M-sec:excl_decreasing} of the Methods, 2{,}313 for an inverted Tregs-ON threshold
and 11 for an unlabelled one). Panels C and E of Supplementary Fig.~2 use all 1{,}742 retained bistable
combinations; panel Supplementary Fig.~2D uses all 1{,}180 ON\_mono combinations. Phase portraits (Supplementary Fig.~2) and shifts are drawn as for
Extended Data Figs.~2A, 2B and Extended Data Figs.~1D, 1E. Because Tregs raise the low-burden attractor's total burden in every
combination of this branch, the rule for panel Supplementary Fig.~2B (largest decrease of the low-burden
total burden) selects the ON\_mono combination with the smallest increase ($+1.01$). The
numbers are the pathogenic harm column and the Supplementary Fig.~2 rows of Tables~\ref{tab:core-s01}
and~\ref{tab:core-s01-ex}; the rescue-gap distribution of this branch is in
Table~\ref{M-tab:core-gap} of the Methods.

\subsection{Legend}

%% ---- from docs/methods_parts/core.tex, lines 731-740
\paragraph{Supplementary Fig.~2.}
As Extended Data Figs.~2A, 2B for the pathogenic harm branch (Tregs raise the low-burden
total burden by more than 1; 5,246 gated combinations). \textbf{A}: \texttt{tau\_id} 20745,
$K_g = 35$, $\beta_r = 0.0065$, $r_P = 0.16$ (gap 17.11).
\textbf{B}: \texttt{tau\_id} 20324, $K_g = 30$,
$\beta_r = 0.0095$, $r_P = 0.08$. In this branch the low-burden attractor's total burden rises
under Tregs in every ON\_mono combination, so the example rule (largest decrease) selects
the smallest increase ($+1.01$).

\paragraph{Supplementary Fig.~2.}
As Extended Data Fig.~1 for the pathogenic harm branch (5,246 gated combinations).
\textbf{A}, \textbf{C}:
$n = 1{,}742$ (66\% of gaps exactly 0, 34\% positive, maximum 17.1); the low-burden
attractor moves to higher total burden under Tregs in every combination (median $+2.49$),
mainly through $P_{LP}$ (median $+6.31$). \textbf{B}: $n = 1{,}180$, all censored.

\section{Generation of Supplementary Fig.~3: Treg treatments on the combinations Tregs harm}
\label{sec:figS07}

% ============================================================================
% docs/methods_parts/harm.tex -- group "harm": Supplementary Figure S07.
% Every number below was recomputed on 2026-10-01 from ./data by the scripts
% named in the % from: comments (paths relative to the repository root).
%   P50  = pipeline/50_harm_rescue_compute.R      (writes data/harm_rescue_data.rds)
%   P51  = pipeline/51_harm_explore_sample.R      (writes data/harm_explore_cache.rds)
%   NUM  = pipeline/checks/harm_methods_numbers.R (gate, dip landmarks, thresholds)
%   S07  = figure_scripts/Supplementary_Fig_3.R console output, source_data/Supplementary_Fig_3_*.csv
%   CHK1 = pipeline/checks/harm_check_delay_without_extinction.R
%   CHK2 = pipeline/checks/harm_check_load_at_dstar.R
% ============================================================================

%% ---- from docs/methods_parts/harm.tex, lines 18-231
Supplementary Fig.~3 asks, for the pathogenic parameter combinations in
which Treg induction worsens the outcome, which modification of the Treg
response removes that harm. Four modifications are compared: delaying the onset
of Treg induction, and scaling each of three Treg rate constants. All model
equations, parameter definitions, the extinction floor and the total burden
are those of Sections~\ref{M-sec:model}, \ref{M-sec:units} and \ref{M-sec:burden} of the Methods; this section describes only what is
specific to the three figures.

\subsection{Population: the gated harmed combinations}

The starting set is the pathogenic harm branch written by the
selection stage (\texttt{pipeline/03\_select.R}; Sections~\ref{M-sec:gate} and~\ref{M-sec:classes} of the Methods,
Table~\ref{M-tab:gate} of the Methods), i.e.\ every pathogenic combination (kinetic parameter row,
$K_g$, $\beta_r$, $r_P$) that passed the Tregs-OFF bistability
gate and whose low-burden attractor is raised by Tregs by more than the effect
tolerance on the total burden ($\Delta b_{\mathrm{lo}} > \delta = 1$). From this set the
combinations treated were selected in three further steps
(Table~\ref{tab:harm_gate}):
\begin{enumerate}
\item the Tregs-OFF threshold $E_h(0)^*_{\mathrm{OFF}}$ had to be finite and
  below 95, because a threshold returned in the last probe interval of the sweep
  ($E_h(0)\le 99$) is bounded by the end of the probed line rather than located
  by it; no combination was removed by this step;
\item combinations whose Tregs-ON resolving side runs in the opposite direction
  (\texttt{dir\_on} $=$ decreasing) were removed, because for them the Tregs-OFF
  and Tregs-ON thresholds are not commensurable; ON\_mono
  combinations (Tregs-ON direction undefined) are unaffected;
\item among combinations that remain bistable under Tregs, those whose rescue
  gap (shift of the threshold between the two conditions) was not at least
  $\texttt{GAP\_MIN}=10^{-9}$ were removed, i.e.\ zero, negative or non-finite
  gaps; ON\_mono combinations are exempt because their recorded gap is
  the whole $E_h(0)$ range rather than a distance between two crossings.
\end{enumerate}
The 1,773 surviving combinations span 185 kinetic parameter rows, all seven
$K_g$ values (15--45), all six non-zero $r_P$ values (0.08--0.48) and all
nine $\beta_r$ values (0.0015--0.0095). Their harm size on the
fixed points, $\Delta b_{\mathrm{lo}}$, has median 7.67 burden units (IQR
3.18--20.77, range 1.00--75.83).
% from: NUM section 1

\begin{table}[H]
\centering
\caption{Selection of the combinations treated in Supplementary Fig.~3.}
\label{tab:harm_gate}
\begin{tabular}{lrr}
\toprule
step & removed & remaining \\
\midrule
pathogenic harm branch (bistable 4{,}066; ON\_mono 1{,}180) & -- & 5,246 \\
$E_h(0)^*_{\mathrm{OFF}}$ finite and $<95$ & 0 & 5,246 \\
Tregs-ON line commensurable (inverted threshold: 2,313; unlabelled: 11) & 2,324 & 2,922 \\
bistable: rescue gap $\ge 10^{-9}$ (gap $=0$: 1,145; $<0$: 4) & 1,149 & 1,773 \\
\quad of which bistable / ON\_mono & & 593 / 1,180 \\
integration: non-finite state on some arm & 2 & 1,771 \\
harmed when integrated untreated ($\Delta b>1$ at delay 0) & 26 & 1,745 \\
\bottomrule
\end{tabular}
% from: NUM sections 1-2; P50 log (non-finite arms); S07 console ("1745 harmed")
\end{table}

\subsection{Arms, initial condition and integration}

For every combination 42 arms were integrated from the same initial condition:
one Tregs-OFF arm ($\theta=0$ throughout) and 41 treated arms
with Treg induction switched on ($\theta=1$). The initial
condition is the one used throughout the pathogenic branch
(Eq.~\ref{M-eq:icline} of the Methods): $E_h(0)=e_0$, $E_n(0)=K_E-e_0$, $M_0(0)=K_M=100$,
$T_0(0)=K_T=100$, $P_L(0)=500$, $P_{LP}(0)=10$ and all other states zero, with
\begin{equation}
e_0 = E_h(0)^*_{\mathrm{OFF}} + f_0\,(99 - E_h(0)^*_{\mathrm{OFF}}),\qquad f_0=0.01,
\label{eq:harm-e0}
\end{equation}
i.e.\ just inside the Tregs-OFF threshold on its resolving side, so that the
untreated system without Tregs resolves to its low-burden attractor. The treated
arms are
\begin{itemize}
\item \textbf{delay}: $\theta=0$ on $[0,d)$ and $1$ from $t=d$
  onward, integrated as two legs with the state at $t=d$ carried over;
  $d\in\{0,\allowbreak 10,\allowbreak 20,\allowbreak 40,\allowbreak 80,\allowbreak 160,\allowbreak 320,\allowbreak 960,\allowbreak 2400\}$ (9 levels);
\item \textbf{$\tau_{10}$ reduced}: $\tau_{10}\to m_{10}\,\tau_{10}$ (the rate at which
  the Treg product $S$ deactivates $M_1$) with Tregs on from $t=0$,
  $m_{10}\in\{1,\allowbreak 0.75,\allowbreak 0.5,\allowbreak 0.35,\allowbreak 0.25,\allowbreak 0.15,\allowbreak 0.1,\allowbreak 0.05,\allowbreak 0.03,\allowbreak 0.02,\allowbreak 0.0125,\allowbreak 0.0075\}$ (12 levels);
\item \textbf{$\tau_{t,d}$ increased}: $\tau_{t,d}\to m_{t,d}\,\tau_{t,d}$ (spontaneous
  Treg deactivation), $m_{t,d}\in\{1,\allowbreak 2,\allowbreak 4,\allowbreak 8,\allowbreak 16,\allowbreak 32,\allowbreak 64,\allowbreak 128\}$ (8 levels);
\item \textbf{$\tau_{t,a}$ reduced}: $\tau_{t,a}\to m_{t,a}\,\tau_{t,a}$ (Treg induction),
  $m_{t,a}\in$ the same 12 values as $m_{10}$.
\end{itemize}
Delay $0$ and multiplier $1$ are the same untreated Tregs-ON arm reached four
ways; their burden differences agree exactly (maximum absolute difference 0 over
1,771 combinations). The two descending ladders end at $1/0.0075\approx 133$-fold
so that the three rate treatments span the same fold range as the 128-fold
$\tau_{t,d}$ ladder.
% from: data/harm_rescue_data.rds $settings; NUM section 2

Every arm was integrated with \texttt{deSolve::lsodar} (absolute tolerance
$10^{-14}$, relative tolerance $10^{-8}$, maximum step 5, at most $10^5$ steps per
output interval; Table~\ref{M-tab:solver} of the Methods) to $T_{\max}=25{,}000$, with the absorbing extinction floor of
Section~\ref{M-sec:units} of the Methods: a root on $\max(P_L,P_{LP})-10^{-11}$ whose event sets
both compartments to exactly zero. Output was stored on a common grid of 301
times, $t=0$ and 300 points log-spaced from 0.05 to 25,000; the extra knot a
delayed arm has at $t=d$ is dropped, so medians across combinations are taken at
identical times. An arm was discarded if the solver failed or returned fewer
rows than requested; no arm was discarded on this ground. Two combinations
carried non-finite states with the full number of rows (the solver diverged
without reporting failure): combination (batch 1, row 1069, $K_g=40$,
$\beta_r=0.0045$, $r_P=0.24$), whose Tregs-OFF arm leaves the
reals at $t=18{,}387$, and (batch 3, row 646, $K_g=40$, $\beta_r=0.0055$,
$r_P=0.40$), whose $\tau_{10}\times0.15$ arm does so at $t=19{,}212$. Both
combinations were removed entirely (not only the arm), leaving 1,771 combinations
$\times$ 42 arms $=$ 74,382 summary rows and 22,388,982 trajectory rows. The
production run took 130.5 min for the arms and 85.9 min for the thresholds below
on a 32-core cluster node.
% from: P50 log (stored compute log of pipeline/50_harm_rescue_compute.R; identical counts in data/harm_rescue_data.rds; NUM section 2)

\paragraph{Endpoints and outcome class.}
For each arm the total burden at $T_{\max}$ is
$b=\sqrt{(K_E-E_h)^2+\max(P_{LP},0)^2}$ (Eq.~\ref{M-eq:burden} of the Methods), and the
burden difference is $\Delta b = b_{\mathrm{arm}} - b_{\mathrm{OFF}}$ against the
Tregs-OFF arm of the same combination. With the effect tolerance
$\delta=1$ burden unit, an arm is \emph{harm} if $\Delta b>\delta$,
\emph{neutral} if $|\Delta b|\le\delta$ and \emph{help} if $\Delta b<-\delta$.
Integrated untreated, 1,745 of the 1,771 combinations are harmed (median
$\Delta b=7.53$, IQR 3.19--20.38), 7 are neutral and 19 are helped; the 26 that
are not harmed (their sweep label was decided on the fixed points, not on this
trajectory) are carried in the time-series figures but have no threshold.
% from: NUM section 2

\paragraph{Pathogen-dip landmarks.}
On every arm the trajectory $P_{LP}(t_i)$ on the output grid is first floored,
$z_i = P_{LP}(t_i)$ if $P_{LP}(t_i)\ge 10^{-11}$ and $z_i=0$ otherwise, using the
same $10^{-11}$ as the extinction floor. With $k=\arg\max_i z_i$:
\begin{itemize}
\item $t_{\mathrm{peak}}=t_k$, $P_{LP}^{\mathrm{peak}}=z_k$;
\item $t_{\mathrm{dip}}$ is the first time $t_j$, $j\ge k$, at which $z$ attains
  its minimum over $\{t_k,\dots,t_n\}$, and $P_{LP}^{\mathrm{dip}}=z_j$; on an arm
  that reaches the floor this is the first output time at which $P_{LP}$ is zero;
  on an arm that does not, it is the trough between the peak and any rebound (or
  the last output time if $P_{LP}$ declines monotonically to the end);
\item $t_{\mathrm{ctrl}}$ is the first $t_i$ after the peak with
  $z_i<0.01\,z_k$ (undefined if none);
\item $t_{\mathrm{zero}}$ is the first $t_i$ after the peak with $z_i=0$
  (undefined if the arm never reaches the floor).
\end{itemize}
On the Tregs-OFF arm of the 1,745 harmed combinations: $t_{\mathrm{peak}}$ median
9.69 (IQR 7.44--14.38), $P_{LP}^{\mathrm{peak}}$ median 39.1 (IQR 24.7--62.1),
$t_{\mathrm{dip}}$ median 69.8 (IQR 31.7--338.9, range 15.7--25,000; 7 arms at the
last output time), $P_{LP}^{\mathrm{dip}}$ median 1.42 (IQR 0.20--3.40, range
0--42.1). $t_{\mathrm{ctrl}}$ is defined for 542 and $t_{\mathrm{zero}}$ for 246
(median 852, IQR 300--2,213). On every one of these 246 arms both $P_L$ and
$P_{LP}$ are exactly zero at $T_{\max}$ (the absorbing event fired), and
$t_{\mathrm{dip}}=t_{\mathrm{zero}}$; they are called the \emph{extinct}
combinations below. Over all 1,771 combinations the extinct set has 267 members.
% from: NUM section 3; CHK1 ("267 of 1771", "disagree on 0")

\subsection{Bisected thresholds}

For each harmed combination (1,745) and each treatment, the levels were ordered
by increasing intervention (delay ascending; $m_{10}$ and $m_{t,a}$ descending; $m_{t,d}$
ascending), and two boundaries were located: \emph{neutral}, the first level at
which $\Delta b\le\delta$ (harm has ended), and \emph{help}, the first level at
which $\Delta b<-\delta$. A boundary has a bracket if its predicate holds at some
ladder level $i>1$; the bracket is $[\ell_{i-1},\ell_i]$ at the first such level.
The response is flagged non-monotone if the predicate fails again at any level
beyond $\ell_i$. Inside the bracket the threshold was bisected by integrating
new arms: arithmetically for the delay until the bracket width was at most
$\mathrm{DTOL}=1$ time unit, geometrically (midpoint $\sqrt{\ell_{\mathrm{lo}}\ell_{\mathrm{hi}}}$ of the bracket
$[\ell_{\mathrm{lo}},\ell_{\mathrm{hi}}]$) for the multipliers until
$|\log(\ell_{\mathrm{hi}}/\ell_{\mathrm{lo}})|\le\log(1.01)$, i.e.\ a relative tolerance of 1\%.
The value returned is the end of the final bracket on the satisfied side, so the
reported threshold always satisfies the predicate; the delay threshold of the
neutral boundary is denoted $d^\ast$, the delay at which a combination stops
being harmed. Where the class measured at the neutral threshold was already help,
that value was reused as the help threshold without a second bisection
(delay 188, $\tau_{10}$ 8, $\tau_{t,d}$ 7, $\tau_{t,a}$ 8 combinations).
The counts are in Table~\ref{tab:harm_thr}. For the delay, every one of the 188
combinations with a help threshold crosses the whole dead band in one step:
its neutral and help thresholds are identical. One further combination
(batch 1, row 1905, $K_g=30$, $\beta_r=0.0075$, $r_P=0.08$)
reaches only neutral: its burden difference falls from $+4.46$ at delays
$\le80$ to $-0.91$ at delays $\ge160$ (just short of help), and its neutral
threshold is $d^\ast=95$. For the three rate treatments the two boundaries coincide in
only 7--8 combinations, with a median separation of 0.10--0.12 decades.
% from: P50 log; NUM section 4

\begin{table}[H]
\centering
\caption{Bisected thresholds per treatment and boundary over the 1,745 harmed
combinations. ``Fraction'' is the number with a threshold divided by 1,745 (the
plateau height of the curves in Supplementary Fig.~2C). Fold change is measured away from the
untreated arm: $d^\ast$ itself for the delay, $m_{t,d}$ for $\tau_{t,d}$, $1/m_{10}$ and
$1/m_{t,a}$ for $\tau_{10}$ and $\tau_{t,a}$.}
\label{tab:harm_thr}
\footnotesize
\setlength{\tabcolsep}{3pt}
\begin{tabular}{llrrrlll}
\toprule
treatment & boundary & $n$ & fraction & non-mono. & class at threshold & median [IQR] & range \\
\midrule
delay & neutral & 189 & 10.83\% & 0 & 188 help, 1 neutral & 448.8 [214.4--1{,}276] & 53.1--2{,}207 \\
delay & help & 188 & 10.77\% & 0 & 188 help & 450.0 [215.3--1{,}284] & 53.1--2{,}207 \\
$\tau_{10}$ reduced & neutral & 1{,}727 & 98.97\% & 1 & 9 help, 1{,}718 neutral & 3.56 [1.69--8.51] & 1.01--132.3 \\
$\tau_{10}$ reduced & help & 715 & 40.97\% & 280 & 715 help & 3.62 [1.87--7.19] & 1.08--125.1 \\
$\tau_{t,d}$ increased & neutral & 1{,}679 & 96.22\% & 1 & 7 help, 1{,}672 neutral & 4.63 [2.00--14.13] & 1.005--128.0 \\
$\tau_{t,d}$ increased & help & 705 & 40.40\% & 179 & 705 help & 5.05 [2.32--13.02] & 1.09--121.3 \\
$\tau_{t,a}$ reduced & neutral & 1{,}681 & 96.33\% & 1 & 8 help, 1{,}673 neutral & 4.62 [2.00--14.14] & 1.01--131.2 \\
$\tau_{t,a}$ reduced & help & 707 & 40.52\% & 183 & 707 help & 5.04 [2.32--13.29] & 1.09--121.2 \\
\bottomrule
\end{tabular}
% from: NUM section 4
\end{table}

Combinations without a delay threshold are still harmed at the longest delay
tested (2,400): 1,556 of the 1,745, of which 57 are extinct. The delay brackets
were $[40,80]$ (2), $[80,160]$ (26), $[160,320]$ (47), $[320,960]$ (53) and
$[960, 2{,}400]$ (61).
% from: CHK1; NUM section 4

\subsection{Panels of Supplementary Fig.~3}

%% ---- from docs/methods_parts/harm.tex, lines 235-325
\paragraph{Panel A: help against harm.}
This panel contrasts the combinations Tregs harm with those they help and
therefore draws on a separate sample (\texttt{pipeline/51\_harm\_explore\_sample.R}).
From all gated pathogenic combinations labeled help (17,495) or harm (5,246) by
the selection stage, the same three steps as above were applied
($E_h(0)^*_{\mathrm{OFF}}<95$: 17,494 help / 5,246 harm; Tregs-ON line
commensurable: 14,994 / 2,922; rescue gap: 6,932 / 1,773), and 250 combinations were
drawn at random per group (seed 3). For each, the Tregs-OFF and the Tregs-ON arm
were integrated from the same $e_0$ as above, with the same solver, tolerances
and extinction floor, to $T_{\max}=12{,}000$ on a grid of 420 times (300 linear
points on $[0,300]$, 120 log-spaced on $[300, 12{,}000]$). All 500 combinations
integrated under both conditions. The panel shows, per group and condition, the
median across combinations of $\max(P_{LP},1)$ at each output time, with the
interquartile range as a band (lighter for Tregs OFF), on log--log axes. The
display floor of 1 state unit ($10^{11}$ organisms) is a drawing choice only; at
$t=12{,}000$, 31.6\% of Tregs-OFF and 14.6\% of Tregs-ON arms are at or below
it, and the fraction exactly extinct is 28.8\% (help, OFF), 20.8\% (help, ON),
16.4\% (harm, OFF) and 1.2\% (harm, ON). Regenerating the sample (same seed) reproduces
the stored trajectories bit-identically.
% from: P51 log; NUM section 5; scratch re-run compared with data/harm_explore_cache.rds

\paragraph{Panel B: the delay at which harm ends against the elimination time.}
One point per harmed combination with a delay threshold (189): $x=t_{\mathrm{dip}}$
of its Tregs-OFF arm, $y=d^\ast$ (the bisected threshold of the class the
combination first reaches on the ladder: help for 188, neutral for 1), both on
log axes, with the identity line dashed. All 189 are extinct combinations
(100\%). Reported on the panel: Spearman rank correlation between $d^\ast$ and
$t_{\mathrm{dip}}$, $\rho=0.994$; the ratio $d^\ast/t_{\mathrm{dip}}$, median
0.907 (IQR 0.816--0.954); and the ordinary least-squares slope of
$\log_{10}d^\ast$ on $\log_{10}t_{\mathrm{dip}}$, 1.126 (95\% confidence interval
1.107--1.146). No combination whose Tregs-OFF arm is not extinct has a delay
threshold (Table~\ref{tab:harm_cross}), so the selection to extinct points is a
property of the data and not a filter.
% from: source_data/Supplementary_Fig_3_delay_fit.csv; S07 console

\paragraph{Panel C: cumulative conversion by the three rate treatments.}
For each rate treatment and boundary, the bisected thresholds were converted to
fold change away from the untreated arm ($1/m_{10}$, $m_{t,d}$, $1/m_{t,a}$), sorted, and drawn
as a step function $F(x)=\#\{\text{thresholds}\le x\}/n_0$ with $n_0=1{,}745$
(the number harmed at the untreated level, identical for all treatments), starting
at $(1,0)$. Color is the boundary (harm$\to$neutral, harm$\to$help), line type the
treatment; markers on the $\tau_{t,d}$ and $\tau_{t,a}$ curves (circles and
triangles, nine per curve, at log-spaced fold values offset by half a spacing) separate
the two curves, which nearly coincide because the Treg set point depends on
$\tau_{t,a}/\tau_{t,d}$. Final heights are the fractions in Table~\ref{tab:harm_thr}.
% from: NUM section 4

\paragraph{Panels D--G: one combination.}
The demonstration combination was chosen among the 178 harmed combinations that
have a help threshold under all four treatments (all 178 are extinct): ranked by
$|\log_{10}(t_{\mathrm{dip}}/\tilde t_{\mathrm{dip}})|$ with
$\tilde t_{\mathrm{dip}}=481.4$ the median over the 178, the first is batch 2,
row 1062, $K_g=25$, $\beta_r=0.0015$, $r_P=0.08$
($t_{\mathrm{dip}}=481.4$). Five arms were integrated for the figure from the
same $e_0$ and with the same solver settings, on 401 output times ($t=0$ and 400
log-spaced points from 0.05 to 25,000): Tregs OFF, and each treatment at its
help threshold (delay 451.2; $\tau_{10}\times0.6008$; $\tau_{t,d}\times1.874$;
$\tau_{t,a}\times0.5327$). Panels D--G show $T_r$, $M_1$, $P_{LP}$ and $E_h$
(linear) against time (log); $S$ is omitted because it is a first-order filter
of $T_r$ ($\dot S=\sigma_S T_r-\tau_S S$).
% from: S07 console ("row 2 demonstrates ...")

\begin{table}[H]
\centering
\caption{Least sufficient ladder level per combination (first level on the
ladder, counted from the untreated arm as step 1, at which the class is
reached), and its Spearman correlation with $t_{\mathrm{dip}}$ of the Tregs-OFF
arm. Here ``harm$\to$neutral'' counts the combinations that land exactly in the
neutral class ($|\Delta b| \le \delta$) at some ladder level; it differs from the
neutral boundary of Table~\ref{tab:harm_thr}, which is neutral or better, by the
combinations that cross the dead band in one step (219, 318 and 222 for the three
rate treatments; 188 for the delay). ``Extinct'' counts combinations whose Tregs-OFF arm
reaches the floor (or has $P_{LP}^{\mathrm{dip}}<10^{-3}$).}
\label{tab:harm_dose}
\begin{tabular}{llrrrrrr}
\toprule
treatment & class & $n$ & extinct & least step & most step & median step & Spearman \\
\midrule
delay & harm$\to$neutral & 1 & 1 & 6 & 6 & 6 & -- \\
delay & harm$\to$help & 188 & 188 & 5 & 9 & 8 & 0.960 \\
$\tau_{10}$ & harm$\to$neutral & 1{,}508 & 211 & 2 & 12 & 5 & 0.053 \\
$\tau_{10}$ & harm$\to$help & 715 & 244 & 2 & 12 & 5 & $-0.203$ \\
$\tau_{t,d}$ & harm$\to$neutral & 1{,}361 & 190 & 2 & 8 & 4 & 0.043 \\
$\tau_{t,d}$ & harm$\to$help & 705 & 237 & 2 & 8 & 4 & $-0.065$ \\
$\tau_{t,a}$ & harm$\to$neutral & 1{,}459 & 211 & 2 & 12 & 6 & 0.053 \\
$\tau_{t,a}$ & harm$\to$help & 707 & 238 & 2 & 12 & 6 & $-0.073$ \\
\bottomrule
\end{tabular}
% from: source_data/Supplementary_Fig_3_dose_table.csv
\end{table}
On the ladder (not bisected), the class reached first is later left again at a
larger intervention in 432 ($\tau_{10}$), 318 ($\tau_{t,d}$) and 466
($\tau_{t,a}$) combinations for neutral and 280, 179, 183 for help; never for the
delay.
% from: S07 console ("reached then left again")

%% ---- from docs/methods_parts/harm.tex, lines 448-506
\subsection{Diagnostics}

\paragraph{Delay requires extinction.}
For each of the 1,745 harmed combinations, the best class reached at any ladder
delay was cross-tabulated against extinction of its Tregs-OFF arm (both $P_L$ and
$P_{LP}$ below $10^{-11}$ at $T_{\max}$; over all 1,771 combinations this set is
exactly the set on which $t_{\mathrm{zero}}$ is defined) (Table~\ref{tab:harm_cross}). There are
\textbf{zero} combinations that a delay moves out of harm without an extinct
Tregs-OFF arm, over 1,499 non-extinct harmed combinations, including the 452 of
them whose Tregs-OFF dip falls below 1 state unit. The rate treatments do help
there: 492 ($\tau_{t,d}$), 493 ($\tau_{t,a}$) and 499 ($\tau_{10}$) of the 1,499
reach help at some ladder level. Among extinct combinations, 58 are never helped
by a delay (57 remain harmed, 1 reaches neutral with best $\Delta b=-0.91$); for
9 of them no delay is better than delay 0.
% from: CHK1

\begin{table}[H]
\centering
\caption{Best class reached at any ladder level, by extinction of the Tregs-OFF arm (1,745 harmed combinations).}
\label{tab:harm_cross}
\begin{tabular}{llrrr}
\toprule
treatment & Tregs-OFF arm & harm & neutral & help \\
\midrule
delay & not extinct (1,499) & 1{,}499 & 0 & 0 \\
delay & extinct (246) & 57 & 1 & 188 \\
$\tau_{10}$ & not extinct & 17 & 983 & 499 \\
$\tau_{10}$ & extinct & 1 & 29 & 216 \\
$\tau_{t,d}$ & not extinct & 58 & 949 & 492 \\
$\tau_{t,d}$ & extinct & 8 & 25 & 213 \\
$\tau_{t,a}$ & not extinct & 57 & 949 & 493 \\
$\tau_{t,a}$ & extinct & 7 & 25 & 214 \\
\bottomrule
\end{tabular}
% from: CHK1
\end{table}

\paragraph{What orders the departure from the identity in Supplementary Fig.~2B.}
For each of the 189 points the Tregs-OFF arm was re-integrated (same settings)
and the pathogen load it carries at $t=d^\ast$ was read off. All 189 arms reach
the absorbing state and none regrows after first crossing the floor. The ratio
$d^\ast/t_{\mathrm{dip}}$ has Spearman correlation $-0.939$ with $P_{LP}(d^\ast)$,
the strongest of the candidates examined (others: $t_{\mathrm{dip}}$ $0.864$,
killing margin $(\alpha_{1,e}M_1+\alpha_{1,r}R)/r_P$ at $d^\ast$ $-0.800$,
peak margin $-0.667$; $K_g$ 0.12, $\beta_r$ $-0.003$, $r_P$ 0.20).
$P_{LP}(d^\ast)$ has median $7.6\times10^{-11}$ (IQR $2.4\times10^{-11}$--$9.5\times10^{-10}$):
for most combinations the pathogen is at or within one decade of the floor when
the delayed Tregs arrive. The extreme point (batch 2, row 745, $K_g=30$,
$\beta_r=0.0035$, $r_P=0.08$; ratio 0.297, $t_{\mathrm{dip}}-d^\ast=199.9$)
carries the largest load at $d^\ast$ (0.35).
% from: CHK2

\paragraph{Numerical robustness of the extreme point.}
An earlier diagnostic re-bisected $d^\ast$ of that point under tighter tolerances, a
smaller maximum step, denser output, other integrators and other extinction floors, and
concluded that $d^\ast$ is not a numerical artefact. That diagnostic reads a data
format that the final pipeline no longer writes and needs several hundred integrations
per setting; it has \emph{not} been repeated on the final data, and its conclusion is
reported here as unverified.

\subsection{Legend}

%% ---- from docs/methods_parts/harm.tex, lines 510-532
\paragraph{Supplementary Fig.~3.}
Treg treatments on the combinations Tregs harm.
\textbf{(A)} Median pathogen load $P_{LP}$ (floored at 1 for display) of 250
randomly drawn gated combinations Tregs help (blue) and 250 they harm (red),
each under Tregs OFF (dashed) and Tregs ON (solid), from the same initial
condition; bands are interquartile ranges across combinations.
\textbf{(B)} For each of the 189 harmed combinations a Treg-induction delay can
rescue, the bisected delay at which harm ends ($d^\ast$) against the time its
Tregs-OFF arm eliminates the pathogen ($t_{\mathrm{dip}}$); green, the
combination reaches help (188), cyan, neutral (1); dashed, identity. Inset:
$n=189$, 100\% with an extinct Tregs-OFF arm, Spearman 0.99, median ratio
$d^\ast/t_{\mathrm{dip}}$ 0.91 (IQR 0.82--0.95), log--log slope 1.13 (95\% CI
1.11--1.15).
\textbf{(C)} Cumulative fraction of the 1,745 harmed combinations converted to
neutral or better (cyan) and to help (green) against fold change of
$\tau_{10}$ (reduced, solid), $\tau_{t,d}$ (increased, dashed with circles) and
$\tau_{t,a}$ (reduced, dotted with triangles) away from the untreated arm, using
bisected thresholds.
\textbf{(D--G)} One combination (batch 2, row 1062, $K_g=25$,
$\beta_r=0.0015$, $r_P=0.08$): $T_r$, $M_1$, $P_{LP}$ and $E_h$
under Tregs OFF (black) and under each treatment at the least level that
reaches help: delay 451.2, $\tau_{10}\times0.60$, $\tau_{t,d}\times1.87$,
$\tau_{t,a}\times0.53$.

\section{Figures, scripts and data}
\label{sec:fig-script-data}

\begin{longtable}{@{}p{0.13\linewidth}p{0.29\linewidth}p{0.26\linewidth}p{0.24\linewidth}@{}}
\caption{Figure, figure script (in \texttt{figure\_scripts/}; shared drawing code in \texttt{lib/}),
data files read from \texttt{data/}, and the pipeline stages that write them. All
scripts also read the kinetic grids \texttt{data/tau\_grid\_v6\_1..5.rds} where
parameters are needed.}\label{tab:fig-script-data}\\
\toprule
Figure & Script & Data read & Written by\\\midrule
\endfirsthead
\toprule
Figure & Script & Data read & Written by\\\midrule
\endhead
\bottomrule
\endfoot
Fig.~2 & \texttt{Fig\_2.R} & \texttt{effero\_knockout.rds}; for F, G also \texttt{effero\_fg\_extra.rds}, \texttt{effero\_fg\_extra2.rds} & \texttt{pipeline/06\_effero\_knockout.R} (the additional units: with \texttt{UNITS\_FILE} from \texttt{06b\_effero\_fg\_extra\_units.R})\\
Figs.~3A, 3B; Extended Data Figs.~2A, 2B & \texttt{Fig\_3.R}, \texttt{Extended\_Data\_Fig\_2.R}, \texttt{lib/fig\_skeleton.R} & \texttt{skeleton\_sterile.rds}, \texttt{rescue\_sterile.rds}, \texttt{skeleton\_pathogenic\_help.rds}, \texttt{rescue\_pathogenic\_help.rds} & \texttt{pipeline/03\_select.R}\\
Fig.~3C; Extended Data Fig.~2C & \texttt{Fig\_3.R}, \texttt{Extended\_Data\_Fig\_2.R}, \texttt{lib/fig\_load\_balance.R} & \texttt{load\_balance\_sterile.rds}, \texttt{load\_balance\_pathogenic\_help.rds} & \texttt{pipeline/20\_load\_balance.R} (from \texttt{rescue\_*}, \texttt{skeleton\_*}, \texttt{selected\_big.rds})\\
Fig.~4 & \texttt{Fig\_4.R}, \texttt{lib/fig\_damage\_matched.R}, \texttt{lib/fig\_demo\_case.R} & \texttt{fig04\_gated\_units.rds}, \texttt{fig04\_damage\_matched\_gated.rds}, \texttt{fig04b\_v3\_trajectories.rds} (Extended Data Fig.~5 also reads \texttt{fig06\_first\_peak.rds}; the fourth arm of the Methods, \texttt{aux\_gated\_fourth\_arm.rds}, is read by the numbers script only) & \texttt{pipeline/43\_damage\_matched\_gated\_units.R}, \texttt{40\_damage\_matched.R} (\texttt{cluster/damage\_matched\_gated.sh}), \texttt{41\_damage\_matched\_trajectories.R}, \texttt{42\_first\_peak.R}; \texttt{pipeline/48\_gated\_fourth\_arm.R}\\
Fig.~5; Extended Data Fig.~6 & \texttt{Fig\_5.R}, \texttt{Fig\_5\_assemble.py}, \texttt{lib/fig\_imperfect\_treg.R} & \texttt{fig05\_switch\_prob\_cells.rds}, \texttt{fig05\_switch\_prob\_units.rds}, \texttt{fig05\_units.csv} & \texttt{pipeline/60}--\texttt{63\_imperfect\_*.R}\\
Extended Data Fig.~1 & \texttt{Extended\_Data\_Fig\_1.R}, \texttt{lib/fig\_skeleton.R} & \texttt{skeleton\_sterile.rds}, \texttt{rescue\_sterile.rds}, \texttt{skeleton\_pathogenic\_help.rds}, \texttt{rescue\_pathogenic\_help.rds} & \texttt{pipeline/03\_select.R}\\
Extended Data Fig.~3 & \texttt{Extended\_Data\_Fig\_3.R}, \texttt{lib/fig\_load\_balance.R} & as Fig.~3C and Extended Data Fig.~2C & \texttt{pipeline/20\_load\_balance.R}\\
Extended Data Fig.~4 & \texttt{Extended\_Data\_Fig\_4.R} & \texttt{fig\_beta\_effero.rds} & \texttt{pipeline/45\_beta\_effero\_tradeoff.R} (from \texttt{fp\_big\_batch\_*.rds})\\
Extended Data Fig.~5 & \texttt{Extended\_Data\_Fig\_5.R}, \texttt{lib/fig\_damage\_matched.R} & as Fig.~4 & as Fig.~4\\
Extended Data Fig.~6 & \texttt{Extended\_Data\_Fig\_6.R}, \texttt{lib/fig\_imperfect\_treg.R} & as Fig.~5 & \texttt{pipeline/60}--\texttt{63\_imperfect\_*.R}\\
Extended Data Fig.~7 & \texttt{Extended\_Data\_Fig\_7.R}, \texttt{lib/fig\_imperfect\_treg.R} & as Fig.~5 & \texttt{pipeline/60}--\texttt{63\_imperfect\_*.R}\\
Supplementary Fig.~1 & \texttt{Supplementary\_Fig\_1.R}, \texttt{lib/fig\_boundary\_gain.R}, \texttt{lib/corridor\_defs.R} & \texttt{corridor\_q\_sterile\_round1.rds}, \texttt{corridor\_q\_pathogenic\_help.rds}, \texttt{cache\_bnd\_v6\_\ldots\_cheb1.rds} & \texttt{pipeline/21}--\texttt{23\_corridor\_*.R}; cache by \texttt{lib/corridor\_defs.R} when \texttt{Supplementary\_Fig\_1.R} first runs\\
Supplementary Fig.~2 & \texttt{Supplementary\_Fig\_2.R}, \texttt{lib/fig\_skeleton.R} & \texttt{skeleton\_pathogenic\_harm.rds}, \texttt{rescue\_pathogenic\_harm.rds} & \texttt{pipeline/03\_select.R}\\
Supplementary Fig.~3 & \texttt{Supplementary\_Fig\_3.R} & \texttt{harm\_rescue\_data.rds}, \texttt{harm\_explore\_cache.rds}, \texttt{skeleton\_pathogenic\_harm.rds} (panels D--G re-integrated) & \texttt{pipeline/50\_harm\_rescue\_compute.R}, \texttt{51\_harm\_explore\_sample.R}\\
\end{longtable}

Upstream of all branch files: \texttt{pipeline/00\_parameter\_grids.R} (the stored grids;
refuses to overwrite them), \texttt{01\_sweep.R} $\to$ \texttt{fp\_big\_batch\_<grid>.rds}
via \texttt{02\_assemble.R}, and \texttt{03\_select.R} $\to$ \texttt{selected\_big.rds},
\texttt{skeleton\_<branch>.rds}, \texttt{rescue\_<branch>.rds}. Every number of
Sections~\ref{M-sec:model}--\ref{M-sec:classification} and of
Figs.~3A, 3B, Extended Data Figs.~2A, 2B, Extended Data Fig.~1 and Supplementary Fig.~2 is printed by \texttt{pipeline/census\_core.R};
those of Fig.~3C, Extended Data Figs.~2C and~3 and Supplementary Fig.~1 by \texttt{pipeline/24\_numbers\_load\_balance\_boundary.R}; of Fig.~4
and Extended Data Fig.~5 by \texttt{pipeline/checks/damage\_matched\_numbers.R}; of Fig.~5 and Extended Data Figs.~6 and~7 by \texttt{pipeline/checks/imperfect\_numbers.R}; of Supplementary Fig.~3 by
\texttt{pipeline/checks/harm\_methods\_numbers.R} and the two diagnostic scripts in
\texttt{pipeline/checks/}; of Fig.~2 by \texttt{figure\_scripts/Fig\_2.R}
(\texttt{source\_data/Fig\_2\_*.csv}).

\end{document}

